% AA Classe Del Documento
\documentclass[11pt, a4paper]{article}

% AA Tipografia E Lingua
% BB Lingua E Codifica Font
\usepackage[italian]{babel}
\usepackage[T1]{fontenc}
\usepackage{lmodern}

% BB Geometria Pagina E Microtipografia
\usepackage[margin=2cm, top=1.5cm, bottom=1.5cm]{geometry}
\usepackage{parskip}
\usepackage{microtype}

% AA Gestione Dei Colori
% BB Pacchetto Colori
\usepackage[table]{xcolor}

% BB Palette Colori Personalizzata
\definecolor{darkgreen}{RGB}{0,100,0}
\definecolor{theoryblue}{RGB}{0,84,147}
\definecolor{backcolour}{RGB}{245,245,245}
\definecolor{codegray}{rgb}{0.5,0.5,0.5}

\definecolor{codebg}{HTML}{F8F9FA}
\definecolor{codeframe}{HTML}{E2E8F0}
\definecolor{codekeywords}{HTML}{7C3AED}
\definecolor{codedocker}{HTML}{0284C7}
\definecolor{codecomments}{HTML}{64748B}
\definecolor{codestrings}{HTML}{0D9488}
\definecolor{codenumbers}{HTML}{94A3B8}
\definecolor{codetext}{HTML}{1E293B}

% AA Matematica E Notazione Scientifica
\usepackage{amsmath, amssymb, amsfonts}
\usepackage{mathtools}
\usepackage{siunitx}

% BB OTTIMIZZAZIONE MATRICI LUNGHE PER LA STAMPA
\setlength{\arraycolsep}{10pt}
\usepackage{bm}

% AA Tabelle Ed Allineamenti
\usepackage{booktabs}
\usepackage{array}
\usepackage{multirow}
\usepackage{tabularx}
\usepackage{arydshln}

% AA Grafica E Disegno Vettoriale
% BB Gestione Figure E Posizionamento
\usepackage{graphicx}
\graphicspath{{./Images/Orchestrazione/}}
\usepackage{subcaption}
\usepackage{float}

% BB TikZ E Pgfplots
\usepackage{tikz}
\usetikzlibrary{arrows.meta, patterns, angles, quotes, calc}
\usetikzlibrary{positioning}
\usepackage{pgfplots}
\pgfplotsset{compat=1.18}

% AA Listati Di Codice (Listings)
% BB Configurazione Generale
\usepackage{listings}
\usepackage{upquote}

% BB Configurazione Generale
\lstset{
    basicstyle=\ttfamily\footnotesize\color{codetext},
    backgroundcolor=\color{codebg},
    frame=single,
    rulecolor=\color{codeframe},
    framesep=6pt,
    xleftmargin=8pt,
    xrightmargin=8pt,
    breaklines=true,
    breakatwhitespace=true,
    showstringspaces=false,
    keepspaces=true,
    columns=flexible,
    upquote=true,
    tabsize=2,
    captionpos=b,
    % Stili e Colori Predefiniti Globali
    keywordstyle=\bfseries\color{codekeywords},
    commentstyle=\itshape\color{codecomments},
    stringstyle=\color{codestrings}
}

% BB Definizione Dei Linguaggi

% CC YAML
\lstdefinelanguage{yaml}{
    keywords={true,false,null,yes,no,on,off},
    keywordstyle=\bfseries\color{codekeywords},
    comment=[l]{\#},
    commentstyle=\itshape\color{codecomments},
    stringstyle=\color{codestrings},
    morestring=[b]',
    morestring=[b]",
    sensitive=false
}

% CC Dockerfile
\lstdefinelanguage{dockerfile}{
    keywords={FROM, RUN, CMD, LABEL, EXPOSE, ENV, ADD, COPY, ENTRYPOINT,
              VOLUME, USER, WORKDIR, ARG, ONBUILD, STOPSIGNAL, HEALTHCHECK,
              SHELL, AS},
    keywordstyle=\bfseries\color{codedocker},
    comment=[l]{\#},
    commentstyle=\itshape\color{codecomments},
    stringstyle=\color{codestrings},
    morestring=[b]",
    morestring=[b]',
    sensitive=true
}

% CC INI
\lstdefinelanguage{ini}{
    comment=[l]{\#},
    morecomment=[l]{;},
    commentstyle=\itshape\color{codecomments},
    stringstyle=\color{codestrings},
    morestring=[b]",
    morestring=[b]',
    keywordstyle=\bfseries\color{codekeywords},
    sensitive=false
}

% CC JSON
\lstdefinelanguage{json}{
    keywords={true,false,null},
    keywordstyle=\bfseries\color{codekeywords},
    comment=[l]{//},
    morecomment=[s]{/*}{*/},
    commentstyle=\itshape\color{codecomments},
    stringstyle=\color{codestrings},
    morestring=[b]",
    sensitive=true
}

% CC PYTHON (Dialetto nativo pulito)
\lstset{defaultdialect=[3]Python}
\lstdefinelanguage{python}[]{Python}{
    morekeywords={as, async, await, True, False, None},
    sensitive=true
}

% CC BASH / SHELL
\lstdefinelanguage{bash}[]{sh}{
    morekeywords={docker, cd, run, build, stop, rm, rmi, push, pull, network, compose, openstack, rq, worker, keystone-manage, time, sleep, echo, exit, sudo, export, systemctl},
    sensitive=true
}

% CC HTML (Autonomo, previene il loop infinito del parser XML/HTML)
\lstdefinelanguage{html}{
    keywords={img, src, alt, class, id, div, span, table, tr, td, th, href, a, p, h1, h2, h3, form, input, button},
    keywordstyle=\bfseries\color{codekeywords},
    comment=[s]{<!--}{-->},
    commentstyle=\itshape\color{codecomments},
    stringstyle=\color{codestrings},
    morestring=[b]",
    morestring=[b]',
    sensitive=false
}

% CC SQL / CYPHER (Autonomo con supporto per // di Cypher)
\lstdefinelanguage{sql}{
    keywords={MATCH, RETURN, DETACH, DELETE, WHERE, CREATE, SET, WITH, ORDER, BY, ASC, DESC, LIMIT, SELECT, INSERT, INTO, VALUES, FROM, UPDATE, DROP, TABLE, JOIN, ON, AND, OR, NOT},
    keywordstyle=\bfseries\color{codekeywords},
    comment=[l]{//},
    morecomment=[l]{--},
    commentstyle=\itshape\color{codecomments},
    stringstyle=\color{codestrings},
    morestring=[b]',
    morestring=[b]",
    sensitive=false
}

% AA Algoritmi E Pseudocodice
\usepackage[ruled,vlined,linesnumbered]{algorithm2e}
\SetKwData{Dati}{Input}
\SetKwData{Risultato}{Output}
\SetKw{Inizializza}{Inizializza}
\SetKwFor{Per}{per}{fare}{fine per}
\SetKwComment{Commento}{// }{}

% AA Struttura, Titoli Ed Ambienti
% BB Personalizzazione Titoli Sezioni
\usepackage{titlesec}
\titleformat{\section}{\normalfont\LARGE\bfseries}{\thesection}{1em}{}
\titleformat{\subsection}{\normalfont\Large\bfseries}{\thesubsection}{1em}{}
\titleformat{\subsubsection}{\normalfont\large\bfseries}{\thesubsubsection}{1em}{}

% BB Teoremi Ed Ambienti Ausiliari
\usepackage{comment}
\newtheorem{theorem}{Definizione}

% BB Liste Ed Elenchi Puntati
\usepackage{enumitem}
\setlist{noitemsep}
\usepackage{pifont}

% AA Riferimenti, Link Ipertestuali E Segnalibri
\usepackage[unicode]{hyperref}
\hypersetup{
    colorlinks=true,
    linkcolor=blue,
    urlcolor=blue,
    citecolor=blue,
    pdftitle={3) Orchestrazione Delle Risorse In Sistemi Distribuiti},
    pdfauthor={Francesco Lo Verde},
    pdfpagemode=FullScreen,
}
\usepackage{bookmark}

% AA Metadati Del Documento
\title{3) Orchestrazione Delle Risorse In Sistemi Distribuiti}

\date{\today}

\author{Francesco Lo Verde}
\begin{document}
\maketitle

% AA Orchestrazione_1 [1-33]    ------------------------------------------------------------

\section{Introduzione A Kubernetes E Alla Containerizzazione}

\subsection{Richiamo Fondamentale: Dai Modelli Di Virtualizzazione Ai Container}
Per comprendere la genesi e l'architettura dei sistemi di orchestrazione moderni, è indispensabile analizzare l'evoluzione delle tecniche di isolamento delle risorse di calcolo. Nel paradigma della virtualizzazione tradizionale basata su macchina virtuale (\emph{Virtual Machine}, VM), l'astrazione viene operata al livello dell'hardware fisico mediante un componente denominato \textbf{Hypervisor} (o \emph{Virtual Machine Monitor}, VMM). L'hypervisor può risiedere direttamente a contatto con il ferro (Tipo 1, \emph{bare-metal}) oppure operare al di sopra di un sistema operativo ospitante (Tipo 2, \emph{hosted}). In entrambi i casi, ciascuna macchina virtuale istanziata incapsula un intero sistema operativo \emph{Guest}, dotato del proprio kernel, dei propri moduli di avvio, di librerie di sistema e dei binari applicativi. Questo approccio garantisce una robusta separazione dei domini di sicurezza mediante estensioni hardware della CPU (quali Intel VT-x o AMD-V), ma introduce un considerevole costo computazionale (\emph{overhead}):
\begin{itemize}[noitemsep]
    \item Lentezza nelle fasi di bootstrap a causa del ciclo di avvio completo del sistema operativo ospitato;
    \item Duplicazione ridondante delle strutture dati e delle routine del kernel;
    \item Occupazione non trascurabile di memoria RAM e cicli CPU dedicati esclusivamente al sostentamento del kernel \emph{Guest}.
\end{itemize}

% [IMMAGINE DEL FILE DI SLIDE CON NOME "Orchestrazione_1"] Confronto architetturale tra virtualizzazione tradizionale basata su Hypervisor e containerizzazione basata su Container Engine — Slide 5 / Pagina 5 del PDF

Al contrario, la tecnologia dei \textbf{Container} realizza la virtualizzazione a livello del sistema operativo (\emph{OS-level virtualization}). I container condividono il medesimo kernel del sistema operativo host, eliminando l'overhead del sistema operativo guest. L'isolamento e la segregazione delle risorse all'interno dell'ambiente Linux sono garantiti nativamente da due primitive fondamentali del kernel:
\begin{enumerate}[noitemsep]
    \item \textbf{Control Groups (\emph{cgroups})}: meccanismo preposto al monitoraggio, alla limitazione gerarchica e alla contabilità delle risorse computazionali (allocazione di quote CPU, quantitativi massimi di RAM, throughput I/O su disco, contatori di memoria swap e priorità di scheduling);
    \item \textbf{Namespaces}: primitiva volta a circoscrivere la visibilità delle strutture di sistema per un gruppo di processi. I namespace isolano le tabelle di routing e le interfacce di rete (\texttt{net}), l'albero dei processi (\texttt{pid}), i punti di montaggio del filesystem (\texttt{mnt}), la comunicazione tra processi (\texttt{ipc}), i nomi di dominio e host (\texttt{uts}), e i mapping degli identificativi utente (\texttt{user}).
\end{enumerate}

Grazie a questo disaccoppiamento, un container non è altro che un processo Linux standard (o un gruppo di processi) eseguito con vincoli rigidi e visibilità confinata. Il \textbf{Container Engine} (ad esempio Docker o containerd) funge da intermediario logico per orchestrare la creazione delle immagini a livelli sovrapposti (\emph{overlay filesystems}, come UnionFS o OverlayFS) e configurare le primitive di sistema all'atto dell'avvio.

\subsection{Origine, Filosofia Progettuale Ed Etimologia}
Sebbene i container risolvano il problema dell'impacchettamento uniforme delle applicazioni (risolvendo il noto paradosso \emph{"funziona sulla mia macchina"}), l'esecuzione di architetture distribuite composte da decine di migliaia di container operanti su scala industriale richiede uno strato di supervisione centralizzato. Google affrontò questa sfida internamente a partire dai primi anni 2000 attraverso lo sviluppo di due sistemi di gestione cluster proprietari:
\begin{itemize}[noitemsep]
    \item \textbf{Borg}: il framework monolitico di prima generazione, scritto in C++, che tuttora alimenta l'infrastruttura di elaborazione batch e i servizi di produzione di Google, capace di coordinare miliardi di container alla settimana;
    \item \textbf{Omega}: l'evoluzione di Borg, orientata a un'architettura decentralizzata con uno stato condiviso a consistenza ottimistica (\emph{transaction-based state store}).
\end{itemize}

Nel 2014, capitalizzando oltre un decennio di esperienza operativa accumulata con Borg e Omega, Google ha avviato il progetto \textbf{Kubernetes} (spesso abbreviato nell'alfanumerico \textbf{K8s}, denotando le otto lettere interposte tra la `K' e la `s'), sviluppandolo interamente nel linguaggio Go. Nel 2015, Google ha donato Kubernetes come progetto fondante alla neonata \textbf{Cloud Native Computing Foundation (CNCF)}, un'organizzazione neutrale sotto l'egida di Linux Foundation volta a standardizzare l'ecosistema del cloud computing orientato ai microservizi.

Etimologicamente, il termine deriva dal sostantivo greco antico \emph{"$\kappa\upsilon\beta\epsilon\rho\nu\acute{\eta}\tau\eta\varsigma$"} (\emph{kybernētēs}), traducibile come \emph{"timoniere"} o \emph{"pilota di una nave"}. Questo concetto riflette in modo iconografico il ruolo primario della piattaforma: fungere da guida e nocchiero nel governare flotte eterogenee di container all'interno di infrastrutture complesse.

% [IMMAGINE DEL FILE DI SLIDE CON NOME "Orchestrazione_1"] Illustrazione simbolica del timoniere ed etimologia greca del termine Kubernetes — Slide 3 / Pagina 3 del PDF

\subsection{La Necessità Dell'Orchestrazione Dei Container}
In uno scenario di produzione reale basato su microservizi, un'applicazione web complessa non è confinata a un singolo processo isolato. Essa richiede una combinazione organica di livelli frontend, logica applicativa di backend, code di messaggi asincrone e motori di database. Gestire tali contesti operativi manualmente su larga scala comporta criticità non banali:
\begin{itemize}[noitemsep]
    \item \textbf{Gestione dinamica dei picchi di carico}: si considerino, ad esempio, le fluttuazioni di traffico che investono una piattaforma di commercio elettronico durante eventi promozionali straordinari o periodi festivi. In presenza di un repentino incremento di traffico, il sistema deve essere in grado di istanziare istantaneamente decine o centinaia di repliche dell'applicazione; al termine dell'evento, le risorse computazionali devono essere prontamente rilasciate per ottimizzare i costi e i consumi dell'infrastruttura sottostante;
    \item \textbf{Tolleranza ai guasti ed elevata disponibilità (\emph{High Availability})}: in un data center, i singoli server fisici sono soggetti a guasti hardware improvvisi, interruzioni di connettività di rete o esaurimento delle risorse. Un orchestratore ha il compito di rilevare istantaneamente tali anomalie, riassegnando automaticamente i carichi di lavoro dai nodi degradati o non responsivi verso nodi sani;
    \item \textbf{Bilanciamento del carico e instradamento}: l'infrastruttura deve distribuire equamente il traffico utente in ingresso tra le repliche attive dei servizi, esponendo un punto d'accesso unificato e statico astraendo la natura effimera dei singoli container;
    \item \textbf{Aggiornamenti continui a zero disservizio (\emph{Zero-Downtime Rollouts})}: possibilità di distribuire nuove versioni applicative mediante strategie a rilascio progressivo (\emph{Rolling Updates}), ripristinando istantaneamente la versione precedente (\emph{Rollback}) qualora si riscontrino errori critici a runtime.
\end{itemize}

Queste esigenze hanno determinato l'affermazione dei sistemi di orchestrazione. Sebbene esistano soluzioni alternative — quali \textbf{Docker Swarm}, focalizzato sulla semplicità d'uso in cluster ridotti ma privo di logiche avanzate di autoscaling e controllo fine dei volumi, o \textbf{Apache Mesos}, orientato alla condivisione di risorse distribuite a grana fine per carichi scientifici eterogenei — Kubernetes si è imposto come standard \emph{de facto} industriale per la sua estendibilità, la robustezza del suo modello di rete e l'efficacia del paradigma dichiarativo.

\subsection{Paradigma Imperativo Vs Paradigma Dichiarativo}
La differenza concettuale fondamentale tra la gestione tradizionale dei sistemi operativi e l'approccio ingegneristico di Kubernetes risiede nel paradigma di controllo:
\begin{itemize}[noitemsep]
    \item \textbf{Approccio Imperativo}: l'amministratore impartisce sequenze rigide di comandi indicando \emph{come} ottenere un risultato (ad esempio: \emph{"alloca una macchina virtuale, esegui il pull dell'immagine, avvia il processo sulla porta 80, verifica che il processo sia attivo"}). Tale paradigma è intrinsecamente fragile: qualsiasi errore intermedio lascia il sistema in uno stato parziale non deterministico;
    \item \textbf{Approccio Dichiarativo}: l'amministratore specifica \emph{cosa} deve esistere nel sistema, ovvero il \textbf{Desired State} (stato desiderato), demandando alla piattaforma il compito di determinare i passi necessari a raggiungerlo. In Kubernetes, la definizione dell'infrastruttura e delle applicazioni avviene mediante manifest strutturati in formato YAML o JSON.
\end{itemize}

Il motore di Kubernetes implementa costantemente un ciclo di riconciliazione a retroazione (\emph{reconciliation loop}):
\[
\text{Errore di Stato } e(t) = \text{Stato Desiderato } S_{\text{des}} - \text{Stato Corrente } S_{\text{act}}(t)
\]
Fintanto che $e(t) \neq 0$, i meccanismi di controllo autonomi intervengono per applicare modifiche all'infrastruttura (istanziare pod, eliminare container anomali, riconfigurare endpoint di rete), garantendo la proprietà di \textbf{Self-Healing} attraverso:
\begin{itemize}[noitemsep]
    \item \emph{Autoplacement}: posizionamento ottimale dei container sui nodi idonei;
    \item \emph{Autorestart}: riavvio immediato dei processi in caso di terminazione inaspettata;
    \item \emph{Autoreplication}: mantenimento del numero esatto di repliche attive;
    \item \emph{Autoscaling}: espansione o contrazione orizzontale e verticale in base alle metriche telemetriche.
\end{itemize}

\newpage

% ==============================================================================
% SEZIONE: ARCHITETTURA DEL CLUSTER KUBERNETES
% ==============================================================================

\section{L'Architettura Del Cluster Kubernetes}

Un cluster Kubernetes è composto da una gerarchia distribuita di macchine (server fisici o macchine virtuali) ripartite funzionalmente in due macro-ambiti: il \textbf{Control Plane} (storicamente denominato \textbf{Master Node}) e l'insieme dei nodi operativi di calcolo, detti \textbf{Worker Nodes} (un tempo indicati come \emph{Minions}).

% [IMMAGINE DEL FILE DI SLIDE CON NOME "Orchestrazione_1"] Architettura di un cluster Kubernetes con nodo Master e multipli Worker Nodes — Slide 9 / Pagina 9 del PDF

% [IMMAGINE DEL FILE DI SLIDE CON NOME "Orchestrazione_1"] Scomposizione dei componenti interni del Master Node e del Worker Node — Slide 10 / Pagina 10 del PDF

\subsection{Il Control Plane (Master Node)}
Il Control Plane agisce come sistema operativo distribuito del cluster. Esso ha la responsabilità esclusiva di memorizzare e validare lo stato globale, coordinare le decisioni di scheduling e orchestrare gli eventi di gestione del ciclo di vita dei carichi di lavoro. Nei deployment industriali orientati all'alta disponibilità (\emph{High Availability}), i componenti del Control Plane vengono replicati su più istanze (tipicamente 3 o 5 nodi fisici distinti) per prevenire singoli punti di fallimento (\emph{Single Point of Failure}, SPoF).

I componenti logici che costituiscono il Master Node sono:

\subsubsection{1. Kube-Apiserver}
Il \textbf{kube-apiserver} costituisce il componente cardine e il punto di accesso unificato a tutte le funzionalità del cluster.
\begin{itemize}[noitemsep]
    \item Espone l'interfaccia di programmazione applicativa RESTful ufficiale di Kubernetes su protocollo HTTPS;
    \item Intercetta tutte le comunicazioni in ingresso, siano esse inviate da client umani esterni (tramite la CLI `kubectl`), da dashboard grafiche di management o dai componenti interni dei nodi;
    \item Implementa le catene di elaborazione relative ad autenticazione (\emph{Authentication}), autorizzazione basata su ruoli (\emph{Role-Based Access Control}, RBAC) e mutazione/validazione delle richieste tramite \emph{Admission Controllers};
    \item È progettato per essere completamente \textbf{stateless}: non memorizza alcuno stato localmente, ma delega la persistenza a un database transazionale distribuito. Questa proprietà permette a \texttt{kube-apiserver} di scalare orizzontalmente in modo trasparente mediante l'interposizione di un bilanciatore di carico di rete (L4).
\end{itemize}

\subsubsection{2. Datastore Distribuito etcd}
La memorizzazione durevole dello stato del cluster è affidata a \textbf{etcd}, un database chiave-valore (\emph{key-value store}) distribuito, consistente e open-source.
\begin{itemize}[noitemsep]
    \item etcd rappresenta l'unica fonte autoritativa di verità (\emph{Single Source of Truth}) dell'intero cluster; in esso risiede la totalità dei manifest, dei segreti, delle configurazioni e degli stati operativi dei componenti;
    \item La consistenza forte dei dati è garantita dall'implementazione dell'algoritmo di consenso distribuito \textbf{Raft}. Raft impone la presenza di un leader eletto a maggioranza; per tale ragione, al fine di garantire la tolleranza a $f$ fallimenti di nodi senza interruzione di servizio, il cluster etcd deve essere composto da un numero dispari di membri pari a:
    \[
    N_{\text{nodi}} = 2f + 1
    \]
    Di conseguenza, per resistere alla perdita di un nodo ($f=1$) è richiesto un cluster a 3 istanze, mentre per tollerare il fallimento simultaneo di due nodi ($f=2$) si impiega un'architettura a 5 nodi;
    \item etcd implementa primitive native di lock distribuito e meccanismi di notifica degli eventi a lunga durata (\emph{Watchers}), che permettono a componenti come l'API server di ricevere notifiche push immediate su qualsiasi variazione di stato delle chiavi.
\end{itemize}

\subsubsection{3. Kube-Scheduler}
Il \textbf{kube-scheduler} è il processo di controllo deputato alla selezione vincolata del nodo ottimale sul quale eseguire ciascun Pod di nuova creazione. Il processo decisionale di scheduling è altamente complesso e opera sequenzialmente mediante due fasi logiche: \emph{Filtering} (o predicati, per escludere i nodi non idonei) e \emph{Scoring} (o priorità, per classificare i nodi rimasti ed eleggere il vincitore).

I vincoli primari considerati dallo schedulatore includono:
\begin{itemize}[noitemsep]
    \item \textbf{Requisiti di risorse hardware}: verifica della disponibilità effettiva di CPU, memoria RAM ed eventuali acceleratori (GPU) rispetto alle soglie di \texttt{requests} indicate nel manifest del Pod;
    \item \textbf{Requisiti di servizio e porte}: verifica che le porte host richieste dal container non siano già occupate da altri processi sul nodo;
    \item \textbf{Policy hardware e software}: restrizioni legate a licenze, architettura della CPU (x86\_64 vs ARM64) o kernel patch specifici;
    \item \textbf{Affinità e Anti-Affinità dei Nodi (\emph{Node Affinity/Anti-Affinity})}: regole dichiarative per forzare o preferire l'assegnazione di un Pod a specifici nodi contrassegnati da etichette (\emph{labels}), ad esempio per collocare carichi pesanti su nodi dotati di dischi NVMe;
    \item \textbf{Affinità e Anti-Affinità tra Pod (\emph{Inter-Pod Affinity/Anti-Affinity})}: regole volte a massimizzare la co-locazione di microservizi ad alta interazione per minimizzare la latenza di rete (affinità), oppure volte a segregare le repliche del medesimo servizio su macchine fisiche o zone di disponibilità distinte per scongiurare disservizi generalizzati (anti-affinità);
    \item \textbf{Taints e Tolerations}: proprietà duale all'affinità dei nodi. Mentre l'affinità attrae i Pod verso un insieme di nodi, un \textbf{Taint} applicato a un nodo consente a quest'ultimo di respingere (\emph{repel}) l'allocazione di Pod. Un Pod può essere schedulato su un nodo marcato con un Taint solo ed esclusivamente se possiede una corrispondente \textbf{Toleration};
    \item \textbf{Località dei dati (\emph{Data Locality})}: predilezione di nodi in cui sono già presenti i volumi fisici o le immagini di container richieste, al fine di azzerare i tempi di trasferimento su rete;
    \item \textbf{Scadenze operative (\emph{Deadlines})}: priorità di esecuzione associate a Job o batch computation con tempi di completamento prefissati.
\end{itemize}

\subsubsection{4. Kube-Controller-Manager}
Il \textbf{kube-controller-manager} è un singolo processo binario che raggruppa ed esegue concorrentemente i cicli di retroazione e controllo (\emph{controllers}) nativi di Kubernetes. Ciascun controller monitora costantemente lo stato del cluster interrogando l'API server e applica correzioni continue affinché lo stato effettivo converga verso lo stato desiderato.

La tabella seguente illustra la suddivisione tassonomica dei controller, discriminando tra i controller nativi integrati nel core e gli add-on specializzati per estendere le funzionalità del cluster:

\begin{table}[H]
\centering
\small
\renewcommand{\arraystretch}{1.2}
\begin{tabularx}{\textwidth}{l p{2.2cm} X p{3.8cm}}
\toprule
\textbf{Nome Controller} & \textbf{Tipologia} & \textbf{Funzione Primaria} & \textbf{Esempi Gestiti} \\
\midrule
\texttt{ReplicaSet} & Built-in & Mantiene il numero esatto di repliche dei Pod & \texttt{rs-nginx}, \texttt{rs-redis}, \texttt{rs-mysql} \\
\texttt{Deployment} & Built-in & Gestisce rollout dichiarativi, rollback e ReplicaSet & \texttt{deploy-frontend}, \texttt{deploy-backend} \\
\texttt{DaemonSet} & Built-in & Garantisce l'esecuzione di una replica per nodo & \texttt{fluentd}, \texttt{node-exporter} \\
\texttt{StatefulSet} & Built-in & Gestisce Pod con identità stabile e ordine sequenziale & \texttt{postgres-cluster}, \texttt{kafka-cluster} \\
\texttt{Job} & Built-in & Esegue Pod fino al completamento con successo & \texttt{data-migration}, \texttt{image-scan} \\
\texttt{CronJob} & Built-in & Esegue Job schedulati periodicamente a tempo & \texttt{backup-job}, \texttt{report-job} \\
\texttt{Node} & Built-in & Monitora la reattività dei nodi e avvia l'eviction & \texttt{node-1}, \texttt{node-2}, \texttt{node-3} \\
\texttt{Service} & Built-in & Configura Load Balancer nativi su cloud provider & \texttt{my-service}, \texttt{api-svc}, \texttt{web-svc} \\
\texttt{EndpointSlice} & Built-in & Mantiene mapping scalabili IP/porta dei Pod per i Service & \texttt{endpoints per Service} \\
\texttt{Namespace} & Built-in & Gestisce il ciclo di vita e la pulizia dei namespace & \texttt{default}, \texttt{prod}, \texttt{staging}, \texttt{dev} \\
\texttt{PersistentVolume} & Built-in & Gestisce il binding e il provisioning tra PV e PVC & \texttt{pv-nfs}, \texttt{pv-ssd}, \texttt{pvc-db} \\
\texttt{ServiceAccount} & Built-in & Gestisce identità e token di accesso API per Pod & \texttt{default}, \texttt{build-robot}, \texttt{deployer} \\
\texttt{HPA} & Built-in & Scala i Pod in base al consumo di CPU/RAM & \texttt{hpa-frontend}, \texttt{hpa-backend} \\
\texttt{ClusterRole} & Built-in & Gestisce l'aggregazione di permessi RBAC & \texttt{view}, \texttt{edit}, \texttt{admin} \\
\midrule
\texttt{Ingress} & Add-on & Attua instradamento HTTP/HTTPS Layer 7 & \texttt{test-ingress}, \texttt{api-ingress} \\
\texttt{Cert Manager} & Add-on & Gestisce il rilascio e rinnovo di certificati TLS & \texttt{myapp-tls}, \texttt{api-cert} \\
\texttt{Metrics Server} & Add-on & Fornisce telemetria di risorse (CPU/RAM) per HPA & Metriche per pod/node \\
\texttt{ExternalDNS} & Add-on & Sincronizza record DNS esterni con i Service & \texttt{example.com} $\rightarrow$ IP \\
\texttt{Cluster Autoscaler} & Add-on & Ridimensiona i pool di nodi del cloud sottostante & Worker pool scaling \\
\texttt{Gateway} & Add-on & Implementa le specifiche evolute Gateway API (L4--L7) & \texttt{gateway}, \texttt{httproute} \\
\texttt{Istio} & Add-on & Amministra configurazioni di Service Mesh avanzate & \texttt{virtualservice}, \texttt{destinationrule} \\
\texttt{Prometheus Operator}& Add-on & Gestisce l'infrastruttura di monitoraggio distribuita & \texttt{servicemonitor}, \texttt{prometheusrule} \\
\bottomrule
\end{tabularx}
\caption{Quadro sinottico dei controller Kubernetes built-in ed estensioni add-on.}
\label{tab:controllers}
\end{table}

\subsection{I Nodi Lavoratori (Worker Nodes)}
I Worker Node rappresentano le entità operative delegate all'esecuzione fisica dei carichi di lavoro. Ciascun nodo ospita i seguenti moduli:

% [IMMAGINE DEL FILE DI SLIDE CON NOME "Orchestrazione_1"] Componenti esecutivi e di rete del Worker Node — Slide 17 / Pagina 17 del PDF

\subsubsection{1. Kubelet}
Il \textbf{kubelet} è l'agente principale in esecuzione su ogni nodo del cluster. Non gestisce container generati arbitrariamente al di fuori della piattaforma, ma garantisce rigorosamente che i container descritti all'interno dei \textbf{PodSpecs} assegnati siano attivi e conformi. Le mansioni primarie includono:
\begin{itemize}[noitemsep]
    \item Ricezione delle specifiche dei Pod inoltrate dal \texttt{kube-apiserver};
    \item Download e sincronizzazione dei segreti crittografici (\emph{Secrets}) e delle mappe di configurazione (\emph{ConfigMaps});
    \item Coordinamento del montaggio dei volumi fisici associati allo storage del Pod;
    \item Dialogo con il container runtime locale per pilotare il download delle immagini e l'avvio o la terminazione dei container;
    \item Rilevamento dello stato di salute dei nodi e dei singoli Pod, trasmettendo i report telemetrici al Control Plane;
    \item Esecuzione periodica dei controlli di integrità applicativa (\emph{Probes}):
    \begin{itemize}[noitemsep]
        \item \emph{Startup Probe}: determina l'avvenuta inizializzazione dell'applicazione, disattivando temporaneamente gli altri probe durante i lenti processi di avvio;
        \item \emph{Liveness Probe}: intercetta situazioni di deadlock o stallo interno dell'applicazione; in caso di fallimento prolungato, il kubelet termina il container applicando la policy di riavvio (\texttt{restartPolicy});
        \item \emph{Readiness Probe}: verifica se il container è abilitato a ricevere traffico di rete; in caso di fallimento, l'indirizzo IP del Pod viene temporaneamente rimosso dagli endpoint dei Service di bilanciamento.
    \end{itemize}
\end{itemize}

\subsubsection{2. Kube-Proxy}
Il \textbf{kube-proxy} è un demone di rete in esecuzione su ciascun worker node, incaricato di implementare l'astrazione di rete dei \textbf{Service}.
\begin{itemize}[noitemsep]
    \item Consente la comunicazione di rete tra Pod situati su nodi distinti o verso client esterni, mantenendo la coerenza logica della topologia distribuita;
    \item Monitora l'API server per intercettare la nascita o la distruzione di oggetti Service ed EndpointSlice;
    \item Riflette localmente le regole di inoltro dei pacchetti programmando gli stack di filtraggio del kernel Linux (principalmente mediante tabelle \texttt{iptables} o strutture ad alte prestazioni \texttt{IPVS}, oppure mediante una modalità legacy puramente applicativa in \emph{userspace});
    \item Realizza il bilanciamento di traffico Layer 4 (TCP e UDP) verso gli IP effettivi dei container, traducendo il \emph{ClusterIP} virtuale nell'indirizzo di uno specifico Pod.
\end{itemize}

\subsubsection{3. Container Runtime}
Il software esecutivo sottostante che isola le risorse computazionali e avvia i container (trattato dettagliatamente nella Sezione \ref{sec:cri}).

\subsubsection{4. Add-On Di Cluster E Servizi Ausiliari}
Gli Add-on sono componenti architetturali complementari distribuiti tipicamente sotto forma di Pod di sistema, essenziali per la piena funzionalità operativa:
\begin{itemize}[noitemsep]
    \item \textbf{Cluster DNS (CoreDNS)}: server DNS interno ad alta affidabilità che assegna automaticamente record di risoluzione dei nomi (record A, SRV) ai Service di Kubernetes. Tutti i container istanziati nel cluster vengono configurati automaticamente con il puntamento a CoreDNS nel proprio \texttt{/etc/resolv.conf}, permettendo la \emph{service discovery} nominale anziché vincolata a indirizzi IP mutevoli;
    \item \textbf{Componenti CNI (\emph{Container Network Interface})}: plugin di rete (quali Calico, Flannel o Weave) incaricati di allestire la rete overlay piatta (\emph{flat network}) che interconnette tutti i Pod del cluster;
    \item \textbf{Stack di Monitoraggio e Logging}: aggregazione di metriche tramite Prometheus e visualizzazione avanzata su dashboard Grafana per la telemetria continua;
    \item \textbf{Ingress Controllers}: router applicativi Layer 7 (ad es. NGINX Ingress Controller, Envoy, Traefik) che consentono l'instradamento basato su percorsi URL, terminazione TLS e bilanciamento avanzato del traffico HTTP/HTTPS proveniente dall'esterno del cluster.
\end{itemize}

\newpage

% ==============================================================================
% SEZIONE: IL CONTAINER RUNTIME E L'INTERFACCIA CRI
% ==============================================================================

\section{Il Container Runtime E L'Interfaccia CRI}
\label{sec:cri}

\subsection{Principio Di Disaccoppiamento E Standardizzazione}
Nelle versioni embrionali di Kubernetes, il supporto all'esecuzione dei container era vincolato rigidamente al motore monolitico Docker. Il codice sorgente del \texttt{kubelet} conteneva riferimenti espliciti e chiamate proprietarie all'API di Docker. Questa architettura monolitica violava i principi di modularità ingegneristica e precludeva l'adozione di container runtime alternativi orientati a contesti specifici (ad esempio runtime leggeri, sandbox di isolamento spinto o sistemi per hypervisor).

Per risolvere strutturalmente questa rigidità architetturale, la comunità di Kubernetes ha introdotto la \textbf{Container Runtime Interface (CRI)}. La CRI stabilisce un principio cardine: \textbf{Kubernetes deve essere completamente disaccoppiato dall'implementazione specifica del runtime}. La CRI si configura come un'interfaccia a plugin standardizzata che permette al \texttt{kubelet} di dialogare con molteplici runtime eterogenei senza alcuna necessità di ricompilare il binario del kubelet o modificare il Control Plane.

Tra le principali implementazioni conformi allo standard CRI figurano:
\begin{itemize}[noitemsep]
    \item \textbf{containerd} (e il plugin integrato \texttt{CRI-containerd}): il runtime industriale standard de-facto, derivato dallo scorporo del core di Docker;
    \item \textbf{CRI-O}: un runtime minimale conforme agli standard OCI (\emph{Open Container Initiative}), sviluppato specificamente per Kubernetes;
    \item \textbf{Docker Engine} (mediante shim di compatibilità dedicato, quale \texttt{dockershim} o \texttt{cri-dockerd});
    \item \textbf{Frakti}: implementazione per ambienti ad elevato isolamento che avvia container all'interno di micro-VM basate su hypervisor;
    \item Soluzioni storiche o sperimentali quali \textbf{rkt} (tramite l'adattatore \texttt{rktlet}).
\end{itemize}

\subsection{Il Protocollo gRPC E Le Operazioni Fondamentali}
La specifica CRI definisce un'architettura client-server basata interamente sul framework open-source ad alte prestazioni \textbf{gRPC} (\emph{Google Remote Procedure Call}), che adotta \textbf{Protocol Buffers} (Protobuf) come linguaggio neutro e compatto di serializzazione binaria dei dati.

In questa architettura locale di nodo:
\begin{itemize}[noitemsep]
    \item Il \textbf{kubelet} agisce come \textbf{Client gRPC};
    \item Il modulo di interfaccia del runtime (CRI Shim / Runtime Daemon) espone un \textbf{Server gRPC} in ascolto su una socket Unix locale (ad esempio \texttt{/run/containerd/containerd.sock}).
\end{itemize}

% [IMMAGINE DEL FILE DI SLIDE CON NOME "Orchestrazione_1"] Schema della comunicazione gRPC client-server tra Kubelet e CRI Shim — Slide 21 / Pagina 21 del PDF

Le interfacce remote esposte dal server CRI sono suddivise in due servizi logici:
\begin{enumerate}[noitemsep]
    \item \textbf{ImageService}: dedicato alla gestione delle immagini di sistema:
    \begin{itemize}[noitemsep]
        \item \texttt{PullImage}: effettua il download dell'immagine da un registry remoto pubblico o privato;
        \item \texttt{ListImages}: restituisce l'elenco delle immagini già presenti nella cache locale del nodo;
        \item \texttt{RemoveImage}: elimina un'immagine dal disco locale del nodo.
    \end{itemize}
    \item \textbf{RuntimeService}: preposto al governo del ciclo di vita dei container e delle sandbox di Pod:
    \begin{itemize}[noitemsep]
        \item \texttt{RunPodSandbox}: prepara l'ambiente di isolamento condiviso (crea il network namespace del Pod);
        \item \texttt{CreateContainer}: istanzia un container specificando parametri di configurazione, filesystem, variabili d'ambiente e montaggi;
        \item \texttt{StartContainer}: avvia il processo primario all'interno del container creato;
        \item \texttt{StopContainer}: notifica il segnale di terminazione (\texttt{SIGTERM} seguito da \texttt{SIGKILL} dopo un periodo di grazia) arrestando il container;
        \item \texttt{RemoveContainer}: distrugge definitivamente il container rilasciando le risorse associate.
    \end{itemize}
\end{enumerate}

\subsection{Workflow Dettagliato Di Creazione Di Un Container}
La sequenza operativa con cui il cluster trasforma una definizione dichiarativa in un processo in esecuzione segue un flusso deterministico rigoroso:

\begin{figure}[H]
\centering
\begin{tikzpicture}[
    node distance=1.3cm,
    block/.style={rectangle, draw=theoryblue, fill=backcolour, text width=12cm, rounded corners, inner sep=6pt, font=\small},
    arrow/.style={->, >=Latex, draw=theoryblue, thick}
]
\node[block] (step1) {\textbf{1. Ricezione PodSpec}: \texttt{kubelet} intercetta la notifica di schedulazione dall'API server e valida la specifica del Pod (immagini, variabili, mount point, porte).};
\node[block, below=of step1] (step2) {\textbf{2. Traduzione CRI}: \texttt{kubelet} serializza le direttive del manifest in una richiesta gRPC formattata in Protocol Buffers (ad esempio \texttt{CreateContainerRequest}).};
\node[block, below=of step2] (step3) {\textbf{3. Inoltro gRPC}: La richiesta viene inviata attraverso la socket Unix locale verso il servizio di runtime (containerd/CRI-O).};
\node[block, below=of step3] (step4) {\textbf{4. Esecuzione sul Runtime}: Il runtime scarica l'immagine (se non presente), allestisce i filesystem sovrapposti, associa il network namespace già preparato e configura le limitazioni cgroups.};
\node[block, below=of step4] (step5) {\textbf{5. Avvio e Risposta}: Il processo applicativo viene eseguito; il runtime restituisce un messaggio \texttt{CreateContainerResponse} confermando lo stato attivo.};
\node[block, below=of step5] (step6) {\textbf{6. Supervisione Telemetrica}: \texttt{kubelet} esegue periodicamente le probe di liveness/readiness, monitorando la salute del processo e notificando gli esiti all'API server.};

\draw[arrow] (step1) -- (step2);
\draw[arrow] (step2) -- (step3);
\draw[arrow] (step3) -- (step4);
\draw[arrow] (step4) -- (step5);
\draw[arrow] (step5) -- (step6);
\end{tikzpicture}
\caption{Workflow end-to-end di istanziazione e supervisione di un container mediante CRI.}
\label{fig:cri-workflow}
\end{figure}

\newpage

% ==============================================================================
% SEZIONE: I POD: L'UNITÀ ATOMICA DI ESECUZIONE
% ==============================================================================

\section{I Pod: L'Unità Atomica Di Esecuzione}

\subsection{Definizione Formale E Relazione Container-Pod-Node}
In Kubernetes, i container non vengono mai istanziati ed eseguiti direttamente sui nodi lavoratori in modo isolato. L'unità minima di calcolo, distribuzione e scheduling prende il nome di \textbf{Pod}.

% [IMMAGINE DEL FILE DI SLIDE CON NOME "Orchestrazione_1"] Relazione gerarchica tra multipli container all'interno di un Pod collocato sul Worker Node — Slide 7 / Pagina 7 del PDF

\begin{theorem}
Un \textbf{Pod} è l'astrazione fondamentale e l'atomo computazionale di Kubernetes. Esso rappresenta un gruppo co-locato di uno o più container che condividono in modo deterministico lo stack di rete, lo spazio di memoria condivisa (IPC) e i volumi di memorizzazione su disco.
\end{theorem}

L'architettura impone precise garanzie formali:
\begin{itemize}[noitemsep]
    \item \textbf{Integrità di co-locazione su singolo nodo}: qualora un Pod sia costituito da una molteplicità di container, la totalità di essi è garantita risiedere sul \textbf{medesimo nodo lavoratore}. Un singolo Pod non può in alcun caso estendersi o essere partizionato attraverso nodi fisici o virtuali distinti;
    \item \textbf{Stack di Rete Condiviso}: a ciascun Pod allocato viene assegnato un \textbf{indirizzo IP univoco all'interno del cluster}. Tutti i container appartenenti al medesimo Pod condividono lo stesso network namespace:
    \begin{itemize}[noitemsep]
        \item I container interni comunicano tra loro a bassissima latenza attraverso l'interfaccia di loopback locale (\texttt{localhost});
        \item Lo spazio delle porte di rete (porte TCP e UDP) è comune. Ne consegue che due container operanti nel medesimo Pod non possono effettuare il bind della stessa porta di rete (ad esempio, entrambi non possono occupare contemporaneamente la porta TCP 80);
    \end{itemize}
    \item \textbf{Filesystem e Volumi Condivisi}: i volumi di archiviazione dichiarati a livello di Pod possono essere montati contemporaneamente all'interno dei filesystem isolati dei vari container membri, consentendo la condivisione diretta di file e dati di stato;
    \item \textbf{Natura Effimera}: i Pod sono entità intrinsecamente volatili. In caso di guasto hardware del nodo o terminazione pianificata, il Pod non viene ripristinato con la medesima identità ma viene distrutto e ricreato \emph{ex-novo} (spesso con un nuovo identificatore e un nuovo indirizzo IP) ad opera dei controller di livello superiore (\texttt{Deployment}, \texttt{DaemonSet}, \texttt{Job}).
\end{itemize}

\subsection{Dichiarazione Di Un Pod Tramite Manifest YAML}
La definizione di un Pod avviene redigendo un file di configurazione strutturato in sintassi YAML, il quale rispetta lo schema ufficiale dell'API di Kubernetes. Di seguito viene riportato il manifest standard per la distribuzione di un web server NGINX:

\begin{lstlisting}[language=yaml, caption={Manifest dichiarativo per la creazione di un Pod semplice (\texttt{nginx.yml}).}]
apiVersion: v1
kind: Pod
metadata:
  name: nginx
  namespace: default
spec:
  containers:
  - name: nginx
    image: nginx
    ports:
    - containerPort: 80
\end{lstlisting}

L'analisi semantica del manifest evidenzia i campi strutturali:
\begin{itemize}[noitemsep]
    \item \texttt{apiVersion: v1}: dichiara la versione dello schema REST dell'API utilizzato per serializzare e validare l'oggetto;
    \item \texttt{kind: Pod}: formalizza la tipologia di risorsa che si intende istanziare;
    \item \texttt{metadata}: blocco di dati descrittivi che assegna un nome univoco all'oggetto (\texttt{name: nginx}) e il contesto di isolamento logico (\texttt{namespace: default});
    \item \texttt{spec}: definisce lo stato desiderato (\emph{Desired State}), includendo la lista dei container esecutivi (\texttt{spec.containers}), l'immagine da scaricare dal registro (\texttt{image: nginx}) e l'esposizione informativa della porta applicativa (\texttt{containerPort: 80}).
\end{itemize}

\subsection{Operatività Con la CLI kubectl E Diagnostica}
Lo strumento ufficiale da riga di comando per interagire con l'API server del cluster è \textbf{kubectl}. Il comando consente la gestione dell'intero ciclo di vita degli oggetti:

\begin{lstlisting}[language=bash, caption={Comandi operativi per la creazione e ispezione dello stato del Pod.}]
# Creazione della risorsa a partire dal file di specifica YAML
[root@controller ~]# kubectl create -f nginx.yml
pod/nginx created

# Verifica dello stato di allocazione, indirizzo IP e nodo assegnato
[root@controller ~]# kubectl get pods -o wide
NAME    READY   STATUS    RESTARTS   AGE   IP          NODE     ...
nginx   1/1     Running   0          43s   10.36.0.4   worker

# Ispezione analitica dei dettagli, eventi di sistema e probes
[root@controller ~]# kubectl describe pod nginx
\end{lstlisting}

Il comando \texttt{kubectl describe} riveste un ruolo diagnostico cruciale: aggrega in un unico output i metadati, lo stato dettagliato dei container, le condizioni dei probe di integrità e, soprattutto, l'elenco cronologico degli eventi (\emph{Events}) generati dal Control Plane (avvenuta schedulazione, pull dell'immagine, esecuzione di mount e avvio del processo).

\subsection{Accesso Diretto Ai Pod E Meccanismi Di Port-Forwarding}
In condizioni operative standard, l'indirizzo IP del Pod (\texttt{10.36.0.4} nell'esempio sopra riportato) appartiene a un piano di indirizzamento virtuale privato non direttamente instradabile dalla macchina operatore esterna al cluster. Se un amministratore cerca di effettuare una richiesta diretta:
\begin{lstlisting}[language=bash]
curl http://localhost:80
# Esito: Failed to connect to localhost port 80
\end{lstlisting}
La connessione fallisce poiché il servizio non è in ascolto sullo stack TCP/IP della macchina locale. Diversamente, se la chiamata viene invocata direttamente dall'interno del Pod:
\begin{lstlisting}[language=bash]
kubectl exec -it nginx -- curl http://localhost:80
\end{lstlisting}
La richiesta ha successo, in quanto eseguita all'interno del medesimo network namespace del server web.

Per finalità di collaudo, sviluppo o \emph{troubleshooting} di un Pod anomalo senza esporlo pubblicamente su Internet, Kubernetes mette a disposizione la funzionalità di **port-forwarding**:

\begin{lstlisting}[language=bash]
# Sintassi generale:
# kubectl port-forward <resource-type/resource-name> [local_port]:<pod_port>

kubectl port-forward pods/nginx 8080:80
\end{lstlisting}

\subsubsection{Meccanismo Operativo Del Port-Forwarding}
Il comando \texttt{kubectl port-forward} stabilisce una connessione HTTP/SPDY bidirezionale o WebSocket persistente tra la macchina client locale e il \texttt{kube-apiserver}. L'API server agisce da proxy inoltrando i byte verso l'agente \texttt{kubelet} residente sul nodo worker di destinazione, il quale instrada i pacchetti direttamente nella porta interna del Pod. In questo modo:
\begin{enumerate}[noitemsep]
    \item Il client locale apre una porta in ascolto (ad esempio la TCP \texttt{8080});
    \item Una chiamata locale verso tale porta:
    \begin{lstlisting}[language=bash]
curl http://localhost:8080
    \end{lstlisting}
    riceve risposta corretta dal server web interno al container;
    \item Tutto il flusso di rete è incapsulato e tunnellato in modo sicuro attraverso la connessione autenticata all'API server.
\end{enumerate}

La medesima astrazione di inoltro locale può essere applicata ad altri oggetti dell'API di Kubernetes:
\begin{lstlisting}[language=bash, caption={Esempi di port-forwarding su risorse di livello superiore.}]
# Port-forwarding verso un Deployment
kubectl port-forward deployment/my-deployment 8080:80

# Port-forwarding verso un ReplicaSet
kubectl port-forward replicaset/my-replicaset 8080:80

# Port-forwarding verso un'astrazione di Service
kubectl port-forward service/my-service 8080:80
\end{lstlisting}

\newpage

% ==============================================================================
% SEZIONE: GESTIONE DELLO STATO E KUBERNETES VOLUMES
% ==============================================================================

\section{Gestione Dello Stato E Kubernetes Volumes}

\subsection{La Criticità Dell'Effimerità Dei Container}
Nei motori di containerizzazione nativi, il filesystem associato a un'istanza è intrinsecamente effimero. Esso viene costruito sovrapponendo uno strato di lettura/scrittura (\emph{read-write thin layer}) al di sopra delle immagini immutabili a sola lettura. Questo design comporta una criticità architetturale fondamentale:
\begin{itemize}[noitemsep]
    \item Qualora il processo principale all'interno di un container subisca un arresto anomalo (crash) o il kubelet ne comandi il riavvio in seguito al fallimento di una liveness probe, \textbf{l'intero contenuto scritto sul filesystem locale dal container terminato viene irrimediabilmente perduto};
    \item Il container riavviato riparte dallo stato incontaminato dell'immagine base, rendendo impossibile la persistenza di basi di dati, file di sessione o log diagnostici;
    \item Due o più container eseguiti all'interno del medesimo Pod possiedono filesystem confinati e reciprocamente inaccessibili, impedendo lo scambio strutturato di file a livello locale.
\end{itemize}

Per ovviare a questo limite, Kubernetes definisce l'astrazione dei \textbf{Volumi} (\emph{Volumes}). A differenza di oggetti di primo livello quali i Pod o i nodi, i volumi non possono esistere in modo autonomo nel manifest di base, ma sono dichiarati all'interno della sezione \texttt{spec.volumes} della specifica del Pod e montati nei singoli container tramite la direttiva \texttt{volumeMounts}.

\subsection{Volumi Effimeri Vs Volumi Persistenti}
In merito al ciclo di vita dei dati rispetto all'infrastruttura computazionale, i volumi si suddividono in due categorie principali:
\begin{enumerate}[noitemsep]
    \item \textbf{Volumi Effimeri (\emph{Ephemeral Volumes})}:
    \begin{itemize}[noitemsep]
        \item Condividono strettamente il ciclo di vita del Pod che li dichiara (ad esempio la tipologia \texttt{emptyDir});
        \item Vengono allocati nel momento esatto in cui il Pod è avviato sul nodo e distrutti permanentemente all'atto della cancellazione del Pod;
        \item Tuttavia, \textbf{i dati contenuti nel volume persistono attraverso i riavvii dei singoli container} interni al Pod: se un container subisce un crash e viene riavviato dal kubelet, la nuova istanza trova i dati scritti precedentemente sul volume intatti.
    \end{itemize}
    \item \textbf{Volumi Persistenti (\emph{Persistent Volumes})}:
    \begin{itemize}[noitemsep]
        \item Risorse di storage fisiche esterne o gestite via rete (ad esempio SAN, NFS, dischi a blocchi cloud come AWS EBS, GCE Persistent Disk o Azure Disk);
        \item Sono del tutto indipendenti dal ciclo di vita del singolo Pod;
        \item Qualora il Pod venga terminato o rischedulato su un nodo fisico differente, il volume permane inalterato e può essere re-incapsulato e rimontato sul nuovo Pod mediante l'astrazione \texttt{persistentVolumeClaim} (PVC).
    \end{itemize}
\end{enumerate}

\subsection{Pattern Di Condivisione Dello Storage Tra Container}
La possibilità di montare simultaneamente il medesimo volume su container differenti appartenenti allo stesso Pod abilita pattern architetturali avanzati (quali il \emph{Sidecar pattern} o l'\emph{Ambassador pattern}).

% [IMMAGINE DEL FILE DI SLIDE CON NOME "Orchestrazione_1"] Architettura di condivisione di volumi multipli tra container WebServer, ContentAgent e LogRotator — Slide 29 / Pagina 29 del PDF

Si consideri uno scenario composto da tre container cooperanti in un unico Pod:
\begin{enumerate}[noitemsep]
    \item \texttt{WebServer}: demone HTTP che eroga documenti statici risiedenti nella cartella \texttt{/var/htdocs} e scrive i registri di accesso in \texttt{/var/logs};
    \item \texttt{ContentAgent}: processo delegato a generare o sincronizzare periodicamente contenuti web, salvando i file in \texttt{/var/html};
    \item \texttt{LogRotator}: processo asincrono incaricato di prelevare i registri da \texttt{/var/logs}, comprimerli ed eseguire rotazione o analisi.
\end{enumerate}

Senza l'impiego di volumi condivisi, ciascun container opererebbe in un ambiente isolato e i tre processi non potrebbero collaborare. L'orchestrazione risolve il problema dichiarando due distinti volumi a livello di Pod:
\begin{itemize}[noitemsep]
    \item \textbf{Volume \texttt{publicHtml}}: montato contemporaneamente nel container \texttt{WebServer} sul percorso \texttt{/var/htdocs} e nel container \texttt{ContentAgent} sul percorso \texttt{/var/html}. In questo modo, qualsiasi pagina generata da \texttt{ContentAgent} diviene istantaneamente visibile e servibile dal \texttt{WebServer};
    \item \textbf{Volume \texttt{logVol}}: montato sia in \texttt{WebServer} su \texttt{/var/logs} che in \texttt{LogRotator} sul proprio percorso di log. Il container \texttt{ContentAgent}, non avendo montato tale volume, non ha visibilità né accesso ai registri di traffico, preservando il principio del minimo privilegio.
\end{itemize}

\subsection{Tassonomia Delle Tipologie Di Volume Supportate}
Kubernetes supporta un ampio ventaglio di meccanismi di memorizzazione e integrazione back-end:

\begin{table}[H]
\centering
\small
\renewcommand{\arraystretch}{1.2}
\begin{tabularx}{\textwidth}{l l X}
\toprule
\textbf{Tipologia Volume} & \textbf{Storage Provider} & \textbf{Finalità Operativa E Caratteristiche} \\
\midrule
\texttt{emptyDir} & Localhost & Directory effimera vuota legata al ciclo di vita del Pod; utile per caching e condivisione intra-Pod. \\
\texttt{hostPath} & Localhost & Mappa direttamente una directory o file del filesystem del nodo worker dentro il Pod. \\
\texttt{glusterfs} & GlusterFS Cluster & Filesystem distribuito open-source multi-nodo con replica e fault-tolerance. \\
\texttt{downwardAPI} & Kubernetes Pod Info & Espone metadati di runtime del Pod ai container sotto forma di file di testo. \\
\texttt{nfs} & NFS Server & Permette la condivisione di storage di rete read-write concorrente tra molteplici Pod. \\
\texttt{awsElasticBlockStore}& AWS EBS & Volume a blocchi gestito su infrastruttura pubblica Amazon Web Services. \\
\texttt{gcePersistentDisk} & Google Compute Engine & Disco persistente nativo di Google Cloud Platform. \\
\texttt{azureDisk} & Microsoft Azure & Servizio di archiviazione a blocchi gestito su cloud Azure. \\
\texttt{projected} & Kubernetes Resources & Proietta congiuntamente Secret, ConfigMap e Downward API in una cartella unificata. \\
\texttt{secret} & Kubernetes Secret & Inietta credenziali, chiavi private o certificati TLS in memoria RAM protetta (\texttt{tmpfs}). \\
\texttt{vSphereVolume} & VMware vSphere & Montaggio di dischi virtuali VMDK in infrastrutture virtualizzate enterprise. \\
\texttt{gitRepo} & Git Repository & Clona un repository Git all'avvio del Pod (soluzione deprecata a favore di init-container). \\
\bottomrule
\end{tabularx}
\caption{Panoramica delle tipologie di volumi supportate nativamente in Kubernetes.}
\label{tab:volume-types}
\end{table}

\subsection{Esposizione Di Metadati Mediante Downward API}
Spesso un'applicazione containerizzata necessita di conoscere a runtime informazioni relative al contesto infrastrutturale in cui è immersa (ad esempio il proprio nome di Pod, il namespace di appartenenza, l'indirizzo IP assegnato o le quote limite di CPU/RAM) senza dover interrogare direttamente l'API server con credenziali privilegiate.

Kubernetes realizza questa astrazione mediante la \textbf{Downward API}, la quale permette di esporre tali metadati attraverso due canali distinti:
\begin{enumerate}[noitemsep]
    \item \textbf{Variabili d'ambiente}: le informazioni scalari vengono passate all'ambiente del container tramite la direttiva \texttt{valueFrom.fieldRef};
    \item \textbf{Volume files}: la Downward API genera un volume virtuale in sola lettura in cui i metadati del Pod o dei container vengono proiettati come file di testo (\emph{plain text files}). L'applicazione può dunque leggere lo stato dell'infrastruttura semplicemente aprendo un normale percorso nel proprio filesystem.
\end{enumerate}

\subsection{Esempi Notevoli Di Dichiarazione Di Volumi In YAML}

\subsubsection{1. Volume Effimero EmptyDir}
Viene generato vuoto all'avvio del Pod su storage transitorio (disco locale del nodo o RAM se configurato con \texttt{medium: Memory}). I dati persistono attraverso il riavvio dei singoli container ma vengono distrutti all'arresto del Pod.
\begin{lstlisting}[language=yaml, caption={Dichiarazione e montaggio di un volume \texttt{emptyDir}.}]
spec:
  containers:
  - name: my-app
    image: app-image
    volumeMounts:
    - name: data
      mountPath: /cache
  volumes:
  - name: data
    emptyDir: {}
\end{lstlisting}

\subsubsection{2. Volume HostPath}
Mappa direttamente una directory presente sul sistema operativo del nodo worker ospitante all'interno del container. È indicato principalmente per agenti di sistema o strumenti di monitoraggio (quali DaemonSet che raccolgono log di sistema in \texttt{/var/log}), ma presenta forti rischi di sicurezza e vincola rigidamente il Pod allo specifico nodo fisico.
\begin{lstlisting}[language=yaml, caption={Dichiarazione e montaggio di un volume \texttt{hostPath}.}]
spec:
  containers:
  - name: host-accessor
    image: inspector-image
    volumeMounts:
    - name: host-volume
      mountPath: /node-data
  volumes:
  - name: host-volume
    hostPath:
      path: /path/on/host
      type: Directory
\end{lstlisting}

\subsubsection{3. Volume ConfigMap}
Consente di disaccoppiare i parametri di configurazione applicativa dalle immagini dei container. La ConfigMap inietta configurazioni complesse o interi file di testo direttamente nel filesystem del container, evitando la necessità di ricompilare o ricostruire l'immagine al variare delle impostazioni.
\begin{lstlisting}[language=yaml, caption={Dichiarazione e montaggio di un volume basato su \texttt{configMap}.}]
spec:
  containers:
  - name: server
    image: nginx
    volumeMounts:
    - name: config-volume
      mountPath: /etc/nginx/conf.d
  volumes:
  - name: config-volume
    configMap:
      name: my-configmap
\end{lstlisting}

% AA Orchestrazione_1 [34-66]   ------------------------------------------------------------

\section{Gestione Avanzata Dello Storage: Tipologie Di Volume, PV, PVC E StorageClass}

\subsection{Tipologie Di Volume Per Ambienti Di Rete Ed Ecosistemi Cloud}
La gestione della persistenza in un'infrastruttura distribuita richiede meccanismi eterogenei capaci di adattarsi sia all'elaborazione locale, sia alla condivisione concorrente dei dati, sia all'integrazione con infrastrutture di tipo \emph{Infrastructure as a Service} (IaaS) private o pubbliche. Oltre ai volumi puramente effimeri o vincolati al nodo (quali \texttt{emptyDir} e \texttt{hostPath}), Kubernetes supporta nativamente un insieme articolato di driver di storage:

\begin{itemize}[noitemsep]
    \item \textbf{Secret}: Risorsa dedicata all'archiviazione e all'iniezione protetta di dati confidenziali (password, token di autenticazione, certificati TLS, chiavi crittografiche). A differenza delle \texttt{ConfigMap}, i dati memorizzati nei Secret possono essere iniettati come variabili d'ambiente oppure montati come file all'interno di volumi protetti gestiti in memoria volatile (\texttt{tmpfs}), garantendo che i dati sensibili non vengano mai scritti sul disco fisico del nodo lavoratore.
    \begin{lstlisting}[language=yaml, caption={Dichiarazione di un volume basato su \texttt{secret}.}]
volumes:
- name: secret-volume
  secret:
    secretName: my-secret
    \end{lstlisting}

    \item \textbf{NFS (Network File System)}: Consente di montare all'interno dei Pod una directory remota esportata da un server NFS. Trattandosi di un filesystem di rete a livello di file, esso supporta nativamente l'accesso in lettura e scrittura simultaneo da parte di molteplici Pod situati su nodi distinti del cluster (\emph{ReadWriteMany}).
    \begin{lstlisting}[language=yaml, caption={Dichiarazione di montaggio diretto di una risorsa \texttt{nfs}.}]
volumes:
- name: nfs-volume
  nfs:
    path: /path/to/nfs
    server: nfs-server.example.com
    \end{lstlisting}

    \item \textbf{Cinder (OpenStack)}: Rappresenta il servizio di storage a blocchi (\emph{Block Storage Service}) della piattaforma open-source OpenStack. Nelle infrastrutture cloud private basate su architettura congiunta Kubernetes--OpenStack, il plugin monta dinamicamente i volumi virtuali Cinder formattandoli con il filesystem richiesto (ad esempio \texttt{ext4}).
    \begin{lstlisting}[language=yaml, caption={Integrazione di un volume di blocco \texttt{cinder} di OpenStack.}]
volumes:
- name: cinder-volume
  cinder:
    volumeID: my-cinder-volume
    fsType: ext4
    \end{lstlisting}

    \item \textbf{Azure Disk E Azure File}: Nelle distribuzioni ospitate sull'infrastruttura Microsoft Azure, Kubernetes discrimina tra l'accesso a blocchi dedicato mediante \textbf{Azure Disk} (adatto per carichi di lavoro che richiedono un singolo scrittore con elevate prestazioni I/O) e la condivisione di filesystem gestiti tramite \textbf{Azure File} (basati su protocollo SMB/CIFS per accessi multi-nodo concorrenti).
    \begin{lstlisting}[language=yaml, caption={Configurazione di un volume managed \texttt{azureDisk}.}]
volumes:
- name: azure-disk
  azureDisk:
    diskName: myDisk
    diskURI: >-
      /subscriptions/.../resourceGroups/.../providers/Microsoft.Compute/disks/myDisk
    \end{lstlisting}

    \item \textbf{AWS EBS (Elastic Block Store)}: Il servizio di storage a blocchi gestito di Amazon Web Services. Consente l'aggancio (\emph{attach}) persistente di un volume EBS a una specifica istanza virtuale EC2 su cui è in esecuzione il nodo lavoratore.
    \begin{lstlisting}[language=yaml, caption={Dichiarazione di un volume di tipo \texttt{awsElasticBlockStore}.}]
volumes:
- name: ebs-volume
  awsElasticBlockStore:
    volumeID: vol-12345678
    fsType: ext4
    \end{lstlisting}
\end{itemize}

\subsection{Esempio Applicativo Di Storage Condiviso Tra Container Alpine}
Per osservare concretamente come due processi indipendenti possano comunicare mediante un volume comune, si consideri la distribuzione di un Pod multi-container basato su \textbf{Alpine Linux}. Alpine rappresenta una distribuzione Linux minimale, orientata alla sicurezza e basata sulla libreria di sistema \texttt{musl libc} e sulla suite \texttt{BusyBox}; grazie all'ingombro su disco ridotto a pochi megabyte, Alpine costituisce uno standard di riferimento per realizzare container leggeri e privi di complessità superflua.

\begin{lstlisting}[language=yaml, caption={Manifest di condivisione di un volume \texttt{emptyDir} (\texttt{shared-volume-emptydir.yml}).}]
apiVersion: v1
kind: Pod
metadata:
  name: shared-volume-emptyDir
spec:
  containers:
  - name: alpine1
    image: alpine
    command: ["/bin/sleep", "999999"]
    volumeMounts:
    - mountPath: /alpine1
      name: data
  - name: alpine2
    image: alpine
    command: ["/bin/sleep", "999999"]
    volumeMounts:
    - mountPath: /alpine2
      name: data
  volumes:
  - name: data
    emptyDir: {}
\end{lstlisting}

Nel manifest sopra riportato, i due container (\texttt{alpine1} e \texttt{alpine2}) vengono avviati eseguendo un comando di sleep a lungo termine per mantenerne lo stato attivo. Il volume di tipo \texttt{emptyDir}, contrassegnato con il nome logico \texttt{data}, viene montato nel filesystem di \texttt{alpine1} sotto il percorso assoluto \texttt{/alpine1} e nel filesystem di \texttt{alpine2} sotto \texttt{/alpine2}. Qualsiasi file creato o modificato da \texttt{alpine1} in \texttt{/alpine1} risulterà immediatamente leggibile da \texttt{alpine2} all'interno di \texttt{/alpine2}.

L'applicazione e la verifica diagnostica del ciclo di vita del Pod avvengono mediante i comandi:
\begin{lstlisting}[language=bash]
# Distribuzione del manifest dichiarativo
$ kubectl apply -f shared-volume-emptydir.yml

# Ispezione dello stato dei volumi e dei mount point associati ai container
$ kubectl describe pod shared-volume-emptydir
\end{lstlisting}

\subsection{Disaccoppiamento Dello Storage: Persistent Volumes (PV) E Persistent Volume Claims (PVC)}
Nelle architetture di livello industriale, specificare i parametri fisici dei dischi (indirizzi IP dei server NFS, identificativi dei volumi cloud EBS o Cinder) direttamente all'interno del manifest del Pod presenta gravi svantaggi ingegneristici:
\begin{itemize}[noitemsep]
    \item Viola il principio di portabilità: il manifest risulterebbe utilizzabile esclusivamente su una determinata infrastruttura o cloud provider;
    \item Viola la separazione delle responsabilità: lo sviluppatore software non dovrebbe conoscere la topologia di storage del data center, la quale è di competenza dell'amministratore dell'infrastruttura di cluster.
\end{itemize}

Per disaccoppiare formalmente la fornitura dello storage dal suo effettivo consumo, Kubernetes introduce la separazione tra due oggetti:

\begin{theorem}
Un \textbf{Persistent Volume (PV)} è una risorsa di storage a livello di intero cluster (\emph{cluster-wide resource}), approvvigionata manualmente da un amministratore di sistema o creata dinamicamente da uno \emph{StorageClass provisioner}. Esso incapsula i dettagli tecnologici e implementativi del supporto di memorizzazione fisico (capacità, tipologia di protocollo, percorsi fisici, modalità di accesso).
\end{theorem}

\begin{theorem}
Una \textbf{Persistent Volume Claim (PVC)} è una richiesta formale di storage emessa da un utente o da un'applicazione a livello di namespace (\emph{namespaced resource}). Essa astrae i dettagli implementativi specificando unicamente i vincoli operativi desiderati: capacità richiesta (es. 10 GiB) e modalità di accesso.
\end{theorem}

\subsubsection{Modalità Di Accesso E Ciclo Di Associazione (Binding)}
Le modalità di accesso (\texttt{accessModes}) formalizzano i vincoli con cui il supporto fisico può essere montato dai nodi:
\begin{itemize}[noitemsep]
    \item \textbf{ReadWriteOnce (RWO)}: Il volume può essere montato con permessi di lettura e scrittura da un singolo nodo lavoratore alla volta;
    \item \textbf{ReadOnlyMany (ROX)}: Il volume può essere montato contemporaneamente da molteplici nodi in modalità di sola lettura;
    \item \textbf{ReadWriteMany (RWX)}: Il volume può essere montato contemporaneamente da una pluralità di nodi con permessi sia di lettura che di scrittura (caratteristica tipica di filesystem distribuiti quali NFS o CephFS).
\end{itemize}

Il \textbf{PersistentVolume Controller} esegue un ciclo continuo di riconciliazione per individuare una corrispondenza biunivoca (\emph{1-to-1 binding}) tra le richieste espresse in una PVC e i PV disponibili:
\begin{enumerate}[noitemsep]
    \item \textbf{Binding Statico Automatico}: Si realizza quando un PV preesistente soddisfa o eccede i requisiti della PVC in termini di capacità e modalità di accesso. Se i parametri coincidono, il controller passa lo stato di entrambi a \texttt{Bound};
    \item \textbf{Binding Esplicito Forzato}: È possibile forzare l'associazione a uno specifico PV inserendo la direttiva \texttt{spec.volumeName: <nome-PV>} nel manifest della PVC.
\end{enumerate}

\begin{lstlisting}[language=yaml, caption={Esempio di manifest per un PersistentVolume e una corrispondente PersistentVolumeClaim.}]
# Definizione del Persistent Volume (PV) - Ambito Cluster-Wide
apiVersion: v1
kind: PersistentVolume
metadata:
  name: pv-example
spec:
  capacity:
    storage: 1Gi
  accessModes:
  - ReadWriteOnce
  hostPath:
    path: /mnt/data
---
# Definizione della Persistent Volume Claim (PVC) - Ambito Namespaced
apiVersion: v1
kind: PersistentVolumeClaim
metadata:
  name: pvc-example
spec:
  accessModes:
  - ReadWriteOnce
  resources:
    requests:
      storage: 1Gi
\end{lstlisting}

L'amministratore e gli sviluppatori possono monitorare lo stato delle risorse e l'avvenuto binding mediante i comandi:
\begin{lstlisting}[language=bash]
$ kubectl get pv
$ kubectl get pvc
\end{lstlisting}

\subsection{Provisioning Dinamico E Risorsa StorageClass}
Nelle infrastrutture su larga scala, il provisioning statico (in cui l'amministratore crea manualmente a priori decine di PV in attesa che le PVC li consumino) risulta inefficiente e scarsamente scalabile. La risorsa \textbf{StorageClass} consente di automatizzare interamente questo processo realizzando il \textbf{Dynamic Provisioning}: essa definisce quale driver o fornitore (\emph{provisioner}) debba essere invocato per istanziare fisicamente il volume su richiesta al momento della ricezione di una nuova PVC.

\begin{lstlisting}[language=yaml, caption={Definizione di una \texttt{StorageClass} per provisioning dinamico su Google Cloud Platform.}]
apiVersion: storage.k8s.io/v1
kind: StorageClass
metadata:
  name: fast-ssd
# Plugin o driver responsabile della creazione del volume:
provisioner: kubernetes.io/gce-pd
parameters:
  type: pd-ssd                         # SSD ad alte prestazioni
  replication-type: none               # Nessuna replica regionale
reclaimPolicy: Delete                  # Azione alla distruzione della PVC
allowVolumeExpansion: true             # Consente l'espansione del volume
# Differisce il provisioning fino allo scheduling sul nodo:
volumeBindingMode: WaitForFirstConsumer
\end{lstlisting}

L'analisi dei campi evidenzia:
\begin{itemize}[noitemsep]
    \item \texttt{metadata.name}: L'identificatore globale della classe (oggetto cluster-wide, non associato a un singolo namespace). Le PVC faranno riferimento a questo nome specificando il campo \texttt{storageClassName: fast-ssd};
    \item \texttt{provisioner}: Specifica il driver di storage incaricato di dialogare con le API sottostanti (ad es. \texttt{kubernetes.io/gce-pd} per GCP, \texttt{kubernetes.io/aws-ebs} per AWS, \texttt{kubernetes.io/azure-disk} per Azure, o \texttt{rancher.io/local-path} per distribuzioni locali);
    \item \texttt{parameters}: Parametri opachi trasmessi direttamente al provisioner (ad es. tipologia \texttt{pd-ssd} per unità SSD o \texttt{pd-standard} per dischi meccanici convenzionali);
    \item \texttt{volumeBindingMode}: L'impostazione \texttt{WaitForFirstConsumer} è fondamentale in contesti distribuiti: indica al provisioner di non istanziare immediatamente il disco fisico alla creazione della PVC, bensì di attendere che il Pod richiedente sia stato assegnato a un nodo lavoratore dal \texttt{kube-scheduler}. Ciò garantisce che il volume venga creato nella medesima zona geografica o sul medesimo rack del nodo in cui verrà eseguito il carico di lavoro, evitando conflitti topologici;
    \item \texttt{reclaimPolicy}: Definisce il comportamento del volume fisico sottostante qualora la PVC corrispondente venga eliminata:
    \begin{itemize}[noitemsep]
        \item \textbf{Delete} (comportamento predefinito nei cloud provider): Il volume fisico (ad es. il disco virtuale EBS o GCE-PD) viene cancellato irrevocabilmente con conseguente distruzione definitiva dei dati memorizzati;
        \item \textbf{Retain}: Il PersistentVolume fisico non viene distrutto e transita nello stato \texttt{Released}, consentendo all'amministratore di recuperare manualmente i dati o procedere alla bonifica controllata.
    \end{itemize}
\end{itemize}

Ambienti di sviluppo locali (quali Docker Desktop o Minikube) forniscono configurazioni predefinite di StorageClass pronte all'uso:
\begin{lstlisting}[language=bash]
# Ispezione delle StorageClass registrate nel cluster
# kubectl get storageclass
NAME                 PROVISIONER              RECLAIMPOLICY   VOLUMEBINDINGMODE      ALLOWVOLUMEEXPANSION   AGE
hostpath             rancher.io/local-path    Delete          WaitForFirstConsumer   false                  10d
standard (default)   rancher.io/local-path    Delete          WaitForFirstConsumer   false                  10d
\end{lstlisting}

\newpage

% ==============================================================================
% SEZIONE: I NAMESPACE DI KUBERNETES
% ==============================================================================

\section{I Namespace Di Kubernetes: Partizionamento Logico E Sicurezza}

\subsection{Definizione Formale Ed Esigenze Di Partizionamento}
All'interno di un cluster Kubernetes condiviso tra una pluralità di utenti, team di sviluppo o ambienti operativi (sviluppo, collaudo, produzione), sorge l'esigenza di organizzare e confinare le risorse senza dover sostenere l'onere economico e manutentivo derivante dall'allestimento di cluster fisici separati.

\begin{theorem}
Un \textbf{Namespace} in Kubernetes è un meccanismo di isolamento logico che partiziona le risorse del cluster in gruppi distinti e indipendenti. Esso funge da contesto di visibilità entro il quale i nomi degli oggetti devono essere univoci, consentendo a namespace differenti di ospitare risorse aventi il medesimo identificatore.
\end{theorem}

L'utilizzo dei namespace abilita:
\begin{itemize}[noitemsep]
    \item \textbf{Multi-Tenancy logica}: Ripartizione sicura del cluster tra progetti eterogenei o clienti differenti;
    \item \textbf{Quota delle Risorse (\emph{Resource Quotas})}: Applicazione di vincoli per limitare la quantità cumulativa di risorse computazionali (limite globale di core CPU, gigabyte di RAM) o il numero massimo di oggetti (Pod, Service, PVC) istanziabili in ciascun contesto;
    \item \textbf{Delega Amministrativa e RBAC}: Concessione di permessi di accesso granulari circoscritti a uno specifico namespace (ad esempio consentire a un team di sviluppatori di operare esclusivamente nel namespace \texttt{dev} precludendo l'accesso al namespace \texttt{prod}).
\end{itemize}

\subsection{Distinzione Fondamentale Tra Isolamento Logico E Di Rete}
\begin{center}
\fcolorbox{theoryblue}{backcolour}{
\begin{minipage}{0.96\textwidth}
\textbf{Avvertenza Architetturale Critica: Isolamento Logico $\neq$ Isolamento di Rete} \\
Sebbene i namespace isolino rigorosamente i manifest e gli oggetti interrogabili tramite l'API server, \textbf{essi non introducono di per sé alcuna segregazione di rete a livello di traffico IP tra i Pod in esecuzione}.
\end{minipage}}
\end{center}

Il modello di connettività fondamentale di Kubernetes impone che ciascun Pod possa contattare direttamente qualsiasi altro Pod presente nel cluster, indipendentemente dal nodo di esecuzione o dal namespace di appartenenza. Di conseguenza, qualora il plugin CNI adottato non applichi policy restrittive, un'applicazione in esecuzione nel namespace \texttt{app} che conosca l'indirizzo IP di un Pod collocato nel namespace \texttt{prod} può instaurare una connessione TCP/HTTP arbitraria verso di esso. Per ottenere un isolamento di rete effettivo tra namespace è indispensabile configurare esplicitamente oggetti di tipo \textbf{NetworkPolicy} supportati da un provider CNI idoneo (ad esempio Calico).

\subsection{I Quattro Namespace Predefiniti Di Sistema}
All'atto dell'inizializzazione del cluster, il Control Plane genera automaticamente quattro namespace strutturali:

\begin{enumerate}[noitemsep]
    \item \textbf{\texttt{default}}: Il contesto operativo predefinito assegnato a tutte le risorse create qualora l'utente o il manifest non specifichino esplicitamente un namespace alternativo;
    \item \textbf{\texttt{kube-system}}: Riservato a tutti i carichi di lavoro e Pod infrastrutturali del Control Plane e del networking interno (demone CoreDNS, istanze di \texttt{kube-proxy}, plugin CNI);
    \item \textbf{\texttt{kube-public}}: Un namespace speciale dotato di permessi di lettura globali non autenticati (\emph{world-readable}). È concepito per condividere metadati di bootstrapping del cluster o configurazioni che devono risultare visibili a tutti i client prima ancora dell'autenticazione;
    \item \textbf{\texttt{kube-node-lease}}: Namespace dedicato a ospitare gli oggetti di tipo \textbf{NodeLease} per ciascun nodo lavoratore.
\end{enumerate}

\subsubsection{Approfondimento: Gli Heartbeat Ottimizzati Del Kube-Node-Lease}
Nei cluster di grandi dimensioni, monitorare lo stato di salute di migliaia di nodi lavoratori costituisce un collo di bottiglia prestazionale per l'API server e per \texttt{etcd}. Nel modello tradizionale, ciascun agente \texttt{kubelet} trasmetteva periodicamente un intero oggetto \texttt{NodeStatus} all'API server per notificare di essere in vita. Tale operazione risultava computazionalmente onerosa, poiché l'oggetto conteneva metadati complessi, liste di capacità e immagini locali, saturando il canale dati.

Per ovviare a questo problema di scalabilità, Kubernetes adotta l'interfaccia \textbf{Lease API}:
\begin{itemize}[noitemsep]
    \item Ciascun nodo possiede un corrispondente oggetto compatto denominato \texttt{Lease} collocato nel namespace \texttt{kube-node-lease};
    \item Il segnale di heartbeat inviato periodicamente dal \texttt{kubelet} si traduce in un aggiornamento leggero del solo timestamp contenuto nel campo \texttt{spec.renewTime};
    \item Il \textbf{Node Controller} monitora tali scadenze per rilevare tempestivamente il fallimento di un nodo: se il tempo trascorso dall'ultimo rinnovo eccede la soglia limite, il nodo viene marcato come non raggiungibile (\texttt{NotReady}) e il controller avvia le procedure di sgombero ed eliminazione (\emph{eviction}) dei Pod ivi collocati;
    \item Questo approccio riduce drasticamente l'occupazione di banda e il carico di scrittura su \texttt{etcd}.
\end{itemize}

\subsection{Gestione Pratica E Comandi Operativi}
L'interazione con i namespace mediante la CLI \texttt{kubectl} può avvenire in modalità sia imperativa che dichiarativa:

\begin{lstlisting}[language=bash, caption={Ispezione ed estrazione dei dettagli dei namespace.}]
# Visualizzazione dei namespace attivi nel cluster
[root@controller ~]# kubectl get ns
NAME              STATUS   AGE
default           Active   14d
kube-node-lease   Active   14d
kube-public       Active   14d
kube-system       Active   14d

# Ispezione di tutte le risorse presenti nell'intero cluster
[root@controller ~]# kubectl get all --all-namespaces
NAMESPACE     NAME                                   READY   STATUS    RESTARTS   AGE
kube-system   pod/coredns-f9fd979d6-nmsq5            1/1     Running   4          14d
default       service/kubernetes                     ClusterIP 10.96.0.1 <none>   443/TCP   14d
\end{lstlisting}

Per definire un namespace tramite approccio dichiarativo, è possibile determinare le coordinate di versione interrogando l'API server:
\begin{lstlisting}[language=bash]
# Verifica della risorsa tramite espressione regolare
[root@controller ~]# kubectl api-resources | grep -iE 'namespace|KIND'
NAME         SHORTNAMES   APIVERSION   NAMESPACED   KIND
namespaces   ns           v1           false        Namespace

# Ispezione della specifica strutturale
[root@controller ~]# kubectl explain Namespace | head -n 2
KIND:     Namespace
VERSION:  v1
\end{lstlisting}

Si procede quindi alla stesura del file di configurazione \texttt{app-ns.yml}:
\begin{lstlisting}[language=yaml, caption={Manifest dichiarativo per la creazione di un namespace (\texttt{app-ns.yml}).}]
apiVersion: v1
kind: Namespace
metadata:
  name: app
\end{lstlisting}

\begin{lstlisting}[language=bash, caption={Confronto tra approccio dichiarativo e imperativo.}]
# Modalita' Dichiarativa (consigliata per tracciamento in version control)
[root@controller ~]# kubectl create -f app-ns.yml
namespace/app created

# Modalita' Imperativa diretta
[root@controller ~]# kubectl create ns dev
namespace/dev created

# Ispezione analitica del namespace
[root@controller ~]# kubectl describe ns default

# Estrazione della definizione formale in formato YAML
[root@controller ~]# kubectl get ns default -o yaml
\end{lstlisting}

\newpage

% ==============================================================================
% SEZIONE: PATTERN AVANZATI DI COMPOSIZIONE DEI POD
% ==============================================================================

\section{Pattern Avanzati Di Composizione Dei Pod}

\subsection{Richiamo Al Vincolo Di Co-Locazione Multi-Container}
Come evidenziato in precedenza, l'adozione del modello multi-container è soggetta a un invariante architetturale rigido:
\begin{center}
\textbf{Tutti i container appartenenti al medesimo Pod sono vincolati a essere schedulati ed eseguiti sullo stesso identico Worker Node. Un Pod non può mai essere partizionato tra due nodi differenti.}
\end{center}

% [IMMAGINE DEL FILE DI SLIDE CON NOME "Orchestrazione_1"] Schema con divieto assoluto: un Pod multi-container non può essere diviso tra Node 1 e Node 2 — Slide 11 / Pagina 33 del PDF

Nella maggior parte dei casi operativi, si adotta il modello convenzionale \textbf{One-Container-per-Pod}, in cui il Pod funge da semplice wrapper gestionale attorno a un unico processo applicativo. Tuttavia, quando molteplici processi sono fortemente interdipendenti (\emph{tightly coupled}) e richiedono la condivisione dello stack di rete o di spazi di memoria locale, essi vengono incapsulati in un unico Pod multi-container.

In tale ambito, l'ingegneria del software cloud-native distingue tre pattern fondamentali:
\begin{enumerate}[noitemsep]
    \item \textbf{Sidecar Container}: Un container secondario ausiliario che potenzia o estende la logica del container primario senza alterarne il codice applicativo. Si occupa tipicamente di compiti di supporto quali estrazione e inoltro di log verso collector centralizzati, calcolo di metriche telemetriche o proxy di sincronizzazione dati;
    \item \textbf{Ambassador Container}: Agisce come proxy o intermediario per conto del container applicativo primario verso il mondo esterno. Il container principale invia il traffico verso \texttt{localhost} e l'ambassador si occupa di mascherare la complessità di rete, instradando le chiamate verso servizi esterni, gestendo il partizionamento su database con sharding, o attuando logiche di service discovery nell'ambito di una \emph{Service Mesh};
    \item \textbf{Adapter Container}: Funge da componente di conversione e normalizzazione per superare incompatibilità tra protocolli, formati o ambienti. Riceve dati generati dal container principale con interfacce eterogenee o proprietarie (ad es. log legacy formattati arbitrariamente) e li trasforma in un formato conforme agli standard di consumo del cluster (es. endpoint conforme a Prometheus).
\end{enumerate}

\subsection{Caso Di Studio E Laboratorio: Il Pattern Sidecar}
Si illustra un'applicazione completa del pattern Sidecar in cui un container primario produce dati diagnostici e un container secondario espone tali dati tramite interfaccia web su rete locale:

\begin{lstlisting}[language=yaml, caption={Manifest di implementazione del pattern Sidecar (\texttt{create-sidecar.yml}).}]
apiVersion: v1
kind: Pod
metadata:
  name: sidecar-pod
spec:
  volumes:
  - name: logs
    emptyDir: {}
  containers:
  - name: app
    image: busybox
    command: ["/bin/sh"]
    args: ["-c", "while true; do date >> /var/log/date.txt; sleep 10; done"]
    volumeMounts:
    - name: logs
      mountPath: /var/log
  - name: sidecar
    image: centos/httpd
    ports:
    - containerPort: 80
    volumeMounts:
    - name: logs
      mountPath: /var/www/html
\end{lstlisting}

\subsubsection{Dinamica Interna Del Manifest}
\begin{itemize}[noitemsep]
    \item Viene istanziato un volume effimero \texttt{logs} di tipo \texttt{emptyDir: \{\}}, la cui sintassi sintetica indica l'assenza di vincoli opzionali sulla dimensione;
    \item Il container primario \texttt{app} esegue un loop continuo in bash (\texttt{sh -c}) che appende il timestamp di sistema formattato al file \texttt{/var/log/date.txt} a intervalli regolari di dieci secondi;
    \item Il container secondario \texttt{sidecar} (un server web Apache basato su immagine \texttt{centos/httpd}) monta lo stesso identico volume all'interno della directory radice dei documenti web (\texttt{/var/www/html}) esponendo la porta TCP 80;
    \item Poiché entrambi i processi condividono il medesimo filesystem del volume, il file generato da \texttt{app} diviene accessibile tramite protocollo HTTP su porta 80 all'indirizzo \texttt{http://localhost/date.txt}.
\end{itemize}

\subsubsection{Collaudo E Verifica Telemetrica}
È possibile verificare il funzionamento collegandosi in modalità interattiva al container sidecar ed eseguendo una richiesta locale:
\begin{lstlisting}[language=bash, caption={Accesso al container sidecar e ispezione dei dati generati.}]
# Connessione al container specifico "sidecar" all'interno del Pod multi-container
[root@controller ~]# kubectl exec -it sidecar-pod -c sidecar -- /bin/bash

# Esecuzione di una richiesta HTTP verso il web server locale su localhost:80
[root@sidecar-pod ~]# curl http://localhost/date.txt
Fri Nov 27 05:43:00 UTC 2020
Fri Nov 27 05:43:10 UTC 2020
Fri Nov 27 05:43:20 UTC 2020
Fri Nov 27 05:43:30 UTC 2020
Fri Nov 27 05:43:40 UTC 2020
Fri Nov 27 05:43:50 UTC 2020
Fri Nov 27 05:44:00 UTC 2020
Fri Nov 27 05:44:11 UTC 2020
Fri Nov 27 05:44:21 UTC 2020
\end{lstlisting}
L'output dimostra chiaramente che il file scritto dal container primario viene servito con successo dall'istanza HTTP del container sidecar.

\newpage

% ==============================================================================
% SEZIONE: I CONTROLLER DI KUBERNETES E LE RISORSE DI CARICO DI LAVORO
% ==============================================================================

\section{I Controller Di Kubernetes E Le Risorse Di Carico Di Lavoro}

\subsection{Teoria Del Controllo Retroazionato Negli Oggetti Kubernetes}
Nell'ingegneria dell'automazione e della robotica, un ciclo di controllo (\emph{control loop}) è un processo non-terminante che monitora continuativamente le variabili di stato di un impianto e attua azioni correttive al fine di minimizzare l'errore rispetto a un riferimento, analogamente a un termostato che regola la temperatura di un ambiente termico:

% [IMMAGINE DEL FILE DI SLIDE CON NOME "Orchestrazione_1"] Diagramma del ciclo di controllo retroazionato tra stato desiderato e stato attuale — Slide 24 / Pagina 45 del PDF

Gli **Oggetti Kubernetes** costituiscono entità persistenti che rappresentano la configurazione e lo stato del cluster. Ciascun oggetto contiene convenzionalmente due sezioni logiche fondamentali:
\begin{itemize}[noitemsep]
    \item \texttt{spec} (\emph{Specification}): Rappresenta lo \textbf{Stato Desiderato} prescritto dall'amministratore (quali immagini impiegare, numero di repliche, porte aperte, vincoli di tolleranza ai guasti);
    \item \texttt{status} (\emph{Status}): Rappresenta lo \textbf{Stato Attuale} osservato in tempo reale dai controller mediante l'API server.
\end{itemize}

I controller integrati in \texttt{kube-controller-manager} attuano ciclicamente la riconciliazione tra \texttt{status} e \texttt{spec}, aggiornando al contempo i campi descrittivi degli oggetti stessi (ad esempio, quando un calcolo batch termina con successo, il Job Controller muta lo stato del Job marcandolo formalmente come \texttt{Finished}).

I manifest possono essere generati a partire da risorse già caricate nel cluster mediante il comando:
\begin{lstlisting}[language=bash]
kubectl get <object-type>/<object-name> -o yaml
\end{lstlisting}
consentendo all'operatore di impiegare la risorsa estratta come modello dichiarativo da applicare successivamente tramite \texttt{kubectl apply -f}.

\subsection{I Job: Esecuzione Di Carichi Di Lavoro Con Terminazione Finita}
I Pod ordinari generati da controller di servizio sono progettati per rimanere in esecuzione continuativa (\emph{run-forever}). Tuttavia, molti carichi computazionali (operazioni di backup periodico, migrazione di schemi di database, calcoli scientifici per batch processing o scansioni di vulnerabilità) possiedono una durata finita. Per questi scenari si impiega la risorsa **Job**.

Un Job garantisce che uno o più Pod vengano istanziati e continuino a operare fino a quando un numero prefissato di essi non termina con esito positivo (codice di uscita zero).

\subsubsection{Policy Di Riavvio Ammesse}
Per un Pod subordinato a un Job, la direttiva \texttt{restartPolicy} ammette unicamente due valori (il valore \texttt{Always} è categoricamente vietato):
\begin{itemize}[noitemsep]
    \item \textbf{\texttt{OnFailure}}: Se il processo termina con codice di errore (non-zero exit code), il kubelet tenta di riavviare il container all'interno del \textbf{medesimo Pod};
    \item \textbf{\texttt{Never}}: In caso di fallimento o crash, il container non viene riavviato sul posto; il Job controller si occupa di creare un \textbf{nuovo Pod sostitutivo}, preservando il Pod fallito per finalità di ispezione diagnostica (\emph{post-mortem analysis}).
\end{itemize}

\subsubsection{Tipologie Di Job E Parametri Di Concorrenza}
Le dinamiche di esecuzione dei Job sono governate da due parametri interi definiti in \texttt{spec}:
\begin{itemize}[noitemsep]
    \item \texttt{completions} ($n$): Il numero totale di esecuzioni con successo richieste per dichiarare completato il Job;
    \item \texttt{parallelism} ($m$): Il numero massimo di Pod che possono operare simultaneamente in parallelo.
\end{itemize}

Qualora tali parametri non vengano dichiarati esplicitamente, assumono entrambi il valore predefinito pari a 1 ($n=1, m=1$). Si distinguono tre modelli operativi:
\begin{enumerate}[noitemsep]
    \item \textbf{Job Non-Paralleli}: $n=1, m=1$. Viene istanziato un singolo Pod. Se termina con successo il Job è concluso; in caso di errore, il Pod viene rieseguito fino a completamento;
    \item \textbf{Job Paralleli a Conteggio Fisso}: $n > 1, m \ge 1$. Il Job istanzia batch concorrenti composti da al più $m$ Pod, arrestandosi solo quando $n$ Pod hanno terminato con successo;
    \item \textbf{Job Paralleli Con Coda Di Lavoro (\emph{Work Queue})}: $n=1, m > 1$. Vengono avviati $m$ Pod concorrenti che attingono task da una coda esterna comune. Nel momento in cui il primo Pod notifica il completamento positivo dell'intera coda, tutti i Pod rimasti terminano e il Job si conclude.
\end{enumerate}

La pulizia automatica dei Pod generati da Job terminati può essere demandata al controller integrato \textbf{TTL-after-finished}, che distrugge le risorse storiche al superare di un intervallo di tempo (\emph{Time-to-Live}) configurato tramite \texttt{spec.ttlSecondsAfterFinished}.

\begin{lstlisting}[language=yaml, caption={Esempio di manifest per un Job semplice con policy OnFailure (\texttt{job1.yml}).}]
apiVersion: batch/v1
kind: Job
metadata:
  name: hello
spec:
  template:
    spec:
      containers:
      - name: hello
        image: busybox
        command: ['sh', '-c', 'echo "Hello, Kubernetes!" && sleep 30']
      restartPolicy: OnFailure
\end{lstlisting}

\begin{lstlisting}[language=yaml, caption={Job sequenziale con tre completamenti successivi (\texttt{pod-simple-job.yml}).}]
apiVersion: batch/v1
kind: Job
metadata:
  name: pod-simple-job
spec:
  completions: 3
  template:
    spec:
      containers:
      - name: sleepy
        image: alpine
        command: ["/bin/sleep"]
        args: ["5"]
      restartPolicy: Never
\end{lstlisting}

Al termine dell'esecuzione, l'ispezione evidenzia l'avvenuto completamento di tutti i passi previsti:
\begin{lstlisting}[language=bash]
[root@controller ~]# kubectl get jobs
NAME             COMPLETIONS   DURATION   AGE
pod-simple-job   3/3           35s        5m7s
\end{lstlisting}

\subsection{ReplicationController E ReplicaSet}

\subsubsection{Architettura Del ReplicationController}
Il **ReplicationController** (RC) rappresenta il controller storico di prima generazione preposto a garantire che un numero specificato di repliche di un determinato Pod sia sempre attivo e operativo all'interno del cluster. Se un nodo si disconnette o un Pod viene espulso (\emph{evicted}), l'RC rileva immediatamente l'inconsistenza e alloca una nuova replica su un nodo sano.

Un ReplicationController è composto strutturalmente da tre elementi cardine:
\begin{enumerate}[noitemsep]
    \item \textbf{Label Selector}: Stabilisce i criteri di selezione con cui il controller riconosce i Pod che ricadono sotto la sua supervisione;
    \item \textbf{Replica Count}: Intero indicante il numero desiderato di istanze che devono essere mantenute attive;
    \item \textbf{Pod Template}: La specifica formale utilizzata per fabbricare nuove repliche nel caso in cui il numero effettivo risulti inferiore al riferimento desiderato.
\end{enumerate}

% [IMMAGINE DEL FILE DI SLIDE CON NOME "Orchestrazione_1"] Ciclo di riconciliazione decisionale del ReplicationController — Slide 24 / Pagina 45 del PDF

La logica decisionale del ciclo di riconciliazione implementa il seguente algoritmo:
\begin{itemize}[noitemsep]
    \item Se $\text{Repliche Attuali} < \text{Repliche Desiderate}$: Istanzia nuovi Pod applicando il \emph{Pod Template};
    \item Se $\text{Repliche Attuali} > \text{Repliche Desiderate}$: Seleziona ed elimina i Pod eccedenti;
    \item Se $\text{Repliche Attuali} = \text{Repliche Desiderate}$: Permane in stato di quiescenza monitorando gli eventi.
\end{itemize}

\begin{lstlisting}[language=yaml, caption={Dichiarazione classica di un \texttt{ReplicationController} (\texttt{replication-controller.yml}).}]
apiVersion: v1
kind: ReplicationController
metadata:
  name: myapp-rc
  labels:
    app: myapp
    type: dev
spec:
  replicas: 3
  selector:
    app: myapp
  template:
    metadata:
      name: myapp-pod
      labels:
        app: myapp
        type: dev
    spec:
      containers:
      - name: nginx-container
        image: nginx
\end{lstlisting}

Il ridimensionamento orizzontale (\emph{horizontal scaling}) di un ReplicationController può essere attuato a runtime in modo imperativo senza interrompere l'erogazione del servizio:
\begin{lstlisting}[language=bash]
# Espansione a 6 repliche concorrenti
[root@controller ~]# kubectl scale rc myapp-rc --replicas=6
replicationcontroller/myapp-rc scaled

# Contrazione a 3 repliche
[root@controller ~]# kubectl scale rc myapp-rc --replicas=3
replicationcontroller/myapp-rc scaled
\end{lstlisting}

\subsubsection{L'Evoluzione A ReplicaSet E La Potenza Dei Set-Based Selectors}
Nelle attuali versioni di Kubernetes, il ReplicationController è formalmente **deprecato** e integralmente sostituito dal **ReplicaSet** (RS). A livello comportamentale i due controller perseguono la medesima finalità, ma il ReplicaSet introduce una capacità espressiva notevolmente superiore nella definizione dei selettori di etichette:
\begin{itemize}[noitemsep]
    \item \textbf{Selettori basati su uguaglianza (\emph{Equality-Based Selectors})}: Supportati sia da RC che da RS. Consentono unicamente il filtraggio basato su corrispondenza esatta (ad esempio: $\texttt{app} = \texttt{myapp}$ oppure $\texttt{tier} \neq \texttt{frontend}$). Un ReplicationController non è in grado di raggruppare sotto la medesima gestione Pod recanti etichette eterogenee;
    \item \textbf{Selettori basati su insiemi (\emph{Set-Based Selectors})}: Prerogativa esclusiva del ReplicaSet. Consentono di filtrare le etichette valutando l'appartenenza a insiemi discreti mediante il costrutto \texttt{matchExpressions}.
\end{itemize}

Gli operatori validi supportati da \texttt{matchExpressions} sono:
\begin{enumerate}[noitemsep]
    \item \textbf{\texttt{In}}: Il valore associato alla chiave dell'etichetta deve appartenere obbligatoriamente alla lista specificata (es. \texttt{values: ["myapp", "backend"]});
    \item \textbf{\texttt{NotIn}}: Il valore dell'etichetta non deve comparire nell'insieme dei valori specificati;
    \item \textbf{\texttt{Exists}}: Il Pod deve obbligatoriamente contenere l'etichetta specificata nel campo \texttt{key}, a prescindere dal valore ad essa assegnato. Con questo operatore, il campo \texttt{values} non deve essere fornito;
    \item \textbf{\texttt{DoesNotExist}}: Il Pod non deve presentare alcuna etichetta avente la chiave specificata nel campo \texttt{key}. Anche in questo caso il campo \texttt{values} deve rimanere non valorizzato.
\end{enumerate}

\begin{lstlisting}[language=yaml, caption={ReplicaSet avanzato basato su selettori \texttt{matchExpressions} (\texttt{replica-set2.yml}).}]
apiVersion: apps/v1
kind: ReplicaSet
metadata:
  name: myapp-replicaset2
  labels:
    app: myapp
    type: dev
spec:
  replicas: 3
  selector:
    matchExpressions:
    - key: app
      operator: In
      values:
      - myapp
  template:
    metadata:
      name: myapp-pod
      labels:
        app: myapp
        type: dev
    spec:
      containers:
      - name: nginx-container
        image: nginx
\end{lstlisting}

\subsubsection{Vincolo Fondamentale Di Coerenza Sul Template E Adozione Dei Pod}
La specifica di un ReplicaSet impone un vincolo strutturale rigido:
\begin{center}
\textbf{Le etichette dichiarate nel blocco \texttt{spec.template.metadata.labels} devono obbligatoriamente soddisfare tutte le condizioni espresse nel blocco \texttt{spec.selector}.}
\end{center}
Qualora il template del Pod definisca etichette che non corrispondono al selettore, l'API server rifiuta la creazione della risorsa oppure il ReplicaSet genera un errore critico bloccandosi con zero repliche attive (ispezionabile tramite \texttt{kubectl describe}). Questo meccanismo previene cicli infiniti di instanziazione in cui il controller creerebbe Pod per poi non riconoscerli come propri.

Inoltre, un ReplicaSet è capace di praticare l'\textbf{adozione di Pod preesistenti} (\emph{Adoption}): se nel cluster esistono già Pod orfani le cui etichette soddisfano il selettore, il ReplicaSet ne assume immediatamente il controllo contandoli nel proprio bilancio di repliche, senza ricrearli.

\subsection{StatefulSet: Gestione Di Carichi Di Lavoro Con Stato}

\subsubsection{Il Limite Dei ReplicaSet Con Applicazioni Statefull}
I ReplicaSet e i Deployment sono progettati per microservizi totalmente privi di stato (\emph{stateless}). In tali contesti:
\begin{itemize}[noitemsep]
    \item Le repliche sono rigorosamente interscambiabili, fungibili e prive di memoria storica;
    \item Ciascun Pod riceve un identificatore alfanumerico casuale (ad es. \texttt{myapp-deployment-7b4876b5-xyt6z});
    \item Gli indirizzi IP interni sono volatili e non prevedibili, richiedendo sempre l'interposizione di un Service per il bilanciamento.
\end{itemize}

Tale approccio collassa qualora si orchestrino cluster distribuiti dotati di stato (\emph{Stateful Workloads}), quali database relazionali (PostgreSQL), cluster transazionali distribuite (Kafka, ZooKeeper, Elasticsearch, Cassandra). Tali sistemi richiedono tassativamente:
\begin{enumerate}[noitemsep]
    \item \textbf{Identità di rete stabile e deterministica};
    \item \textbf{Associazione biunivoca e durevole tra la specifica istanza e il proprio storage fisico};
    \item \textbf{Ordine sequenziale rigoroso} nelle fasi di bootstrap, arresto e aggiornamento progressivo.
\end{enumerate}

Per soddisfare questi requisiti si impiega la risorsa **StatefulSet**.

% [IMMAGINE DEL FILE DI SLIDE CON NOME "Orchestrazione_1"] Confronto strutturale tra Pod anonimi generati da ReplicaSet e Pod ordinati (A-0, A-1, A-2) gestiti da StatefulSet — Slide 29 / Pagina 50 del PDF

\subsubsection{Identità Ordinale E Risoluzione DNS Con Headless Service}
A differenza dei ReplicaSet, ciascun Pod generato da uno StatefulSet possiede un'\textbf{identità persistente} (\emph{sticky identity}):
\begin{itemize}[noitemsep]
    \item A ciascun Pod viene assegnato un indice ordinale intero a base zero, compreso nell'intervallo $[0, N-1]$ (con $N$ pari al numero di repliche desiderate);
    \item Il nome del Pod e il rispettivo hostname vengono derivati deterministicamente concatenando il nome dello StatefulSet e l'indice: ad esempio, per uno StatefulSet denominato \texttt{myapp-statefulset}, i Pod generati saranno rigorosamente denominati:
    \[
    \texttt{myapp-statefulset-0}, \quad \texttt{myapp-statefulset-1}, \quad \texttt{myapp-statefulset-2}
    \]
    \item Se il nodo su cui opera il Pod \texttt{myapp-statefulset-0} si guasta, la nuova istanza rischedulata sul nuovo nodo assumerà \textbf{esattamente lo stesso identico nome} (\texttt{myapp-statefulset-0}).
\end{itemize}

Poiché l'indirizzo IP del nuovo Pod rischedulato può mutare, lo StatefulSet impone tassativamente l'associazione a un **Headless Service** (dichiarato tramite la direttiva \texttt{serviceName} e configurato con \texttt{clusterIP: None}). Il CoreDNS del cluster registra record DNS di tipo \texttt{A} statici dedicati a ciascuna replica individuale:
\[
\texttt{<nome-pod>.<nome-servizio>.<namespace>.svc.cluster.local}
\]
Ad esempio:
\[
\texttt{myapp-statefulset-0.myapp-service.default.svc.cluster.local}
\]
Anche qualora l'indirizzo IP cambi a seguito di una rischedulazione, il record DNS nominale viene aggiornato istantaneamente dal Control Plane, consentendo ai membri di un cluster quorum (ad es. leader ed eletti) di comunicare senza interruzione e senza ambiguità.

\subsubsection{Gestione Dello Storage Con volumeClaimTemplates}
La caratteristica più rilevante dello StatefulSet risiede nel meccanismo dei **volumeClaimTemplates**:

% [IMMAGINE DEL FILE DI SLIDE CON NOME "Orchestrazione_1"] Flusso di conservazione e associazione dello storage al fallimento e riavvio di una replica in uno StatefulSet — Slide 30 / Pagina 51 del PDF

A differenza di un Pod ordinario che condivide la medesima definizione di volume con le proprie repliche, lo StatefulSet instanzia dinamicamente una \textbf{PVC dedicata e indipendente per ciascun indice ordinale}, denominata secondo lo schema:
\[
\texttt{<nome-claim>-<nome-statefulset>-<indice>}
\]
Ad esempio: \texttt{data-myapp-statefulset-0}, \texttt{data-myapp-statefulset-1}, ecc.

Se la replica con indice 1 fallisce o viene rischedulata su un altro server fisico:
\begin{itemize}[noitemsep]
    \item Il volume persistente ad essa associato \textbf{non viene distrutto né riassegnato};
    \item La nuova istanza \texttt{myapp-statefulset-1} viene automaticamente ricollegata alla PVC preesistente \texttt{data-myapp-statefulset-1}, preservando lo stato dei dati in modo trasparente.
\end{itemize}

\begin{lstlisting}[language=yaml, caption={Manifest strutturato di uno StatefulSet con volumeClaimTemplates (\texttt{stateful-set2.yml}).}]
apiVersion: apps/v1
kind: StatefulSet
metadata:
  name: myapp-statefulset
  labels:
    app: myapp
    type: dev
spec:
  replicas: 3
  selector:
    matchLabels:
      app: myapp
  serviceName: "myapp-service"          # Associazione a Headless Service obbligatoria
  template:
    metadata:
      name: myapp-pod
      labels:
        app: myapp
        type: dev
    spec:
      containers:
      - name: nginx-container
        image: nginx
        volumeMounts:
        - name: data
          mountPath: /usr/share/nginx/html
  volumeClaimTemplates:                # Generazione dinamica di PVC dedicate per replica
  - metadata:
      name: data
    spec:
      accessModes: ["ReadWriteOnce"]
      resources:
        requests:
          storage: 10Gi
      storageClassName: standard       # Delega il provisioning dinamico alla StorageClass
\end{lstlisting}

\subsubsection{Limitazioni Architetturali E Politiche Di Gestione}
L'amministrazione degli StatefulSet presenta vincoli ingegneristici precisi a tutela dell'integrità dei dati:
\begin{itemize}[noitemsep]
    \item \textbf{Nessuna cancellazione automatica dei volumi}: La cancellazione di uno StatefulSet o la sua riduzione di scala (\emph{scaling down}) \textbf{non elimina in alcun caso le PersistentVolumeClaim o i volumi fisici associati}. Questa scelta progettuale garantisce la sicurezza del patrimonio informativo, evitando purghing accidentali a seguito di errori di battitura;
    \item \textbf{Pulizia manuale dei dati}: Per bonificare definitivamente lo storage a valle della rimozione di uno StatefulSet, l'operatore deve cancellare esplicitamente le relative claim:
    \begin{lstlisting}[language=bash]
kubectl delete pvc data-myapp-statefulset-0 \
  data-myapp-statefulset-1 \
  data-myapp-statefulset-2
    \end{lstlisting}
    \item \textbf{Garanzie di ordinamento (\emph{Pod Management Policy})}: Con la politica predefinita \texttt{OrderedReady}, i Pod vengono creati rigorosamente in ordine sequenziale crescente (da $0$ a $N-1$) attendendo che la replica precedente sia pienamente attiva (\emph{Ready}) prima di avviare la successiva. In fase di spegnimento o contrazione, l'arresto procede in senso inverso rigoroso (da $N-1$ discendendo verso $0$).
\end{itemize}
% AA Orchestrazione_1 [67-99]   ------------------------------------------------------------

\section{Rete E Identità Nei Workload Stateful: Headless Service}

Nei sistemi distribuiti orchestrati da Kubernetes, la gestione delle applicazioni si divide concettualmente tra carichi di lavoro \emph{stateless} (privi di stato persistente proprio) e carichi \emph{stateful} (in cui ogni istanza mantiene uno stato locale, dati su disco o un ruolo architetturale ben definito all'interno di un cluster, come nei nodi master/replica di database relazionali o nei sistemi di consenso quorum-based quali Raft o Paxos).

Per comprendere l'esigenza di un **Headless Service**, è necessario rinfrescare brevemente il meccanismo di instradamento di rete offerto da un canonico \textbf{Service} di tipo \texttt{ClusterIP}. In un Service tradizionale, il control plane alloca un indirizzo IP virtuale stabile (\emph{Virtual IP} o VIP) registrato nel componente \textbf{kube-proxy}. Quando un client interno al cluster risolve il nome DNS del Service, il DNS di cluster (CoreDNS) restituisce unicamente tale VIP. Il traffico inviato al VIP viene quindi intercettato a livello di kernel Linux (tramite regole di \texttt{iptables} o tabelle hash di \textbf{IPVS}) e distribuito in modo casuale o \emph{round-robin} verso uno dei Pod sottostanti (i cui endpoint sono mantenuti aggiornati nel registro di cluster). Questo comportamento è ideale per applicazioni disaccoppiate e fungibili, ma risulta distruttivo per carichi di lavoro gestiti da un **StatefulSet**, nei quali il client (o un nodo peer) non deve interagire con un'istanza qualsiasi, bensì con un Pod specifico e univoco.

\subsection{Principio Di Funzionamento Di Un Headless Service}

Uno **StatefulSet** assegna a ciascun Pod un'identità ordinale deterministica e immutabile (ad esempio \texttt{myapp-statefulset-0}, \texttt{myapp-statefulset-1}, \texttt{myapp-statefulset-2}), preservata anche in caso di riavvio o rischedulazione del container su un nodo fisico differente. Tuttavia, affinché questa identità univoca sia raggiungibile attraverso lo stack di rete, Kubernetes richiede che lo StatefulSet sia associato a un Service privo di ClusterIP, definito formalmente **Headless Service**.

Un Headless Service viene dichiarato esplicitamente impostando all'interno della specifica il campo:
\begin{lstlisting}[language=yaml]
spec:
  clusterIP: None
\end{lstlisting}

Configurando \texttt{clusterIP: None}, il control plane disattiva l'assegnazione dell'IP virtuale e, di conseguenza, \textbf{kube-proxy} non crea alcuna regola di instradamento o bilanciamento del carico nel kernel del nodo worker. Il compito di instradare il traffico viene interamente demandato al server DNS interno al cluster. Invece di restituire un singolo VIP condiviso, il server DNS crea e mantiene record di tipo \textbf{A} (e \textbf{SRV}) per ciascun singolo Pod associato al selettore del Service.

La cooperazione architetturale tra le due entità è vincolante:
\begin{itemize}[noitemsep]
    \item Lo **StatefulSet** garantisce la stabilità dei nomi e la persistenza dei volumi associati;
    \item L'**Headless Service** fornisce l'infrastruttura DNS necessaria a risolvere tali identità in indirizzi IP fisici diretti.
\end{itemize}

Senza il Service di supporto, i Pod dello StatefulSet possiederebbero nomi stabili all'interno del database dello scheduler, ma non esisterebbe alcun meccanismo per risolverli in modo deterministico via rete.

\subsection{Risoluzione DNS Diretta E Manifest Di Configurazione}

Quando l'Headless Service è attivo, ogni Pod appartenente allo StatefulSet ottiene un record DNS fully qualified (\emph{FQDN}) strutturato secondo la convenzione gerarchica:
\[
\texttt{<nome-pod>.<nome-servizio>.<namespace>.svc.cluster.local}
\]

All'interno dello stesso namespace, la risoluzione può avvenire in forma abbreviata mediante la sintassi:
\[
\texttt{<nome-pod>.<nome-servizio>}
\]

Di seguito si riporta il manifest YAML che definisce l'Headless Service a servizio dello StatefulSet:

\begin{lstlisting}[language=yaml]
# Headless service necessario per lo StatefulSet
apiVersion: v1
kind: Service
metadata:
  name: myapp-service
spec:
  clusterIP: None # Disabilita il Virtual IP (servizio headless)
  selector:
    app: myapp    # Deve corrispondere esattamente al selettore dello StatefulSet
  ports:
    - port: 80
      targetPort: 80
\end{lstlisting}

Grazie a tale configurazione, da qualsiasi carico di lavoro interno al cluster è possibile invocare direttamente una specifica istanza mediante chiamate puntuali a livello applicativo:
\begin{lstlisting}[language=bash]
curl http://myapp-statefulset-0.myapp-service:80
curl http://myapp-statefulset-1.myapp-service:80
curl http://myapp-statefulset-2.myapp-service:80
\end{lstlisting}

Ciascuna richiesta HTTP raggiungerà direttamente l'interfaccia di rete del Pod target, abilitando protocolli peer-to-peer, repliche di database primario-secondario e comunicazioni inter-nodo a bassa latenza senza intermediari di proxying.

\newpage

% ==============================================================================
% SEZIONE 2: IL CONTROLLER DAEMONSET: GESTIONE DEI SERVIZI DI NODO
% ==============================================================================

\section{Il Controller DaemonSet: Gestione Dei Servizi Di Nodo}

Nei sistemi cloud-native si presenta frequentemente la necessità di eseguire servizi di infrastruttura a livello host, ovvero componenti software che non scalano in base al carico applicativo degli utenti finali, ma devono presenziare in maniera ubiqua su ogni singolo nodo worker del cluster. Esempi canonici comprendono demoni di memorizzazione distribuita (quali Ceph o GlusterFS), agenti di monitoraggio delle metriche hardware (come Prometheus Node Exporter) o demoni di raccolta ed esportazione dei log di sistema (come Fluentd o Logstash).

Per rispondere a questo requisito architetturale, Kubernetes introduce il controller **DaemonSet**.

\subsection{Ciclo Di Vita E Differenze Rispetto Ai ReplicaSet}

La responsabilità primaria del controller DaemonSet consiste nel garantire che una copia di un determinato Pod sia in esecuzione su tutti i nodi del cluster (oppure su un sottoinsieme selezionato di essi). Il ciclo di vita dei Pod gestiti da un DaemonSet è rigidamente vincolato al ciclo di vita dei nodi sottostanti:
\begin{enumerate}[noitemsep]
    \item **Aggiunta di un nodo**: non appena un nuovo nodo entra a far parte del cluster e passa allo stato \texttt{Ready}, il controller DaemonSet pianifica e avvia immediatamente un'istanza del Pod sul nuovo host.
    \item **Rimozione di un nodo**: quando un nodo viene rimosso dal cluster (o posto in \emph{drain}), i Pod associati vengono avviati alla \emph{garbage collection} e deallocati.
    \item **Cancellazione del DaemonSet**: eliminando la risorsa DaemonSet, il control plane effettua la cancellazione atomica e la bonifica di tutti i Pod generati sui nodi.
\end{enumerate}

La differenza concettuale tra un **ReplicaSet** e un **DaemonSet** risiede nel grado di accoppiamento con l'infrastruttura sottostante:
\begin{itemize}[noitemsep]
    \item Il **ReplicaSet** modella applicazioni completamente disaccoppiate dalla topologia dei nodi. L'obiettivo dell'algoritmo di schedulazione consiste nel mantenere costante il numero cardinale $N$ di repliche desiderate; lo scheduler sceglie liberamente i nodi in base a criteri di bilanciamento delle risorse, potendo allocare più copie del medesimo Pod su un singolo nodo e lasciandone altri privi di istanze.
    \item Il **DaemonSet** modella applicazioni vincolate all'infrastruttura hardware o di sistema operativo. Non esiste un conteggio arbitrario di repliche specificabile dall'utente: il numero di Pod attivi è sempre proporzionale al numero di nodi ammissibili, con una e una sola istanza allocata per ciascun nodo.
\end{itemize}

% [IMMAGINE DEL FILE DI SLIDE CON NOME "Orchestrazione_1"] Schema comparativo dell'allocazione dei Pod tra ReplicaSet (repliche multiple concentrate o distribuite su vincoli di carico) e DaemonSet (esattamente una replica per ciascun nodo del cluster) — Slide 56 / Pagina 3 del PDF

\subsection{Schedulazione Su Sottoinsiemi Di Nodi}

Sebbene l'impostazione predefinita preveda la copertura totale dei nodi del cluster, è possibile limitare l'esecuzione del DaemonSet a un sottoinsieme ristretto di macchine. Ciò si ottiene definendo all'interno del Pod template la proprietà \texttt{nodeSelector}, la quale effettua una valutazione di corrispondenza rispetto alle etichette (\emph{labels}) assegnate ai singoli nodi (oppure avvalendosi di chiavi standard del runtime, come l'hostname del nodo). Questa funzionalità è essenziale in cluster eterogenei in cui, ad esempio, determinati agenti di telemetria devono risiedere unicamente sui nodi dotati di acceleratori hardware (GPU) o interfacce di rete dedicate a bassa latenza.

\subsection{Analisi Del Manifest Di Un DaemonSet: Raccolta Log Con Fluentd}

Il seguente listato illustra la configurazione reale di un DaemonSet deputato alla raccolta dei log dei container tramite l'agente **Fluentd**, con montaggio diretto di percorsi del file system dell'host tramite volumi \texttt{hostPath}:

\begin{lstlisting}[language=yaml]
apiVersion: apps/v1
kind: DaemonSet
metadata:
  name: fluentd
  labels:
    app: fluentd
spec:
  selector:
    matchLabels:
      app: fluentd
  template:
    metadata:
      labels:
        app: fluentd
    spec:
      containers:
        - name: fluentd
          image: fluent/fluentd:v0.14.10
          resources:
            limits:
              memory: 200Mi
            requests:
              cpu: 100m
              memory: 200Mi
          volumeMounts:
            - name: varlog
              mountPath: /var/log
            - name: varlibdockercontainers
              mountPath: /var/lib/docker/containers
              readOnly: true
      terminationGracePeriodSeconds: 30
      volumes:
        - name: varlog
          hostPath:
            path: /var/log
        - name: varlibdockercontainers
          hostPath:
            path: /var/lib/docker/containers
\end{lstlisting}

I volumi di tipo \texttt{hostPath} collegano la directory host `/var/log` e `/var/lib/docker/containers` (quest'ultima montata rigorosamente in modalità di sola lettura \texttt{readOnly: true} per prevenire corruzioni accidentali) all'interno dell'ambiente isolato del container, consentendo a Fluentd di aggregare i flussi standard di output di tutti i processi in esecuzione sul nodo.

I comandi operativi associati al ciclo di vita del componente sono:
\begin{lstlisting}[language=bash]
# Distribuzione dichiarativa del DaemonSet
kubectl apply -f ./fluentd-daemonset.yml

# Rimozione del DaemonSet e bonifica automatica dei pod
kubectl delete daemonset fluentd-elasticsearch --namespace=kube-system
\end{lstlisting}

\subsection{Gestione Dichiarativa Vs Imperativa: \texttt{kubectl apply} Vs \texttt{kubectl create}}

La gestione delle risorse in Kubernetes supporta due paradigmi operativi distinti:
\begin{itemize}[noitemsep]
    \item **Approccio Imperativo (\texttt{kubectl create})**: indica al sistema \emph{cosa fare} nell'immediato. Il comando invia una richiesta diretta di creazione all'API Server; se la risorsa specificata nel file manifest esiste già all'interno dello stato del cluster, l'operazione fallisce sollevando un errore di collisione.
    \item **Approccio Dichiarativo (\texttt{kubectl apply})**: definisce lo \emph{stato desiderato} del sistema. Il comando applica una strategia di riconciliazione basata su un algoritmo di unione a tre vie (\emph{three-way merge patch}) che confronta la configurazione corrente in esecuzione sul cluster, la configurazione registrata nell'ultima annotazione applicata e il nuovo file manifest locale. Se l'oggetto non esiste, viene creato; se esiste già, viene aggiornato in modo trasparente senza interrompere le proprietà immutate.
\end{itemize}

\subsection{Ciclo Di Terminazione Ed Esecuzione Dei Lifecycle Hooks}

La corretta terminazione dei processi in un ambiente distribuito è cruciale per evitare interruzioni di servizio o corruzioni dello stato transiente. Il parametro:
\[
\texttt{terminationGracePeriodSeconds: 30}
\]
(il cui valore predefinito è di 30 secondi) stabilisce l'intervallo temporale massimo concesso a un Pod per completare le operazioni di \emph{graceful shutdown} prima che il runtime forzi la chiusura dei processi tramite l'invio del segnale non intercettabile \texttt{SIGKILL}.

Kubernetes arricchisce la gestione del ciclo di vita esponendo due tipologie di **Lifecycle Hooks**:
\begin{itemize}
    \item **PostStart**: viene invocato immediatamente dopo l'avvio del container. Tale hook viene eseguito in modalità asincrona rispetto all'\texttt{ENTRYPOINT} del container (non vi è garanzia di quale dei due venga eseguito per primo). Tuttavia, la chiamata all'hook handler all'interno del contesto del Pod è sincrona: se il gestore del \texttt{PostStart} si blocca o entra in stallo, il Pod non raggiungerà mai lo stato \texttt{Running} e rimarrà bloccato nella fase di inizializzazione. All'hook non possono essere passati parametri da riga di comando.
    \item **PreStop**: viene invocato immediatamente prima che il container venga terminato a seguito di una richiesta esplicita all'API (es. rischedulazione, aggiornamento del controller) o di eventi di gestione (fallimento dei probe di vitalità, preemption da parte di Pod a priorità più elevata, saturazione delle risorse del nodo).
\end{itemize}

La gestione temporale del \texttt{PreStop} è rigorosamente sequenziale e sincrona: il segnale operativo di arresto convenzionale (\texttt{SIGTERM}) **non viene inviato** fino a quando l'esecuzione dell'hook \texttt{PreStop} non è terminata con successo. L'intero budget di tempo a disposizione per la terminazione è vincolato dal periodo di grazia $T_{\text{grace}} = \texttt{terminationGracePeriodSeconds}$.

Indicando con $T_{\text{hook}}$ la durata di esecuzione dell'hook \texttt{PreStop} e con $T_{\text{stop}}$ il tempo impiegato dal processo applicativo per arrestarsi fisiologicamente a valle della ricezione del segnale \texttt{SIGTERM}, la condizione di terminazione non traumatica è definita dalla disuguaglianza:
\[
\boxed{T_{\text{hook}} + T_{\text{stop}} \le T_{\text{grace}}}
\]

Qualora l'hook entri in loop infinito o la somma dei tempi ecceda la soglia configurata (ad esempio, ponendo $T_{\text{grace}} = \SI{60}{\second}$, $T_{\text{hook}} = \SI{55}{\second}$ e $T_{\text{stop}} = \SI{10}{\second}$, per un totale di $\SI{65}{\second} > \SI{60}{\second}$), il kernel del container runtime interviene allo scadere dei $\SI{60}{\second}$ uccidendo immediatamente il processo mediante \texttt{SIGKILL}, troncando ogni transazione applicativa ancora pendente.

\newpage

% ==============================================================================
% SEZIONE 3: ESECUZIONE PROGRAMMATA CON I CRONJOB
% ==============================================================================

\section{Esecuzione Programmata Con I CronJob}

Mentre i Pod ordinari gestiti da ReplicaSet o DaemonSet modellano servizi di lunga durata (\emph{long-running services}) progettati per rimanere costantemente in ascolto su una porta di rete, una frazione rilevante dei carichi computazionali aziendali consiste in elaborazioni a lotti (\emph{batch processing}), quali backup di database, aggregazioni periodiche di telemetria, invio di reportistica o pulizia di archivi temporanei.

Nei sistemi operativi UNIX-like, tali attività vengono storicamente demandate al demone di sistema \textbf{cron}. In Kubernetes, questo paradigma viene formalizzato tramite due controller integrati: la risorsa base **Job** (che coordina l'esecuzione di Pod con ciclo di vita a completamento finito, riprovando l'avvio in caso di guasto fino al successo del processo) e la risorsa sovraordinata **CronJob**, che gestisce la pianificazione temporale dei Job stessi.

\subsection{Sintassi Dello Scheduling Temporale}

Il controller CronJob adotta la sintassi standard del file `/etc/crontab`. L'espressione di pianificazione è articolata su cinque campi posizionali distinti, valutati da sinistra verso destra:
\[
\underbrace{\text{minuto}}_{(0-59)} \quad \underbrace{\text{ora}}_{(0-23)} \quad \underbrace{\text{giorno del mese}}_{(1-31)} \quad \underbrace{\text{mese}}_{(1-12)} \quad \underbrace{\text{giorno della settimana}}_{(0-6 \text{ o } 7, 0=\text{Dom})}
\]

Si consideri a titolo esemplificativo la riga di comando Linux:
\begin{lstlisting}[language=bash]
0 5 * * 1 tar -zcf /var/backups/home.tgz /home/
\end{lstlisting}

La decodifica formale dei singoli parametri evidenzia quanto segue:
\begin{itemize}[noitemsep]
    \item **Minuto (\texttt{0})**: specifica che l'esecuzione deve avviarsi al minuto esatto zero dell'ora;
    \item **Ora (\texttt{5})**: individua le ore 5:00 del mattino (il formato orario impiega la notazione h24);
    \item **Giorno del mese (\texttt{*})**: il metacarattere asterisco funge da wildcard ("qualsiasi valore del dominio ammissibile"), svincolando l'esecuzione dalla data solare;
    \item **Mese (\texttt{*})**: wildcard che estende la validità a tutti i dodici mesi dell'anno solare;
    \item **Giorno della settimana (\texttt{1})**: specifica l'esecuzione nel primo giorno della settimana lavorativa, corrispondente al lunedì (essendo lo 0 e il 7 convenzionalmente associati alla domenica).
\end{itemize}
L'intera espressione prescrive pertanto l'esecuzione del comando di archiviazione ogni lunedì alle ore 5:00 del mattino.

Oltre ai singoli valori puntuali e alle wildcard, è possibile specificare intervalli a passi discreti (\emph{step values}) tramite l'operatore barra (\texttt{/}). Applicato a un intervallo esplicito, come ad esempio \texttt{0-23/2} nel campo delle ore, l'operatore salta di due ore consecutive (equivalente all'enumerazione \texttt{0,2,4,6,8,10,12,14,16,18,20,22}). Applicato direttamente all'asterisco nel campo dei minuti (sintassi \texttt{*/2}), denota una ripetizione periodica a cadenza di ogni due minuti.

\subsection{Manifest Di Configurazione Ed Esecuzione Del CronJob}

Il manifest che segue definisce un CronJob denominato \texttt{pod-cronjob}, programmato per eseguire ogni due minuti un'immagine minimale BusyBox contenente un semplice script di echo:

\begin{lstlisting}[language=yaml]
apiVersion: batch/v1
kind: CronJob
metadata:
  name: pod-cronjob
spec:
  schedule: "*/2 * * * *"
  jobTemplate:
    spec:
      template:
        spec:
          containers:
            - name: pod-cronjob
              image: busybox
              args:
                - /bin/sh
                - -c
                - date; echo hello from k8s cluster
          restartPolicy: OnFailure
\end{lstlisting}

Si noti che per i carichi di tipo Job, il parametro \texttt{restartPolicy} non può assumere il valore ordinario \texttt{Always}, in quanto il container non deve riavviarsi una volta terminata con successo l'elaborazione; le opzioni ammesse sono unicamente \texttt{OnFailure} o \texttt{Never}.

\subsection{Comportamento Del Runtime E Ispezione Dello Stato}

Allo scoccare dell'intervallo temporale programmato, il controller CronJob istanzia un oggetto di tipo Job, il quale a sua volta crea il Pod esecutivo conformemente al \texttt{jobTemplate}.

Una volta che il processo principale contenuto nel container esegue l'uscita con codice di terminazione zero (\texttt{exit code 0}), il Pod transita irreversibilmente nello stato **Completed** (fase \texttt{Succeeded}). Poiché il container non è più attivo e i relativi namespace di esecuzione sono stati chiusi, non è più possibile aprire sessioni interattive su di esso (il comando \texttt{kubectl exec} genererà un errore). Tuttavia, i registri di stdout e stderr vengono preservati all'interno del nodo per la consultazione diagnostica tramite \texttt{kubectl logs <pod-name>}.

L'interrogazione dello stato mediante la CLI evidenzia la generazione sequenziale degli artefatti nel tempo:
\begin{lstlisting}[language=bash]
$ kubectl get all
NAME                             SCHEDULE      TIMEZONE   SUSPEND   ACTIVE   LAST SCHEDULE   AGE
cronjob.batch/pod-cronjob        */2 * * * *   <none>     False     0        <none>          23s

# A seguito delle attivazioni periodiche:
NAME                                READY   STATUS      RESTARTS   AGE
pod/pod-cronjob-29006744-l98fg      0/1     Completed   0          2m33s
pod/pod-cronjob-29006746-vw2gg      0/1     Completed   0          33s

NAME                                STATUS      COMPLETIONS   DURATION   AGE
job.batch/pod-cronjob-29006744      Complete    1/1           5s         2m33s
job.batch/pod-cronjob-29006746      Complete    1/1           5s         33s
\end{lstlisting}

I metadati generati associano a ciascun Job derivato un suffisso numerico progressivo derivante dal timestamp UNIX dell'istante di schedulazione, garantendo la tracciabilità univoca delle diverse esecuzioni a lotti nel cluster.

\newpage

% ==============================================================================
% SEZIONE 4: IL CONTROLLER DEPLOYMENT E LE STRATEGIE DI RILASCIO
% ==============================================================================

\section{Il Controller Deployment E Le Strategie Di Rilascio}

Nei data center moderni e nelle architetture cloud native, il deployment manuale o la gestione diretta di Pod e ReplicaSet elementari si rivela insufficiente per governare requisiti complessi quali la scalabilità dinamica, il rilascio continuo senza disservizio (\emph{zero-downtime}) e il rollback immediato a versioni precedenti in caso di regressione. Per assolvere a tali esigenze, Kubernetes espone il controller di livello più elevato: il **Deployment**.

Un Deployment non gestisce direttamente i singoli Pod; esso opera come supervisore dichiarativo al di sopra dei ReplicaSet, definendo lo stato desiderato dell'applicazione. Quando viene creato o modificato un Deployment, il control plane istanzia o adatta un ReplicaSet sottostante, il quale ha l'onere operativo di creare e monitorare i Pod fisici.

% [IMMAGINE DEL FILE DI SLIDE CON NOME "Orchestrazione_1"] Diagramma gerarchico a strati che illustra la relazione di composizione: Deployment (livello superiore dichiarativo) supervisiona il ReplicaSet, che a sua volta supervisiona l'insieme dei Pod replicati — Slide 62 / Pagina 12 del PDF

\subsection{Proprietà Fondamentali Della Specifica Di Un Deployment}

La sezione \texttt{spec} del Deployment disciplina quattro blocchi di configurazione cardine:
\begin{itemize}[noitemsep]
    \item \texttt{replicas}: indica la cardinalità dell'insieme di Pod che devono risultare costantemente in esecuzione e pronti a servire traffico;
    \item \texttt{strategy}: governa l'algoritmo di rimpiazzo e transizione adottato durante l'aggiornamento dell'applicazione;
    \item \texttt{selector}: specifica mediante predicati di etichetta (\texttt{matchLabels}) come identificare i Pod assegnati al governo di quel Deployment;
    \item \texttt{template}: contiene la specifica vera e propria del Pod (\emph{PodTemplateSpec}), comprensiva di immagini dei container, porte, variabili d'ambiente e montaggi di volume.
\end{itemize}

\begin{lstlisting}[language=yaml]
apiVersion: apps/v1
kind: Deployment
metadata:
  name: nginx-deployment
  labels:
    app: nginx
spec:
  replicas: 3
  selector:
    matchLabels:
      app: nginx
  template:
    metadata:
      labels:
        app: nginx
    spec:
      containers:
        - name: nginx
          image: nginx:1.14.2
          ports:
            - containerPort: 80
\end{lstlisting}

\subsection{Topologia Dei Nomi Dei Pod E Hashing Del Template}

Esaminando l'elenco dei Pod generati da un Deployment attraverso il comando \texttt{kubectl get pods}, si riscontra una struttura di denominazione rigorosa e composita:
\begin{lstlisting}[language=bash]
$ kubectl get pods
NAME                                READY   STATUS    RESTARTS   AGE
nginx-deployment-647677fc66-6p45z   1/1     Running   0          71s
nginx-deployment-647677fc66-cv49d   1/1     Running   0          71s
nginx-deployment-647677fc66-txwb8   1/1     Running   0          71s
\end{lstlisting}

La nomenclatura riflette il formalismo gerarchico:
\[
\texttt{<deployment-name>-<pod-template-hash>-<random-suffix>}
\]

\begin{enumerate}[noitemsep]
    \item **Deployment Name**: radice identificativa della risorsa controllore genitrice (ad esempio \texttt{nginx-deployment});
    \item **Template Hash**: stringa alfanumerica (ad esempio \texttt{647677fc66}) derivata dall'applicazione di una funzione di hash crittografico deterministica applicata all'intero corpo della specifica del Pod Template (\texttt{spec.template}). Tale hash funge da chiave identificativa del ReplicaSet di pertinenza: qualora il template venga alterato dall'utente (modificando ad esempio la versione dell'immagine del container), il valore di hash varia, costringendo il Deployment a generare un \emph{nuovo} ReplicaSet parallelo e avviando la sequenza di rilascio;
    \item **Random Suffix**: sequenza casuale generata dallo scheduler (ad esempio \texttt{6p45z}) per garantire la collision avoidance degli identificatori univoci all'interno dell'etcd cluster.
\end{enumerate}

\subsection{Strategie Di Rilascio: Recreate Vs RollingUpdate}

Nell'aggiornamento di una flotta applicativa da una versione iniziale $v_1$ a una nuova versione $v_2$, si possono adottare due strategie mutuamente esclusive:

\begin{itemize}
    \item **Recreate**: Kubernetes forza la terminazione e l'eliminazione istantanea di tutti i Pod associati alla versione $v_1$. Soltanto quando la flotta precedente è stata completamente smantellata, il controller avvia la creazione dei Pod attestati sulla versione $v_2$. Questa strategia comporta una disconnessione temporanea e inevitabile del servizio (\emph{downtime}), ma risulta necessaria qualora il software non supporti l'accesso concorrente di versioni eterogenee al medesimo schema di database relazionale o backend di memorizzazione.
    \item **RollingUpdate**: rappresenta la modalità predefinita e preferibile. I Pod della versione $v_1$ vengono rimpiazzati con istanze della versione $v_2$ in modo incrementale e graduale, garantendo che una quota minima di Pod rimanga costantemente attiva e pronta a processare il traffico utente, azzerando i disservizi.
\end{itemize}

\subsection{Formalizzazione Di \texttt{maxSurge} E \texttt{maxUnavailable}}

La dinamica del **RollingUpdate** è governata analiticamente da due parametri chiave configurabili dall'amministratore (espressi come numeri interi assoluti o frazioni percentuali rispetto al campo \texttt{replicas}):

\begin{itemize}
    \item \textbf{\texttt{maxUnavailable}}: esprime il numero massimo di Pod che possono risultare non operativi (non pronti o in fase di deallocazione) durante l'intera procedura di aggiornamento. Il valore predefinito è pari al 25\%. Se impostato a 0, il sistema non tollera alcuna degradazione della capacità minima di erogazione rispetto al numero nominale di repliche (imponendo conseguentemente $\texttt{maxSurge} > 0$).
    \item \textbf{\texttt{maxSurge}}: esprime il numero massimo di Pod ammessi in eccesso rispetto al numero nominale di repliche richieste. Il valore predefinito è pari al 25\%. Se impostato a 0, la flotta non potrà mai superare la cardinalità di repliche desiderata (imponendo conseguentemente $\texttt{maxUnavailable} > 0$).
\end{itemize}

Sia $R$ la cardinalità nominale dei Pod desiderati configurata nel Deployment, $U$ il valore di \texttt{maxUnavailable} e $S$ il valore di \texttt{maxSurge}. In qualunque istante temporale $t$ durante la procedura di migrazione, il numero totale di Pod in esecuzione (somma di vecchie e nuove istanze) $N_{\text{tot}}(t)$ e il numero di Pod attivi e pronti $N_{\text{ready}}(t)$ sono rigorosamente vincolati dall'intervallo:
\[
\boxed{R - U \le N_{\text{ready}}(t) \le N_{\text{tot}}(t) \le R + S}
\]

Si consideri a titolo d'esempio il manifest:
\begin{lstlisting}[language=yaml]
apiVersion: apps/v1
kind: Deployment
metadata:
  name: rolling-nginx
spec:
  replicas: 4
  strategy:
    type: RollingUpdate
    rollingUpdate:
      maxSurge: 1
      maxUnavailable: 1
  selector:
    matchLabels:
      app: rolling-nginx
  template:
    metadata:
      labels:
        app: rolling-nginx
    spec:
      containers:
        - name: nginx
          image: nginx:1.26.3
\end{lstlisting}

Con $R=4$, $S=1$ e $U=1$:
\begin{itemize}[noitemsep]
    \item Il numero massimo di Pod totali allocabili sull'infrastruttura non eccederà mai $4 + 1 = 5$ Pod;
    \item Il numero minimo di Pod operativi non scenderà mai al di sotto di $4 - 1 = 3$ Pod.
\end{itemize}

% [IMMAGINE DEL FILE DI SLIDE CON NOME "Orchestrazione_1"] Diagramma temporale del ciclo di rollout (RollingUpdate): rappresentazione a blocchi dell'interleaving tra distruzione dei vecchi pod (terminating) e creazione incrementale dei nuovi pod (deploying/ready) con il mantenimento del vincolo di bilanciamento (3 disponibili - 5 massimi) — Slide 66 / Pagina 16 del PDF

\subsection{Evoluzione Dello Stato Nel Rollout: Analisi Dei Log Di Laboratorio}

Nel momento in cui l'immagine del container viene aggiornata da \texttt{nginx:1.26.3} a \texttt{nginx:latest}, il controller avvia la sequenza temporale di riconciliazione registrata negli screenshot di laboratorio:

\begin{enumerate}
    \item **Stato Stazionario Iniziale**: sono attivi quattro Pod attestati sul vecchio ReplicaSet con hash identificativo \texttt{675d74594d}:
\begin{lstlisting}[language=bash]
NAME                                     READY   STATUS    RESTARTS   AGE
pod/rolling-nginx-675d74594d-4jrhz       1/1     Running   0          28s
pod/rolling-nginx-675d74594d-csjgx       1/1     Running   0          27s
pod/rolling-nginx-675d74594d-k924w       1/1     Running   0          27s
pod/rolling-nginx-675d74594d-nrzv4       1/1     Running   0          27s
\end{lstlisting}

    \item **Innesco Del Rilascio**: conformemente a $\texttt{maxSurge}=1$ e $\texttt{maxUnavailable}=1$, il Deployment ordina al nuovo ReplicaSet (avente nuovo template hash \texttt{cd7586f65}) di creare due nuove istanze, portando contemporaneamente a terminazione una vecchia istanza per soddisfare i margini:
\begin{lstlisting}[language=bash]
pod/rolling-nginx-675d74594d-4jrhz       0/1     Completed            0   2m27s
pod/rolling-nginx-675d74594d-csjgx       1/1     Running              0   2m26s
pod/rolling-nginx-675d74594d-k924w       1/1     Running              0   2m26s
pod/rolling-nginx-675d74594d-nrzv4       1/1     Running              0   2m26s
pod/rolling-nginx-cd7586f65-pggqt        0/1     ContainerCreating    0   4s
pod/rolling-nginx-cd7586f65-x56sv        0/1     ContainerCreating    0   4s
\end{lstlisting}

    \item **Progressione Concorrente**: non appena la prima istanza della nuova versione completa l'inizializzazione e passa in stato \texttt{Running 1/1} (es. \texttt{x56sv}), una seconda istanza precedente (\texttt{nrzv4}) riceve il segnale di terminazione, mentre viene avviato un ulteriore nuovo container (\texttt{d44cp}):
\begin{lstlisting}[language=bash]
pod/rolling-nginx-675d74594d-csjgx       1/1     Running              0   2m36s
pod/rolling-nginx-675d74594d-k924w       1/1     Running              0   2m36s
pod/rolling-nginx-675d74594d-nrzv4       1/1     Terminating          0   2m36s
pod/rolling-nginx-cd7586f65-d44cp        0/1     ContainerCreating    0   3s
pod/rolling-nginx-cd7586f65-pggqt        0/1     ContainerCreating    0   14s
pod/rolling-nginx-cd7586f65-x56sv        1/1     Running              0   14s
\end{lstlisting}

    \item **Convergenza Finale**: a rotazione, tutti i Pod della versione precedente vengono posti in \texttt{Terminating} e infine eliminati. Il Deployment converge al nuovo stato di equilibrio in cui l'intero set di 4 repliche appartiene esclusivamente al nuovo ReplicaSet:
\begin{lstlisting}[language=bash]
NAME                                     READY   STATUS    RESTARTS   AGE
pod/rolling-nginx-cd7586f65-d44cp        1/1     Running   0          16s
pod/rolling-nginx-cd7586f65-f5wfs        1/1     Running   0          12s
pod/rolling-nginx-cd7586f65-pggqt        1/1     Running   0          27s
pod/rolling-nginx-cd7586f65-x56sv        1/1     Running   0          27s
\end{lstlisting}

L'intero ciclo si completa senza che l'applicazione abbia mai smesso di erogare il servizio su almeno tre istanze pienamente operative.
\end{enumerate}

\newpage
% ==============================================================================
% SEZIONE 5: MECCANISMI DI SELEZIONE: LABEL E SELECTOR
% ==============================================================================

\section{Meccanismi Di Selezione: Label E Selector}

In Kubernetes, il modello concettuale si basa su un disaccoppiamento radicale tra le risorse controllori e gli oggetti controllati. Un Deployment o un Service non mantengono un puntatore statico o un array di puntatori verso i Pod generati. L'associazione logica viene invece realizzata dinamicamente mediante coppie chiave-valore denominate **Labels** (Etichette) e meccanismi di interrogazione detti **Label Selectors**.

Le etichette presentano proprietà strutturali specifiche:
\begin{itemize}[noitemsep]
    \item Non offrono unicità: a differenza dei nomi o degli UID dei Pod, una stessa etichetta è deliberatamente assegnata a molteplici oggetti che condividono la stessa funzione o il medesimo ambiente (es. \texttt{environment=production}, \texttt{tier=frontend});
    \item Sono flessibili e dinamiche: possono essere definite all'atto della creazione dell'oggetto oppure aggiunte, rimosse o sovrascritte a runtime tramite API senza arrestare il carico operativo.
\end{itemize}

I selettori rappresentano la primitiva di raggruppamento per eccellenza del cluster, consentendo l'isolamento di sottoinsiemi specifici di oggetti per scopi di monitoraggio, instradamento del traffico o governo dei controller.

\subsection{Tipologie Di Selettori: Equality-Based E Set-Based}

Le specifiche dell'API supportano due formalismi di selettori:

\subsubsection{Selettori Basati Sull'Uguaglianza (Equality-Based)}
Operano mediante test di equivalenza diretta su singole stringhe (\emph{IS / IS NOT}):
\begin{itemize}[noitemsep]
    \item \texttt{key = value} (o alternativamente \texttt{key == value}): seleziona tutti gli oggetti che contengono la chiave \texttt{key} associata esattamente al valore \texttt{value} specificato (es. \texttt{tier = frontend});
    \item \texttt{key != value}: seleziona gli oggetti che non possiedono tale chiave oppure che la possiedono associata a un valore disuguale (es. \texttt{tier != frontend});
    \item **Congiunzione logica (AND)**: è possibile concatenare requisiti multipli separandoli da virgola. Ad esempio, l'espressione:
    \[
    \texttt{tier != frontend, game = super-shooter-2}
    \]
    restituisce tutti i Pod associati al gioco \texttt{super-shooter-2} escludendo contestualmente quelli appartenenti al livello frontend.
\end{itemize}

\subsubsection{Selettori Basati Sugli Insiemi (Set-Based)}
Consentono il filtraggio su domini finiti di valori mediante la logica insiemistica (\emph{IN / NOT IN / EXISTS}):
\begin{itemize}[noitemsep]
    \item **Operatore \texttt{In}**: verifica l'appartenenza del valore a una lista enumerata;
    \item **Operatore \texttt{NotIn}**: esclude gli oggetti i cui valori appartengono all'insieme fornito;
    \item **Operatore \texttt{Exists}**: verifica la mera presenza della chiave nei metadati, indipendentemente dal valore ad essa assegnato.
\end{itemize}

Il seguente frammento YAML, tipico di una specifica di ReplicaSet o Deployment moderno, evidenzia l'interazione tra selettori equality-based (\texttt{matchLabels}) e set-based (\texttt{matchExpressions}):

\begin{lstlisting}[language=yaml]
selector:
  matchLabels:
    app: myapp # Selettore basato sull'uguaglianza (Equality-based)
  matchExpressions: # Selettori basati sugli insiemi (Set-based)
    - key: environment
      operator: In
      values: ["production", "qa"]
    - key: tier
      operator: NotIn
      values: ["frontend", "backend"]
    - key: partition
      operator: Exists # Condizione soddisfatta per qualsiasi valore della chiave
\end{lstlisting}

In questo contesto, il controller riconoscerà e aggregherà sotto la propria gestione unicamente i Pod che soddisfano l'intersezione (AND booleano) di tutti i predicati espressi.

\newpage

% ==============================================================================
% SEZIONE 6: GESTIONE, LIMITAZIONE E ISOLAMENTO DELLE RISORSE DI CALCOLO
% ==============================================================================

\section{Gestione, Limitazione E Isolamento Delle Risorse Di Calcolo}

In un'architettura multi-tenant o in cluster distribuiti condivisi, la convivenza di container all'interno del medesimo nodo fisico impone rigidi meccanismi di segregazione e garanzia delle risorse computazionali. Se i processi potessero consumare liberamente le risorse del sistema operativo host, un singolo container con perdite di memoria (\emph{memory leak}) o cicli computazionali non controllati potrebbe saturare l'host, innescando l'effetto noto come \emph{noisy neighbor} e compromettendo i servizi critici adiacenti.

Kubernetes supporta la limitazione su tre classi fondamentali di risorse computazionali:
\begin{itemize}[noitemsep]
    \item **CPU**: capacità computazionale di calcolo;
    \item **Memory**: memoria ad accesso casuale primaria (RAM);
    \item **Hugepages**: supporto per pagine di memoria a dimensionalità estesa (introdotto stabilmente a partire da Kubernetes v1.14).
\end{itemize}

\subsection{Dicotomia Tra Requests E Limits}

Per ciascun container definito all'interno di un Pod, è possibile dichiarare due soglie per ogni risorsa:

\begin{itemize}
    \item \textbf{Requests (\emph{Soft Limit})}: quantifica la quantità minima di risorse richiesta dal container per operare. Il valore di \texttt{requests} rappresenta l'unità atomica considerata dallo **Scheduler** di Kubernetes per determinare la fattibilità del piazzamento su un determinato nodo worker. Uno scheduler posiziona un Pod su un nodo solo se la capacità residua non allocata del nodo è strettamente maggiore o uguale alla somma delle \texttt{requests} di tutti i container del Pod:
    \[
    \sum_{\text{Pod } i} \text{requests}_{\text{cpu}, i} \le \text{Capacit\`a Allocabile}_{\text{cpu, nodo}}
    \]
    Una volta allocato, se il nodo dispone di capacità computazionale libera, il container è autorizzato ad espandere il proprio consumo oltre il valore espresso nella request.
    \item \textbf{Limits (\emph{Hard Limit})}: fissa il tetto massimo invalicabile di consumo ammesso per il container sul nodo in cui è in esecuzione.
\end{itemize}

La specifica adotta le seguenti chiavi all'interno della struttura \texttt{spec.containers[].resources}:
\begin{lstlisting}[language=yaml]
resources:
  requests:
    cpu: <valore>
    memory: <valore>
    hugepages-<size>: <valore>
  limits:
    cpu: <valore>
    memory: <valore>
    hugepages-<size>: <valore>
\end{lstlisting}

\paragraph{Regola Di Ereditarietà Implicita:} Qualora un amministratore definisca esplicitamente il parametro \texttt{limits} per una risorsa (CPU o memoria) omettendo tuttavia il corrispondente blocco \texttt{requests}, il control plane di Kubernetes assegna automaticamente alla request un valore perfettamente identico al limite:
\[
\text{requests}_{\text{assegnata}} \equiv \text{limits}_{\text{dichiarata}}
\]
Al contrario, se non viene specificato alcun limite di CPU, il container non presenta alcun tetto superiore di consumo e potrà potenzialmente assorbire l'intera potenza di calcolo del nodo, a meno che non intervenga un oggetto di tipo \textbf{LimitRange} impostato a livello di namespace per iniettare valori di default vincolanti.

\subsection{Modellizzazione E Unità Di Misura Della CPU}

La risorsa CPU in Kubernetes viene conteggiata in unità scalari di \emph{core}:
\begin{itemize}[noitemsep]
    \item Un'unità intera di CPU (\texttt{1}) equivale esattamente a \textbf{1 vCPU} (o Virtual Core) presso i principali cloud provider (quali AWS, GCP, Azure), oppure a \textbf{1 hyperthread} hardware su architetture server bare-metal Intel/AMD dotate di multithreading simultaneo.
    \item Poiché spesso un microservizio necessita di frazioni modeste di core, è consentito specificare il valore in forma decimale (es. \texttt{0.5}, che rappresenta mezzo core) oppure adottando l'unità dei **millicores** (espressa con il suffisso dimensionale \texttt{m}):
    \[
    1\,\text{CPU} = 1000\,\text{millicores (1000m)} \implies 0.5\,\text{CPU} \equiv 500\text{m}
    \]
    \item La granularità minima assoluta ammissibile dall'API è pari a \textbf{1m} (ovvero $0.001$ CPU core).
\end{itemize}

\subsection{Modellizzazione E Unità Di Misura Della Memoria}

La memoria viene dimensionata computazionalmente in byte ed è esprimibile attraverso due convenzioni sintattiche:
\begin{itemize}
    \item **Prefissi Decimali (Sistema Internazionale - SI)**: basati su potenze di dieci ($10^{3k}$):
    \begin{center}
    \begin{tabular}{llr}
        \toprule
        \textbf{Nome} & \textbf{Suffisso} & \textbf{Valore Equivalente (Byte)} \\
        \midrule
        Kilobyte & \texttt{K} & $1000^1 = 10^3 = \num{1000}$ \\
        Megabyte & \texttt{M} & $1000^2 = 10^6 = \num{1000000}$ \\
        Gigabyte & \texttt{G} & $1000^3 = 10^9 = \num{1000000000}$ \\
        Terabyte & \texttt{T} & $1000^4 = 10^{12} = \num{1000000000000}$ \\
        Petabyte & \texttt{P} & $1000^5 = 10^{15} = \num{1000000000000000}$ \\
        Exabyte  & \texttt{E} & $1000^6 = 10^{18} = \num{1000000000000000000}$ \\
        \bottomrule
    \end{tabular}
    \end{center}

    \item **Prefissi Binari (IEEE/IEC)**: basati su potenze di due ($2^{10k}$), comunemente utilizzati per prevenire ambiguità nei sistemi operativi:
    \begin{center}
    \begin{tabular}{llr}
        \toprule
        \textbf{Nome} & \textbf{Suffisso} & \textbf{Valore Equivalente (Byte)} \\
        \midrule
        Kibibyte & \texttt{Ki} (o KiB) & $2^{10} = \num{1024}$ \\
        Mebibyte & \texttt{Mi} (o MiB) & $2^{20} = \num{1048576}$ \\
        Gibibyte & \texttt{Gi} (o GiB) & $2^{30} = \num{1073741824}$ \\
        Tebibyte & \texttt{Ti} (o TiB) & $2^{40} = \num{1099511627776}$ \\
        Pebibyte & \texttt{Pi} (o PiB) & $2^{50} = \num{1125899906842624}$ \\
        Exbibyte & \texttt{Ei} (o EiB) & $2^{60} = \num{1152921504606846976}$ \\
        \bottomrule
    \end{tabular}
    \end{center}
\end{itemize}

\subsection{Meccanismi Kernel Di Isolamento Ed Enforcement: cgroups, CFS E OOM-Killer}

All'atto pratico, l'agente locale **kubelet** non impone i limiti direttamente sui processi applicativi, bensì trasmette le direttive computazionali al motore di containerizzazione sottostante (**containerd** o **Docker**), il quale a sua volta configura i costrutti di isolamento nativi del kernel Linux: i **cgroups** (\emph{Control Groups}, storicamente differenziati nelle architetture \emph{cgroup v1} e nell'unificata \emph{cgroup v2}).

L'imposizione dei limiti per le due risorse segue logiche fisiche diametralmente opposte, dipendenti dalla natura \emph{comprimibile} (\emph{compressible}) della CPU e \emph{non comprimibile} (\emph{incompressible}) della memoria:

\subsubsection{Enforcement Della CPU E Throttling}
La CPU è una risorsa comprimibile. L'allocazione della potenza di calcolo tra diversi cgroups è governata dal componente dello scheduler del kernel denominato **CFS** (\emph{Completely Fair Scheduler}).
Il runtime fissa convenzionalmente una finestra periodica temporale fondamentale di misura pari a $\tau = \SI{100}{\milli\second}$. La quota di calcolo $Q$ concessa al container in ogni finestra è determinata scalando la frazione di core configurata:
\[
\boxed{Q = (\text{CPU limit in core}) \times \SI{100}{\milli\second}}
\]
Ad esempio, se viene imposto un limite di $\texttt{cpu: "0.5"}$, il container ottiene una quota di esecuzione su CPU pari a:
\[
Q = 0.5 \times \SI{100}{\milli\second} = \SI{50}{\milli\second}
\]
Qualora i thread attivi all'interno del cgroup saturino il tempo di CPU a loro spettante ($\SI{50}{\milli\second}$) prima che sia trascorso l'intero periodo di $\SI{100}{\milli\second}$, il kernel interviene sospendendo (\emph{throttling}) l'esecuzione dei processi. I processi non vengono terminati né danneggiati: essi rimangono in stato di attesa forzata fino all'apertura della finestra successiva di $\SI{100}{\milli\second}$.
Inoltre, qualora i processi di un cgroup rimangano inattivi (\emph{idle}), il CFS ridistribuisce dinamicamente le porzioni temporali inutilizzate ad altri cgroups che presentano code di esecuzione in corso.

\subsubsection{Enforcement Della Memoria E OOM-Killer}
La memoria RAM è una risorsa non comprimibile: una pagina di memoria fisica allocata a un buffer non può essere sottratta temporaneamente per essere restituita a terzi senza alterare lo spazio di indirizzamento virtuale del processo.

Quando la quantità aggregata di memoria consumata dal container raggiunge la soglia rigida definita in \texttt{limits.memory}:
\begin{enumerate}[noitemsep]
    \item Il kernel Linux tenta in prima istanza di recuperare spazio svuotando la memoria cache e i buffer liberabili associati al cgroup;
    \item Se la richiesta di allocazione persiste e il cgroup sfora la quota, il gestore di memoria del kernel invoca l'**OOM-Killer** (\emph{Out-Of-Memory Killer});
    \item L'OOM-Killer seleziona il processo con il più alto punteggio di rischio all'interno del cgroup e ne forza l'arresto immediato inviando un segnale \texttt{SIGKILL} (\emph{Exit Code 137});
    \item In Kubernetes, il Pod transita nello stato di errore \texttt{OOMKilled} e viene eventualmente riavviato dal kubelet se consentito dalla \texttt{restartPolicy}.
\end{enumerate}

Qualora un container non oltrepassi il proprio limite assoluto ma ecceda la quantità dichiarata in \texttt{requests}, e il nodo host nel suo complesso inizi a manifestare condizioni critiche di esaurimento memoria (\emph{node memory pressure}), il kubelet identifica tale container come candidato primario per l'**espulsione** (\emph{eviction}) dal nodo, al fine di salvaguardare la stabilità complessiva del sistema operativo ospitante.

\subsection{Pagine Di Memoria Estese: HugePages}

Nei sistemi calcolatori tradizionali basati sull'architettura x86, la memoria virtuale è strutturata in pagine elementari gestite dalla \emph{Memory Management Unit} (MMU) con dimensione standard fissata a \textbf{4 KB}. La traduzione tra indirizzi logici generati dal software e indirizzi fisici su banchi RAM è coadiuvata dal registro **TLB** (\emph{Translation Lookaside Buffer}), una cache hardware associata alla CPU che memorizza le traduzioni virtuale-fisico più recenti.

Nelle applicazioni intensive (quali database di larga scala, algoritmi di calcolo ad alte prestazioni o workload di tipo High-Performance Computing / HPC), la necessità di mappare svariate decine di gigabyte di memoria comporta la generazione di tabelle delle pagine (\emph{page tables}) imponenti, saturando rapidamente la capacità limitata del TLB e causando continui fallimenti di lookup (\emph{TLB misses}), con pesanti penalizzazioni sulla latenza di accesso alla memoria.

\subsubsection{Vantaggi Operativi Delle HugePages}
Le **HugePages** consistono nell'allocazione di pagine di memoria a dimensionalità maggiorata, tipicamente pari a \textbf{2 MiB} (o \textbf{1 GiB} su architetture server):
\begin{itemize}[noitemsep]
    \item **Abbattimento dei TLB Misses**: una singola entry nel TLB riferita a una pagina da 2 MiB copre lo stesso intervallo di memoria di ben 512 pagine convenzionali da 4 KB, aumentando drasticamente l'\emph{hit ratio} del TLB;
    \item **Riduzione della Frammentazione della Memoria**: riduce l'overhead del sistema operativo dovuto alla frammentazione di blocchi contigui sparsi.
\end{itemize}

A differenza della memoria convenzionale o della CPU, le HugePages **non supportano l'overcommit**: non è possibile richiedere più HugePages di quante siano fisicamente riservate a livello di kernel all'avvio. Inoltre, la richiesta di risorse HugePages impone contestualmente la presenza di un'esplicita dichiarazione di limiti di CPU o di memoria ordinaria.

\subsubsection{Manifest Di Un Pod Con HugePages}
Qualora un Pod consumi HugePages di dimensioni eterogenee, deve utilizzare la notazione formale del volume \texttt{medium: HugePages-<size>}:

\begin{lstlisting}[language=yaml]
apiVersion: v1
kind: Pod
metadata:
  name: huge-pages-example
spec:
  containers:
    - name: example
      image: fedora:latest
      command:
        - sleep
        - inf
      volumeMounts:
        - mountPath: /hugepages-2Mi
          name: hugepage-2mi
        - mountPath: /hugepages-1Gi
          name: hugepage-1gi
      resources:
        limits:
          hugepages-2Mi: 100Mi
          hugepages-1Gi: 2Gi
          memory: 100Mi
        requests:
          memory: 100Mi
  volumes:
    - name: hugepage-2mi
      emptyDir:
        medium: HugePages-2Mi
    - name: hugepage-1gi
      emptyDir:
        medium: HugePages-1Gi
\end{lstlisting}

\subsubsection{Configurazione E Ispezione Sul Nodo Host}
Sul nodo Linux è possibile verificare la disponibilità delle pagine estese ispezionando il file di pseudofs `/proc/meminfo`:
\begin{lstlisting}[language=bash]
cat /proc/meminfo | grep -i huge
# Output tipico:
# HugePages_Total:     0
# HugePages_Free:      0
# HugePages_Rsvd:      0
# HugePages_Surp:      0
# Hugepagesize:     2048 kB
\end{lstlisting}

Per abilitare ad esempio l'allocazione immediata di 64 pagine da 2 MiB (pari a un pool contiguo riservato di $\SI{128}{\mega\byte}$):
\begin{lstlisting}[language=bash]
echo 64 > /proc/sys/vm/nr_hugepages
\end{lstlisting}

Affinché il cluster riconosca la risorsa schedulabile, il demone kubelet deve operare con la feature-gate attiva:
\begin{lstlisting}[language=bash]
--feature-gates=HugePages=true
\end{lstlisting}
La conferma della registrazione della risorsa da parte del control plane si ottiene interrogando il nodo:
\begin{lstlisting}[language=bash]
kubectl describe node <node-name>
\end{lstlisting}
il quale esporrà tra le capacità allocabili i campi \texttt{hugepages-2Mi} e \texttt{hugepages-1Gi}.

\subsection{Esempio Completo Di Pod Multi-Container Con Limiti}

Il seguente manifest illustra un Pod multi-container (composto da un database MySQL e un frontend WordPress) configurato all'interno del namespace \texttt{cpu-limit}, in cui ogni entità dichiara soglie simmetriche di \texttt{requests} (250m CPU, 64Mi memoria) e \texttt{limits} (500m CPU, 128Mi memoria):

\begin{lstlisting}[language=yaml]
apiVersion: v1
kind: Pod
metadata:
  name: frontend
  namespace: cpu-limit
spec:
  containers:
    - name: db
      image: mysql
      env:
        - name: MYSQL_ROOT_PASSWORD
          value: "password"
      resources:
        requests:
          memory: "64Mi"
          cpu: "250m"
        limits:
          memory: "128Mi"
          cpu: "500m"
    - name: wp
      image: wordpress
      resources:
        requests:
          memory: "64Mi"
          cpu: "250m"
        limits:
          memory: "128Mi"
          cpu: "500m"
\end{lstlisting}

La creazione della risorsa tramite CLI imperativa produce il record di sistema:
\begin{lstlisting}[language=bash]
kubectl create -f pod-resource-limit.yml
pod/frontend created
\end{lstlisting}

\subsection{Governance Del Cluster: Le Resource Quotas}

Mentre i limiti a livello di singolo container salvaguardano i singoli nodi worker dal degrado locale, le politiche di multi-tenancy e condivisione sicura di un cluster richiedono strumenti di governo a grana più grossa. A tale scopo, Kubernetes implementa le **Resource Quotas**, incapsulate nell'oggetto API **ResourceQuota**.

Una ResourceQuota stabilisce un tetto massimo invalicabile sull'aggregato complessivo di risorse computazionali consumabili all'interno di un determinato \textbf{Namespace}. Essa consente a team infrastrutturali differenti di operare in autonomia nel medesimo cluster senza il rischio che un gruppo monopolizzi l'intera capacità computazionale.

\subsubsection{Ruolo Dell'Admission Controller E Risposte Di Violazione}
La conformità alle quote viene verificata in tempo reale e in modo sincrono prima della persistenza su etcd da un apposito **Admission Controller** integrato nell'API Server:
\begin{enumerate}[noitemsep]
    \item Ogni volta che un utente inoltra una richiesta di creazione o aggiornamento di un Pod (o di un controller che genera Pod), l'Admission Controller calcola la sommatoria tra le risorse già impegnate nel namespace e quelle richieste dal nuovo manifest.
    \item Se la nuova allocazione provoca lo sforamento di uno qualsiasi dei vincoli imposti dalla quota, la richiesta API viene **rifiutata atomicamente**.
    \item Il server HTTP interrompe la procedura restituendo al client il codice di stato:
    \[
    \boxed{\text{HTTP 403 FORBIDDEN}}
    \]
    corredato da un messaggio esplicativo contenente il dettaglio della risorsa violata, il valore corrente, il quantitativo richiesto e il limite massimo stabilito.
\end{enumerate}

\subsubsection{Tipologie Di Risorse Soggette A Quota}
Le ResourceQuota operano monitorando esclusivamente i Pod che si trovano in uno stato non terminale (escludendo pertanto i Pod in fase \texttt{Succeeded} o \texttt{Failed}). I principali identificatori di risorsa supportati sono riepilogati nella seguente tabella:

\begin{center}
\begin{tabularx}{\textwidth}{lX}
    \toprule
    \textbf{Nome Risorsa} & \textbf{Descrizione Operativa} \\
    \midrule
    \texttt{limits.cpu} & La somma dei limiti di CPU di tutti i Pod attivi nel namespace non può superare questo valore. \\
    \addlinespace
    \texttt{limits.memory} & La somma dei limiti di memoria di tutti i Pod attivi nel namespace non può superare questo valore. \\
    \addlinespace
    \texttt{requests.cpu} (o \texttt{cpu}) & La somma delle richieste di CPU di tutti i Pod attivi non può eccedere questo valore. \\
    \addlinespace
    \texttt{requests.memory} (o \texttt{memory}) & La somma delle richieste di memoria di tutti i Pod attivi non può eccedere questo valore. \\
    \addlinespace
    \texttt{hugepages-<size>} & Il numero complessivo di richieste di HugePages della specifica dimensione (es. 2Mi, 1Gi) non può superare la soglia configurata. \\
    \bottomrule
\end{tabularx}
\end{center}

L'introduzione di una ResourceQuota su un namespace impone indirettamente a tutti gli utenti che rilasciano applicazioni in quel contesto di dichiarare obbligatoriamente le \texttt{requests} e i \texttt{limits} per i propri container: in assenza di tali parametri nei manifest dei Pod, l'Admission Controller respingerebbe le istanze a priori per impossibilità di quantificare l'impatto sulla quota globale del namespace.

% AA Orchestrazione_1 [100-132] ------------------------------------------------------------

\section{Governance E Multi-Tenancy Delle Risorse: ResourceQuota}

Nei contesti aziendali e nelle infrastrutture cloud condivise tra team eterogenei, la gestione delle risorse computazionali non può essere demandata esclusivamente alla disciplina dei singoli sviluppatori. Sebbene la definizione puntuale di \texttt{requests} e \texttt{limits} a livello di singolo container consenta al kernel Linux di prevenire conflitti locali sui nodi worker, tale approccio non impedisce a un singolo gruppo di lavoro di saturare l'intera capacità computazionale del cluster instanziando un numero indefinito di repliche.

Per realizzare un reale isolamento logico multi-tenant, Kubernetes introduce l'oggetto API **ResourceQuota**, applicato alla primitiva di partizionamento per eccellenza: il \textbf{Namespace}.

\subsection{Definizione Operativa Di Una Quota Di CPU Su Namespace}

L'applicazione di una quota si articola su due passaggi operativi: la creazione del namespace dedicato e la successiva applicazione del vincolo aggregato. A titolo esemplificativo, si consideri la creazione del namespace denominato \texttt{quota-example} tramite CLI imperativa:

\begin{lstlisting}[language=bash]
[root@controller ~]# kubectl create namespace quota-example
namespace/quota-example created
\end{lstlisting}

Una volta istanziato il perimetro logico, viene redatto e applicato il manifest della risorsa \texttt{ResourceQuota}, all'interno del quale si fissa il tetto rigido aggregato per il consumo di calcolo:

\begin{lstlisting}[language=yaml]
apiVersion: v1
kind: ResourceQuota
metadata:
  name: resource-quota
  namespace: quota-example
spec:
  hard:
    limits.cpu: "1"
\end{lstlisting}

Nel manifest sopra riportato, la direttiva \texttt{spec.hard.limits.cpu: "1"} impone che la somma algebrica dei limiti di CPU (\texttt{limits.cpu}) dichiarati da tutti i container appartenenti a Pod in stato non terminale all'interno del namespace \texttt{quota-example} non possa in alcun caso superare $1$ core fisico (o vCPU).

Sebbene l'esempio illustri un vincolo focalizzato sulla CPU, la sezione \texttt{spec.hard} supporta la quantificazione combinata di molteplici metriche eterogenee, tra cui:
\begin{itemize}[noitemsep]
    \item **Memoria di sistema**: vincoli su \texttt{requests.memory} e \texttt{limits.memory} (es. limitando l'impegno aggregato a \texttt{4Gi});
    \item **Capacità computazionale minima**: vincoli su \texttt{requests.cpu};
    \item **Cardinalità degli oggetti**: limitazione numerica sul totale dei Pod istanziabili (campo \texttt{pods}), dei Service (campo \texttt{services}), o delle richieste di memorizzazione persistente (\texttt{persistentvolumeclaims});
    \item **Memoria estesa**: vincoli puntuali sulle pagine HugePages (\texttt{hugepages-<size>}).
\end{itemize}

Quando un utente tenta di distribuire un carico di lavoro che determina lo sforamento dei limiti stabiliti in \texttt{hard}, l'**Admission Controller** intercetta la richiesta HTTP a monte della persistenza su \textbf{etcd}, rigettandola con il codice d'errore standard \texttt{HTTP 403 Forbidden}.

\newpage

% ==============================================================================
% SEZIONE 2: DISACCOPPIAMENTO DELLE CONFIGURAZIONI E DATI SENSIBILI: CONFIGMAP E SECRET
% ==============================================================================

\section{Disaccoppiamento Delle Configurazioni E Dati Sensibili: ConfigMap E Secret}

Uno dei principi cardine della progettazione di applicazioni moderne nel cloud (formalizzato dalla metodologia \emph{Twelve-Factor App}) impone la rigida separazione tra il codice sorgente (immutabile e racchiuso nell'immagine del container) e la configurazione specifica dell'ambiente di runtime (dinamica e suscettibile a modifiche tra gli stadi di sviluppo, test e produzione).

Includere parametri applicativi, chiavi crittografiche o stringhe di connessione all'interno dell'immagine Docker comporterebbe gravi rischi di sicurezza e invaliderebbe il requisito di portabilità: per ogni variazione di un singolo flag o certificato, lo sviluppatore sarebbe costretto a ricompilare, ri-taggare e ridistribuire l'intero artefatto binario. Kubernetes risolve tale problematica introducendo due primitive native dedicate: i **ConfigMap** per le informazioni di configurazione non riservate e i **Secret** per i dati sensibili.

\subsection{ConfigMap: Iniezione Dati E Proprietà Del Montaggio}

I ConfigMap consentono di memorizzare dati di configurazione in formato chiave-valore e di renderli accessibili ai container all'avvio o durante l'esecuzione. I principali casi d'uso includono:
\begin{itemize}[noitemsep]
    \item **Variabili d'ambiente**: esportazione di parametri scalari direttamente nel runtime dei processi;
    \item **File di configurazione**: montaggio di interi blocchi di testo come file fisici all'interno del file system del container;
    \item **Argomenti da riga di comando**: passaggio dinamico di flag ai parametri \texttt{command} e \texttt{args} del manifest.
\end{itemize}

\subsubsection{Comportamento Del File System Montato}
Qualora un ConfigMap venga montato come volume all'interno di un percorso di directory del container, i file generati presentano la caratteristica fondamentale di essere **in sola lettura** (\emph{read-only by default}).

Se il processo applicativo tenta di sovrascrivere o scrivere direttamente sui file montati dalla ConfigMap, l'operazione di I/O su disco fallirà a livello di permessi POSIX (oppure, in specifiche configurazioni di overlay file system, la modifica rimarrà confinata nello strato \emph{read-write} effimero del container locale). In nessun caso le modifiche apportate dall'applicazione sul file system locale verranno propagate retroattivamente all'oggetto ConfigMap originale presente nel control plane di Kubernetes.

\subsection{Secret: Protezione Dei Dati Riservati}

I Secret condividono la medesima struttura topologica dei ConfigMap, ma sono espressamente progettati per ospitare credenziali confidenziali, quali password di database, chiavi di autenticazione per API esterne, certificati digitali X.509 e chiavi private TLS, nonché credenziali per l'accesso a registri Docker privati (\emph{ImagePullSecrets}).

\subsubsection{Tipologie Standard Di Secret}
A seconda della natura delle informazioni conservate, l'API Server prevede diverse categorizzazioni semantiche definite nel campo \texttt{type}:

\begin{center}
\begin{tabularx}{\textwidth}{llX}
    \toprule
    \textbf{Tipo di Secret} & \textbf{Formato Dati} & \textbf{Scopo Operativo} \\
    \midrule
    \texttt{Opaque} & Valori arbitrari Base64 & Dati non strutturati a schema libero (impostazione predefinita). \\
    \addlinespace
    \texttt{kubernetes.io/service-account-token} & Token, \texttt{ca.crt}, namespace & Autenticazione dei carichi di lavoro verso l'API Server. \\
    \addlinespace
    \texttt{kubernetes.io/dockercfg} & File \texttt{.dockercfg} serializzato & Credenziali di accesso a registri Docker privati (legacy). \\
    \addlinespace
    \texttt{kubernetes.io/dockerconfigjson} & JSON \texttt{.dockerconfigjson} & Credenziali moderne per il pull di immagini da registri privati. \\
    \addlinespace
    \texttt{kubernetes.io/basic-auth} & \texttt{username}, \texttt{password} & Autenticazione basilare HTTP/servizi. \\
    \addlinespace
    \texttt{kubernetes.io/ssh-auth} & Chiavi SSH pubblica/privata & Autenticazione mediante coppia di chiavi SSH. \\
    \addlinespace
    \texttt{kubernetes.io/tls} & \texttt{tls.crt}, \texttt{tls.key} & Certificati pubblici e chiavi private per terminazione TLS. \\
    \bottomrule
\end{tabularx}
\end{center}

Il tipo **Opaque** rappresenta la scelta più comune: il termine indica che Kubernetes tratta il payload come un blob "opaco" (non trasparente), senza effettuarne alcuna validazione sintattica interna o parsing di schema.

\subsection{La Natura Di Base64 E Le Misure Di Sicurezza Reali}

Un errore concettuale frequente risiede nel confondere l'offuscamento sintattico offerto dalla codifica **Base64** con un meccanismo di cifratura crittografica. Nei manifest YAML dei Secret, i campi del dizionario \texttt{data} devono essere tassativamente espressi in stringhe Base64 (RFC 4648).

Tuttavia, la codifica Base64 non offre alcuna riservatezza: essa costituisce una mera traslitterazione biunivoca progettata per trasportare dati binari arbitrari all'interno di protocolli testuali senza corruzioni di codifica. Chiunque ottenga visibilità sul manifest può decodificare il dato in chiaro in modo banale:

\begin{lstlisting}[language=bash]
# Ispezione del Secret da riga di comando
kubectl get secret my-secret -o yaml
# Output:
# data:
#   password: cGFzc3dvcmQxMjM=
#   username: YWRtaW4=

# Decodifica immediata del payload sensibile:
echo "cGFzc3dvcmQxMjM=" | base64 -d
# Risultato in chiaro: password123
\end{lstlisting}

Per garantire un reale livello di sicurezza enterprise, Kubernetes implementa due meccanismi ortogonali:

\subsubsection{Cifratura A Riposo (Encryption At Rest)}
Di default, i dati dei Secret vengono salvati in formato non cifrato all'interno del database di cluster **etcd**. Gli amministratori possono abilitare la cifratura a riposo configurando il componente \textbf{kube-apiserver} attraverso un file di configurazione di tipo \texttt{EncryptionConfiguration}. Ciò impone l'impiego di algoritmi crittografici simmetrici avanzati (quali AES-GCM o AES-CBC con chiavi a 256 bit) per cifrare il payload prima della scrittura fisica sui blocchi di memoria di etcd, impedendo l'accesso non autorizzato in caso di clonazione illecita dei dischi di storage.

\subsubsection{Controllo Degli Accessi Basato Su Ruoli (RBAC)}
L'accesso in lettura e modifica ai Secret viene rigorosamente filtrato mediante le policy RBAC di cluster. Un'entità (utente o ServiceAccount) non può visualizzare o montare Secret a meno che non sia esplicitamente associata a un oggetto \texttt{Role} o \texttt{ClusterRole} autorizzante:

\begin{lstlisting}[language=yaml]
apiVersion: rbac.authorization.k8s.io/v1
kind: Role
metadata:
  namespace: default
  name: secret-reader
rules:
  - apiGroups: [""]
    resources: ["secrets"]
    verbs: ["get", "list"]
\end{lstlisting}

\subsection{Manifest Di ConfigMap E Secret}

Si consideri l'esempio di una ConfigMap contenente parametri operativi e di un Secret contenente credenziali amministrative:

\begin{lstlisting}[language=yaml]
apiVersion: v1
kind: ConfigMap
metadata:
  name: my-config
data:
  # Definizione di un file di configurazione multi-riga
  app.properties: |
    color=blue
    size=large
  # Valore scalare singolo per il livello di logging
  log.level: info
---
apiVersion: v1
kind: Secret
metadata:
  name: my-secret
type: Opaque
data:
  username: YWRtaW4=            # Base64 di "admin"
  password: MWYyZDFlMmU2N2Rm    # Base64 di "1f2d1e2e67df"
\end{lstlisting}

Il parametro \texttt{log.level: info} istruisce l'applicazione ad avviare la registrazione delle informazioni a un determinato livello di verbosità, consentendo l'alterazione del comportamento del software modificando unicamente l'oggetto Kubernetes.

\subsection{Pattern Di Consumo: Environment Variables Vs Volume Mounts}

L'accesso ai dati definiti in ConfigMap e Secret all'interno dei Pod può avvenire secondo due modalità applicative principali:

\subsubsection{Iniezione Come Variabili D'Ambiente}
Le variabili possono essere importate in modo puntuale (selezionando una singola chiave tramite \texttt{valueFrom}) oppure in modo massivo (iniettando tutte le chiavi presenti mediante \texttt{envFrom}):

\begin{lstlisting}[language=yaml]
apiVersion: v1
kind: Pod
metadata:
  name: app-pod
  labels:
    app: demo
spec:
  containers:
    - name: app-container
      image: busybox
      command: ["/bin/sh", "-c", "sleep 3600"]
      # Metodo 1: Importazione puntuale di singole chiavi
      env:
        - name: LOG_LEVEL
          valueFrom:
            configMapKeyRef:
              name: my-config
              key: log.level
        - name: APP_USERNAME
          valueFrom:
            secretKeyRef:
              name: my-secret
              key: username
        - name: APP_PASSWORD
          valueFrom:
            secretKeyRef:
              name: my-secret
              key: password
      # Metodo 2: Iniezione massiva (ogni chiave diventa una variabile d'ambiente)
      envFrom:
        - configMapRef:
            name: my-config
        - secretRef:
            name: my-secret
\end{lstlisting}

\subsubsection{Montaggio Come File Di Volume}
In alternativa, le risorse possono essere collegate come volumi all'interno dell'albero del file system:

\begin{lstlisting}[language=yaml]
apiVersion: v1
kind: Pod
metadata:
  name: my-pod
spec:
  containers:
    - name: my-container
      image: my-image
      volumeMounts:
        - name: config-volume
          mountPath: /etc/config
        - name: secret-volume
          mountPath: /etc/secret
          readOnly: true
  volumes:
    - name: config-volume
      configMap:
        name: my-config
    - name: secret-volume
      secret:
        secretName: my-secret
\end{lstlisting}

\paragraph{Confronto Di Sicurezza:} Sebbene l'esposizione come variabili d'ambiente risulti computazionalmente comoda, essa è considerata **meno sicura** rispetto al montaggio su file system. Le variabili d'ambiente di un processo Linux sono facilmente ispezionabili da comandi di sistema (es. leggendo il file pseudofs \texttt{/proc/\$PID/environ}), possono trapelare accidentalmente nei report diagnostici in caso di \emph{crash dump} o eccezioni del runtime, e risultano visibili tramite semplici comandi di ispezione del motore dei container (\texttt{docker inspect} o \texttt{crictl inspect}). Di conseguenza, per le informazioni altamente riservate (chiavi private, certificati, password), la prassi architetturale raccomanda l'uso esclusivo di volumi montati con restrizioni stringenti sui permessi POSIX.

\newpage

% ==============================================================================
% SEZIONE 3: IL MODELLO DI RETE DI KUBERNETES: PRINCIPI FONDAMENTALI
% ==============================================================================

\section{Il Modello Di Rete Di Kubernetes: Principi Fondamentali}

L'infrastruttura di rete di Kubernetes è progettata per superare le rigidità e la complessità operativa dei sistemi di containerizzazione precoci, introducendo un modello dichiarativo uniforme che tratta i Pod come entità di prima classe equiparabili a macchine fisiche o virtuali indipendenti.

\subsection{I Quattro Problemi Fondamentali Del Networking Distribuito}

Nel governo di un cluster distribuito, l'architettura deve fornire una risposta solida a quattro classi distinte di connettività:
\begin{enumerate}[noitemsep]
    \item **Comunicazione Container-to-Container**: scambio dati a bassissima latenza tra processi cooperanti collocati all'interno della medesima entità logica atomica (risolto nativamente dal costrutto del **Pod** tramite interfaccia di \emph{loopback});
    \item **Comunicazione Pod-to-Pod**: connettività flat e diretta tra carichi applicativi distribuiti sul medesimo nodo o su nodi differenti dell'infrastruttura;
    \item **Comunicazione Pod-to-Service**: astrazione e bilanciamento del carico interno verso pool di Pod effimeri protetti da un'interfaccia stabile (gestito dalla risorsa **Service** e dal demone **kube-proxy**);
    \item **Comunicazione External-to-Service**: instradamento, esposizione pubblica e ingresso del traffico proveniente dall'esterno del data center verso i servizi interni (gestito tramite Service di tipo \texttt{NodePort}, \texttt{LoadBalancer} o controller \textbf{Ingress}).
\end{enumerate}

% [IMMAGINE DEL FILE DI SLIDE CON NOME "Orchestrazione_1"] Schema topologico della rete del cluster Kubernetes che connette su piani distinti i Nodi fisici, i Pod con indirizzi propri e i Servizi astratti — Slide 87 / Pagina 9 del PDF

\subsection{I Tre Assiomi Fondamentali Del Modello IP-Per-Pod}

Kubernetes rifiuta l'approccio classico di Docker basato sulla mappatura statica delle porte host-to-container (\emph{port forwarding}) o sull'impiego massivo del NAT per la comunicazione interna. Al fine di standardizzare l'ecosistema, il modello impone a qualsiasi plugin di rete (\textbf{CNI}) il rispetto tassativo di tre requisiti fondamentali:

\begin{enumerate}
    \item **I Pod su un nodo possono comunicare con tutti i Pod presenti su qualsiasi nodo senza ricorrere a Network Address Translation (NAT)**: ciascun Pod ottiene un indirizzo IP instradabile a livello di cluster. Quando il Pod $A$ con indirizzo \texttt{10.244.1.5} invia un pacchetto IP al Pod $B$ con indirizzo \texttt{10.244.2.6}, il pacchetto transita sulla rete preservando immutati gli indirizzi di sorgente e di destinazione nell'header IP. Il destinatario riceve il pacchetto e rileva come mittente esattamente l'IP originale \texttt{10.244.1.5}.
    \item **Gli agenti di sistema presenti su un nodo (ad esempio il demone \texttt{kubelet}) possono comunicare con tutti i Pod allocati su quel medesimo nodo senza NAT**: la macchina host dispone di visibilità e accesso diretto alla subnet dei Pod locali per eseguire compiti di monitoraggio, health check (\emph{liveness} e \emph{readiness probes}) e raccolta dei log.
    \item **I Pod eseguiti all'interno del namespace di rete dell'host (\texttt{hostNetwork: true}) possono comunicare con tutti i Pod su tutti i nodi senza NAT**: se un carico opera direttamente nello stack dell'host, il suo traffico verso i restanti Pod del cluster non subisce mascheramenti di indirizzo intermedi.
\end{enumerate}

\subsection{Ripartizione Degli Spazi Di Indirizzamento (Non-Overlapping IP)}

Affinché la rete operi in modo deterministico, il cluster richiede la configurazione di intervalli di indirizzi IP strettamente disgiunti, la cui allocazione è delegata a componenti distinti:
\begin{itemize}[noitemsep]
    \item **IP dei Nodi**: assegnati dall'infrastruttura di rete fisica, dal server DHCP del data center, dal provider cloud o dal componente **cloud-controller-manager**;
    \item **IP dei Pod**: assegnati dinamicamente dal plugin di rete (**CNI**) attingendo da un blocco CIDR globale riservato al cluster (\emph{Pod CIDR}, ad esempio \texttt{10.244.0.0/16});
    \item **IP dei Servizi (Virtual IP)**: generati e assegnati deterministicamente dal control plane tramite il componente **kube-apiserver** attingendo da un intervallo dedicato (\emph{Service CIDR}, ad esempio \texttt{10.96.0.0/12}).
\end{itemize}

\subsection{Analisi Del Parametro \texttt{hostNetwork: true}}

In scenari specifici in cui un'applicazione richiede l'accesso a basso livello alle interfacce fisiche dell'host (ad esempio agenti di telemetria delle schede di rete, demoni BGP o componenti core di networking), è possibile bypassare l'isolamento del namespace configurando la proprietà \texttt{hostNetwork: true}:

\begin{lstlisting}[language=yaml]
apiVersion: v1
kind: Pod
metadata:
  name: nginx
spec:
  hostNetwork: true
  containers:
    - name: nginx
      image: nginx
\end{lstlisting}

Quando \texttt{hostNetwork} è abilitato:
\begin{itemize}[noitemsep]
    \item Il container non riceve un nuovo indirizzo IP isolato, bensì condivide interamente lo stack di rete, le interfacce fisiche e l'indirizzo IP del nodo host;
    \item Le porte aperte dal processo (es. la porta TCP 80 di Nginx) vengono vincolate direttamente sulle interfacce di rete del nodo, inibendo la possibilità di avviare due istanze del medesimo Pod sul medesimo host a causa del conflitto di \emph{bind} sulla porta;
    \item La comunicazione verso gli altri Pod del cluster sfrutta l'instradamento diretto della scheda di rete dell'host: se il Pod $C$ (in esecuzione con \texttt{hostNetwork: true} sul nodo con IP \texttt{192.168.1.10}) contatta il Pod $B$ (avente IP \texttt{10.244.2.6}), il Pod $B$ osserva come indirizzo IP sorgente del frame esattamente \texttt{192.168.1.10}.
\end{itemize}

Questa modalità riduce l'overhead di incapsulamento, ma degrada l'isolamento di sicurezza ed espone il sistema host a vulnerabilità di tipo \emph{port exhaustion} o intercettazione del traffico locale.

\newpage

% ==============================================================================
% SEZIONE 4: NETWORKING INTRA-POD: DAL MODELLO DOCKER AL PAUSE CONTAINER
% ==============================================================================

\section{Networking Intra-Pod: Dal Modello Docker Al Pause Container}

Per comprendere a fondo la genesi del costrutto architetturale del **Pod**, è indispensabile esaminare l'implementazione del networking offerta dal runtime Docker convenzionale su un singolo host Linux, evidenziandone i limiti strutturali che hanno spinto gli ingegneri di Google a progettare l'astrazione del Pod basata sul **Pause Container**.

\subsection{Il Modello Docker Convenzionale Su Host Singolo}

In una tipica installazione Docker su una macchina Linux:
\begin{enumerate}
    \item Il demone crea all'interno del kernel dell'host un'interfaccia bridge di livello 2 denominata convenzionalmente \textbf{docker0}, assegnandole un indirizzo IP privato (ad esempio \texttt{172.17.0.1/24});
    \item All'avvio di un container, il runtime genera una coppia di interfacce virtuali collegate tra loro (\textbf{veth pair}, \emph{virtual ethernet}):
    \begin{itemize}[noitemsep]
        \item Un'estremità della coppia (es. \texttt{veth0}) viene iniettata all'interno del namespace di rete isolato del container, assumendo l'indirizzo IP \texttt{172.17.0.2};
        \item L'altra estremità (es. associata al dispositivo host \texttt{vethXXX}) rimane agganciata al bridge \texttt{docker0} in modalità promiscua.
    \end{itemize}
    \item I processi interni al container possiedono una vista circoscritta: rilevano unicamente la propria interfaccia \texttt{veth0}, non hanno visibilità del bridge \texttt{docker0} né dell'interfaccia fisica \texttt{eth0} dell'host, e utilizzano \texttt{172.17.0.1} come default gateway per l'instradamento;
    \item Quando il container trasmette traffico verso l'esterno (es. verso Internet), il kernel Linux dell'host applica regole di **Source NAT (SNAT)** o **Masquerading**, sostituendo l'indirizzo privato sorgente (\texttt{172.17.0.2}) con l'indirizzo IP pubblico o di rete locale dell'interfaccia fisica \texttt{eth0}. Al ritorno dei pacchetti di risposta, il kernel esegue l'operazione inversa (**Reverse NAT**), reindirizzando il frame al container appropriato.
\end{enumerate}

% [IMMAGINE DEL FILE DI SLIDE CON NOME "Orchestrazione_1"] Modello di connettività Docker locale su singolo host: scheda fisica eth0 collegata al bridge docker0, a cui si attesta la veth0 del container isolato — Slide 89 / Pagina 12 del PDF

\subsection{Il Limite Del Modello Docker Con Container Multipli}

Se su tale macchina si avvia un secondo container, Docker ripete la procedura: crea un nuovo namespace di rete, alloca una seconda coppia di interfacce \texttt{veth1} con IP \texttt{172.17.0.3} e la connette al bridge \texttt{docker0}.

Sebbene i due container possano comunicare tra loro attraverso il bridge mediante i rispettivi indirizzi IP privati, questo approccio non soddisfa il requisito architetturale di **stack di rete condiviso** richiesto da Kubernetes. Se due processi altamente accoppiati (ad esempio un web server primario e un processo secondario di aggiornamento della cache, o un demone di log scraping \emph{sidecar}) devono operare congiuntamente, trattarli come nodi di rete separati costringe ad allocare indirizzi IP distinti, preclude la comunicazione mediante memoria condivisa o socket Unix locali, e impedisce la cooperazione trasparente sull'interfaccia di loopback.

% [IMMAGINE DEL FILE DI SLIDE CON NOME "Orchestrazione_1"] Architettura con due container eseguiti in namespace di rete separati e connessi a docker0 mediante veth distinte — Slide 90 / Pagina 14 del PDF

\subsection{La Condivisione Del Network Namespace E Il Trucco Dei Symlink In Linux}

Il kernel Linux offre una notevole flessibilità nella gestione dei namespace. Un namespace di rete non è vincolato indissolubilmente a un singolo processo; esso è una risorsa di sistema rappresentata all'interno del file system sotto il percorso `/var/run/netns/`.

Quando si crea un namespace tramite comandi standard:
\begin{lstlisting}[language=bash]
sudo ip netns add ns1
\end{lstlisting}
il kernel registra il descrittore all'interno del file `/var/run/netns/ns1`.

Qualora si desideri che un secondo contesto condivida esattamente la medesima configurazione, le medesime tabelle di routing, lo stesso spazio delle porte e lo stesso IP, è possibile creare un collegamento simbolico:
\begin{lstlisting}[language=bash]
sudo ln -s /var/run/netns/ns1 /var/run/netns/ns2
\end{lstlisting}

Eseguendo interrogazioni di rete all'interno di \texttt{ns1} e del presunto \texttt{ns2}:
\begin{lstlisting}[language=bash]
sudo ip netns exec ns1 ip a
sudo ip netns exec ns2 ip a # Restituisce esattamente lo stesso identico output di ns1
\end{lstlisting}
si riscontra che i processi in esecuzione in \texttt{ns2} utilizzano in realtà la medesima infrastruttura di \texttt{ns1}. Docker supporta nativamente questo meccanismo permettendo l'avvio di un container con il flag:
\[
\texttt{--net=container:<container\_id>}
\]

% [IMMAGINE DEL FILE DI SLIDE CON NOME "Orchestrazione_1"] Modello a stack di rete condiviso tra container: container 1 e container 2 attestati sulla medesima veth0 e sul medesimo indirizzo IP — Slide 91 / Pagina 15 del PDF

In questo modello a stack condiviso:
\begin{itemize}[noitemsep]
    \item Tutti i container associati condividono lo stesso identico indirizzo IP (\texttt{172.17.0.2});
    \item La comunicazione reciproca avviene in modalità **Layer-2/MAC** diretta oppure tramite interfaccia di **loopback** invocando semplicemente \texttt{localhost:<porta>};
    \item Non interviene alcun processo di NAT all'interno del gruppo;
    \item Si applica il vincolo di collisione delle porte: due container nello stesso gruppo non possono effettuare il binding della medesima porta TCP/UDP (es. la porta 80).
\end{itemize}

\subsection{Architettura Del Pause Container Nel Pod}

La condivisione del namespace tra container solleva una criticità di gestione del ciclo di vita: se il container che "possiede" originariamente il namespace di rete dovesse arrestarsi, fallire o subire un riavvio, il kernel Linux deallocherebbe il namespace, provocando la disconnessione e la corruzione immediata di tutti gli altri container collegati.

Per superare questo problema, Kubernetes introduce per ciascun Pod un container di infrastruttura dedicato e trasparente: il **Pause Container** (detto anche \emph{sandbox container}).

% [IMMAGINE DEL FILE DI SLIDE CON NOME "Orchestrazione_1"] Diagramma del Sample Pod Network: il container "pause" funge da radice dell'infrastruttura di rete e ospita veth0, mentre container 1 e container 2 vi si agganciano per comunicare via localhost — Slide 92 / Pagina 17 del PDF

\subsubsection{Caratteristiche E Funzioni Del Pause Container}
Il container di pause assolve a compiti architetturali di fondamentale importanza:
\begin{itemize}
    \item **Creazione Delle Primitive Di Isolamento**: all'atto della schedulazione del Pod su un nodo, il runtime crea prima di tutto il pause container, istanziando i **cgroups**, i limiti di allocazione e i namespaces (in particolare il \textbf{Network Namespace} e l'\textbf{IPC Namespace});
    \item **Condivisione Degli Spazi Di Indirizzamento**: tutti i container applicativi dichiarati nel manifest del Pod vengono avviati agganciandosi ai namespace già generati e detenuti dal pause container. Grazie alla condivisione dell'IPC namespace, i container possono utilizzare meccanismi standard di comunicazione tra processi Unix, quali code di messaggi POSIX o segmenti di memoria condivisa System V;
    \item **Governo Del Ciclo Di Vita (Lifetime Management)**: il pause container non esegue codice applicativo. Il suo entrypoint è costituito da un binario minimale compilato in C che effettua unicamente la chiamata di sistema \texttt{pause()}, la quale pone il processo in stato di sospensione forzata in attesa della ricezione di un segnale di sistema (\texttt{SIGTERM}/\texttt{SIGINT}). Poiché il container rimane costantemente attivo e non fallisce mai, esso preserva l'esistenza del network namespace anche qualora i container applicativi adiacenti vadano in crash o vengano riavviati ripetutamente dallo scheduler;
    \item **Consumo Risorse Nullo**: occupando una frazione trascurabile di memoria (tipicamente poche centinaia di kilobyte) e azzerando i cicli di calcolo su CPU, il container non impatta sulle prestazioni dell'host.
\end{itemize}

\subsubsection{Ispezione Operativa Del Pause Container Con Containerd}
Il pause container viene gestito in autonomia a basso livello dai motori CRI (\textbf{containerd} o \textbf{CRI-O}) e non è visualizzabile tramite il comando standard \texttt{kubectl get pods}. Per verificarne l'esistenza fisica sul nodo worker, l'amministratore deve utilizzare l'interfaccia CLI di containerd (\texttt{ctr}) o di Docker:

\begin{lstlisting}[language=bash]
# Ispezione dei container nel namespace k8s di containerd
sudo ctr -n k8s.io container list
sudo ctr -n k8s.io ps
ctr -n k8s.io inspect my-pod-pause

# In alternativa su vecchie installazioni basate su motore Docker:
docker ps -a | grep -i pause
\end{lstlisting}

\subsection{Benefici Dell'Astrazione IP-Per-Pod Nella Migrazione Da Macchine Virtuali}

L'introduzione del modello IP-per-pod apporta enormi vantaggi ingegneristici nelle migrazioni aziendali da architetture a Macchine Virtuali (VM) verso cluster a container. In un'infrastruttura legacy, ciascuna VM disponeva del proprio indirizzo IP e poteva dialogare liberamente con le altre istanze sulla rete aziendale.

Con il modello di Kubernetes:
\begin{itemize}[noitemsep]
    \item I Pod vengono gestiti alla stregua di host fisici o VM native dal punto di vista dell'allocazione delle porte, della risoluzione dei nomi e del bilanciamento del carico;
    \item Viene eliminata la necessità di creare link espliciti tra i container o di ricorrere a complesse matrici di port-mapping statiche sull'host;
    \item Due container all'interno del medesimo Pod comunicano semplicemente invocando:
    \[
    \texttt{curl http://localhost:80}
    \]
    consentendo la coesistenza trasparente di moduli applicativi eterogenei.
\end{itemize}

\newpage

% ==============================================================================
% SEZIONE 5: NETWORKING INTER-NODE: ARCHITETTURA CNI E FLUSSO DI CREAZIONE DEL POD
% ==============================================================================

\section{Networking Inter-Node: Architettura CNI E Flusso Di Creazione Del Pod}

Se all'interno di un singolo nodo worker l'accoppiamento dei container è governato agilmente da bridge locali e dal pause container, l'estensione della connettività su scala di cluster (\emph{Out-Node Networking}) pone problemi di instradamento distribuiti significativamente più complessi.

\subsection{La Criticità Delle Collisioni Di Rete Inter-Node}

Si consideri un cluster basilare composto da due nodi worker fisici connessi tra loro da un router o switch aziendale.

% [IMMAGINE DEL FILE DI SLIDE CON NOME "Orchestrazione_1"] Topologia di rete distribuita a due nodi interconnessi da gateway fisico che evidenzia l'ambiguità degli indirizzi di pod privati locali — Slide 94 / Pagina 19 del PDF

Se ci si limitasse ad adottare il motore di containerizzazione predefinito su entrambi gli host senza alcuna gestione coordinata:
\begin{enumerate}[noitemsep]
    \item Entrambi i nodi creerebbero un bridge locale assegnandogli deterministicamente la medesima subnet privata (ad esempio \texttt{172.17.0.0/24});
    \item I Pod istanziati sul Nodo 1 e sul Nodo 2 riceverebbero indirizzi IP coincidenti (es. entrambi riceverebbero \texttt{172.17.0.2});
    \item Il router fisico dell'infrastruttura non disporrebbe di alcuna rotta per instradare un pacchetto originato sul Nodo 1 e destinato alla subnet privata del Nodo 2.
\end{enumerate}

% [IMMAGINE DEL FILE DI SLIDE CON NOME "Orchestrazione_1"] Schema del conflitto architetturale out-node: entrambi i nodi assegnano la subnet 172.17.0.0/24 in assenza di un sistema di gestione o rete di overlay — Slide 95 / Pagina 20 del PDF

Per risolvere tale ambiguità e rispettare i requisiti di assenza di NAT, Kubernetes adotta un'infrastruttura di bridge dedicati (denominati storicamente **cbr0**, \emph{custom bridge}) e demanda interamente la topologia di rete a una specifica modulare denominata **Container Network Interface (CNI)**.

\subsection{La Specifica CNI (Container Network Interface)}

Promossa all'interno della **Cloud Native Computing Foundation (CNCF)**, la specifica CNI standardizza le modalità con cui i runtime di container interagiscono con i moduli di rete.

Invece di incorporare protocolli di rete complessi all'interno del codice sorgente di Kubernetes, il control plane delega la configurazione delle interfacce a plugin esterni conformi alle direttive CNI. Un plugin CNI si assume la responsabilità di:
\begin{itemize}[noitemsep]
    \item Inserire un'estremità della coppia di interfacce \texttt{veth} all'interno del namespace di rete del container;
    \item Collegare l'altra estremità della \texttt{veth} all'infrastruttura del nodo host (ad esempio collegandola a un bridge \texttt{cbr0}, inserendola nel kernel Open vSwitch o associandola a una policy di routing eBPF);
    \item Allocare un indirizzo IP univoco a livello di cluster interfacciandosi con un modulo di gestione degli indirizzi denominato **IPAM** (\emph{IP Address Management});
    \item Configurare le rotte statiche all'interno del namespace del Pod affinché il default gateway punti all'interfaccia host;
    \item Garantire la propagazione e la visibilità globale delle rotte verso tutti gli altri nodi del cluster.
\end{itemize}

\subsection{Flusso Dettagliato Di Schedulazione Del Pod E Configurazione CNI}

L'interfaccia **CRI** (\emph{Container Runtime Interface}) costituisce il livello di astrazione che separa la logica di alto livello di Kubernetes dai motori di esecuzione fisici (\textbf{containerd}, \textbf{CRI-O}). Quando un Pod viene distribuito, l'interazione coordinata tra API Server, Scheduler, Kubelet, CRI e CNI si sviluppa attraverso i seguenti passi sequenziali:

\begin{enumerate}
    \item **Creazione Del Manifest**: l'amministratore sottopone il manifest YAML del Pod all'API Server, che ne convalida la correttezza e ne salva lo stato su etcd.
    \item **Scelta Del Nodo (Scheduling)**: il componente \textbf{kube-scheduler} rileva il Pod privo di assegnazione, valuta i requisiti di risorse (\texttt{requests}) e seleziona il nodo worker ottimale.
    \item **Ricezione Del Task Da Kubelet**: il demone \textbf{kubelet} residente sul nodo target rileva l'assegnazione e ordina al runtime locale di creare l'ambiente operativo.
    \item **Creazione Del Sandbox E Del Namespace**: il runtime CRI (es. containerd) avvia il pause container, istanziando il nuovo network namespace isolato (\texttt{netns}).
    \item **Invocazione Del Plugin (\texttt{CNI ADD})**: attraverso chiamate standard CNI, il runtime invoca il comando \texttt{CNI ADD}, trasmettendo le informazioni di identificazione del Pod e il path del relativo network namespace.
    \item **Creazione Della Coppia Veth**: il plugin CNI genera la coppia \texttt{veth}. Un'estremità entra nel namespace del Pod (venendo rinominata convenzionalmente \texttt{eth0}), mentre l'altra estremità rimane agganciata all'host (assumendo nomi identificativi quali \texttt{caliXXXX} per Calico o \texttt{vethXXXX} per bridge tradizionali).
    \item **Assegnazione Dell'Indirizzo IP (IPAM)**: il sottomodulo IPAM alloca un IP non in conflitto attinto dal blocco CIDR riservato al nodo.
    \item **Configurazione Di Rotte E Datapath**: il CNI configura l'indirizzo IP su \texttt{eth0}, attiva l'interfaccia di \texttt{loopback}, configura il default gateway del Pod e raccorda il lato host al piano dati (\emph{datapath}) del nodo.
    \item **Messa In Sicurezza Di Rotte E Policy**: l'agente CNI aggiorna le regole di firewall locale (iptables/eBPF) e le tabelle di routing inter-nodo.
    \item **Avvio Dei Container Applicativi**: una volta completata l'infrastruttura di rete del sandbox, il runtime avvia i container applicativi del Pod congiungendoli al namespace. Il Pod diviene pienamente raggiungibile da qualsiasi punto del cluster.
\end{enumerate}

\subsection{Sincronizzazione Dinamica Dello Stato Tramite Meccanismo Di HTTP Watch}

Affinché la rete distribuita converga in tempo reale, i demoni di rete CNI residenti sui nodi worker (ad esempio l'agente \texttt{flanneld} di Flannel o l'agente \texttt{felix} di Calico, generalmente distribuiti come DaemonSet di cluster) devono conoscere costantemente la mappa di allocazione delle subnet di tutti i restanti host del cluster.

Invece di interrogare l'API Server mediante polling ciclico (che provocherebbe un overhead insostenibile all'aumentare dei nodi), l'agente CNI instaura una connessione persistente sfruttando la primitiva **HTTP Watch**:

\begin{enumerate}
    \item L'agente avvia una connessione HTTP/HTTPS verso l'API Server emettendo una richiesta GET esplicita corredata dal parametro \texttt{watch=true}:
    \[
    \texttt{GET /api/v1/nodes?watch=true}
    \]
    \item Il \textbf{kube-apiserver} non chiude il canale socket; esso registra la richiesta all'interno del proprio registro eventi e risponde avviando uno stream continuo formattato in pacchetti JSON (\emph{chunked transfer encoding}).
    \item Il kube-apiserver si registra contestualmente come osservatore (\emph{watcher}) sugli indici del backend distribuito **etcd**.
    \item Nel momento esatto in cui un nuovo nodo viene annesso al cluster, o un'istanza viene rimossa, etcd notifica l'evento atomico all'API Server.
    \item L'API Server serializza l'evento e lo inoltra tempestivamente lungo la connessione HTTP aperta all'agente di nodo.
    \item L'agente riceve il flusso di eventi e aggiorna istantaneamente le tabelle di routing del kernel host, i tunnel di incapsulamento o le regole di filtraggio, garantendo una convergenza a latenza minima.
\end{enumerate}

\newpage

% ==============================================================================
% SEZIONE 6: TECNOLOGIE DI RETE A CONFRONTO: OVERLAY NETWORK VS DIRECT ROUTING
% ==============================================================================

\section{Tecnologie Di Rete A Confronto: Overlay Network Vs Direct Routing}

La realizzazione pratica dei requisiti di connettività inter-node senza NAT si articola nell'ecosistema Kubernetes su due paradigmi architetturali alternativi: le **Reti Overlay** basate su incapsulamento e le soluzioni di **Instradamento Diretto (Direct Routing)** a livello 3 puro.

% [IMMAGINE DEL FILE DI SLIDE CON NOME "Orchestrazione_1"] Schema comparativo del flusso di traffico inter-node: a sinistra, traffico incapsulato in overlay VXLAN attraverso la rete fisica; a destra, instradamento diretto (Direct Routing) con rotte distribuite — Slide 98 / Pagina 26 del PDF

\subsection{Paradigma 1: Reti Overlay (Flannel E Weave Net)}

Nelle reti di tipo **Overlay**, i pacchetti di rete scambiati tra i Pod vengono "impacchettati" (incapsulati) all'interno di normali pacchetti di livello 4 (tipicamente pacchetti UDP) prima di essere immessi sull'infrastruttura di rete fisica sottostante. La rete fisica (detta \emph{underlay network}) osserva unicamente il traffico scambiato tra gli indirizzi IP dei nodi worker fisici, ignorando l'esistenza degli indirizzi IP privati dei Pod.

Il protocollo standard prevalente per questa implementazione è **VXLAN** (\emph{Virtual Extensible LAN}), affiancato da soluzioni moderne quali **WireGuard** per garantire la cifratura integrata del canale.

\subsubsection{Flannel}
Flannel rappresenta una delle prime e più diffuse soluzioni di rete open source sviluppate per Kubernetes:
\begin{itemize}[noitemsep]
    \item Configura una rete IPv4 di livello 3 operando mediante l'agente \textbf{flanneld} in esecuzione su ciascun nodo;
    \item Si appoggia interamente sull'API di Kubernetes (e quindi sullo storage di etcd) per la memorizzazione dello stato, senza richiedere database esterni dedicati;
    \item Alloca a ciascun host un prefisso di rete esclusivo (tipicamente una subnet \texttt{/24} estratta dalla macro-rete \texttt{/16} di cluster);
    \item Crea un'interfaccia virtuale VXLAN per connettere tutti i nodi: quando il Pod sul Nodo 1 trasmette un pacchetto verso il Pod sul Nodo 2, il modulo VXLAN del kernel cattura il frame Ethernet originale del Pod, lo incapsula all'interno di un pacchetto UDP avente come destinazione l'IP fisico del Nodo 2, e lo instrada sull'underlay network; all'arrivo, il kernel del Nodo 2 decapsula il pacchetto UDP ed eroga il payload originario al Pod destinatario.
\end{itemize}

\subsubsection{Weave Net}
Weave è un altro noto add-on CNI che realizza una rete di overlay mesh peer-to-peer tra tutti i container del cluster. Si distingue storicamente per la semplicità di configurazione e la capacità di operare in ambienti distribuiti o multi-cloud eterogenei, offrendo il supporto nativo per la cifratura trasparente di tutto il traffico inter-nodo.

\paragraph{Vantaggi E Limiti Delle Reti Overlay:}
Il vantaggio primario risiede nell'indipendenza totale dall'infrastruttura di rete sottostante: l'overlay funziona su qualsiasi ambiente (cloud pubblico, on-premise, ambienti virtualizzati VMware) senza richiedere la cooperazione degli apparati di rete (switch/router) fisici. Il limite è rappresentato dall'**overhead computazionale e di frame**: l'incapsulamento dei pacchetti comporta consumo di cicli CPU per l'elaborazione degli header e riduce la \emph{Maximum Transmission Unit} (MTU) effettiva a disposizione del carico utile (un header VXLAN sottrae 50 byte al pacchetto), provocando potenziale frammentazione dei pacchetti se non opportunamente compensata.

\subsection{Paradigma 2: Instradamento Diretto E Flat L3 (Project Calico)}

Per eliminare l'overhead dell'incapsulamento e garantire le massime prestazioni di throughput e latenza, il progetto **Calico** implementa un'architettura di instradamento di **Livello 3 puro (Direct Routing)**.

Calico tratta ciascun nodo worker come se fosse un vero e proprio router di transito IP (\textbf{vRouter}). I pacchetti generati dai Pod non vengono incapsulati né subiscono NAT: essi viaggiano sulla rete locale esattamente come pacchetti IP nativi.

\subsubsection{I Componenti Architetturali Di Calico}
L'infrastruttura distribuita di Calico si fonda su tre moduli cooperanti:
\begin{enumerate}
    \item **Felix**: è l'agente primario installato come demone su ciascun nodo del cluster. Felix monitora l'API Server di Kubernetes, alloca gli indirizzi IP per i Pod locali, programma le tabelle di instradamento (\emph{routing tables}) nel kernel Linux del nodo e scrive direttamente le regole di filtraggio del firewall (utilizzando \texttt{iptables}, \texttt{ipset} o programmi \textbf{eBPF}) per applicare le **NetworkPolicy** di sicurezza.
    \item **BIRD (BGP Routing Daemon)**: è un motore di instradamento dinamico open source che implementa il protocollo standard di routing su scala Internet: il **Border Gateway Protocol (BGP)**. BIRD raccoglie le rotte locali generate da Felix (es. "la subnet \texttt{10.244.1.0/24} appartiene al Nodo 1") e le propaga in tempo reale a tutti gli altri demoni BIRD degli altri nodi mediante sessioni di \emph{BGP peering} (architettura \emph{full-mesh} o basata su \emph{Route Reflectors} centralizzati, oppure effettuando il peering diretto con gli switch Top-of-Rack del data center).
    \item **confd**: è un demone di sincronizzazione leggero che monitora il database etcd alla ricerca di variazioni nella configurazione di rete BGP, rigenerando dinamicamente i file di configurazione per il demone BIRD e forzandone il ricaricamento a caldo.
\end{enumerate}

\paragraph{Flessibilità Di Incapsulamento:} Qualora il traffico tra nodi debba attraversare segmenti di rete non cooperanti (ad esempio router intermedi che scartano pacchetti diretti a subnet non registrate nelle routing tables aziendali), Calico offre la modalità ibrida di incapsulamento **IP-in-IP** o **VXLAN**, garantendo il funzionamento su qualsiasi underlay network.

\subsection{Integrazione Cloud Nativa: AWS VPC CNI}

Nei contesti cloud proprietari (come Amazon Elastic Kubernetes Service - EKS), l'uso di overlay o BGP su larga scala può essere sostituito da plugin che integrano strettamente il networking di cluster con il software-defined network del fornitore:
\begin{itemize}[noitemsep]
    \item Il plugin **AWS VPC CNI** alloca indirizzi IP per i Pod prelevandoli direttamente dalla subnet del **Virtual Private Cloud (VPC)** di AWS;
    \item Il plugin aggancia interfacce di rete elastiche secondarie (\textbf{ENI}) alle istanze EC2 e associa gli IP secondari ai singoli Pod;
    \item In questo modo, i Pod diventano entità pienamente native della rete cloud aziendale: essi possono interagire con database gestiti (es. Amazon RDS) e servizi interni senza transitare da proxy o NAT, beneficiando nativamente delle regole di sicurezza e dei controlli di accesso definiti a livello di **AWS Security Groups** e policy IAM.
\end{itemize}

\subsection{Quadro Sinottico Delle Soluzioni Di Networking In Kubernetes}

A conclusione della trattazione sul networking dei sistemi distribuiti orchestrati da Kubernetes, la seguente tabella riepiloga le caratteristiche strutturali, i paradigmi di instradamento e gli ambiti operativi delle principali tecnologie analizzate:

\begin{center}
\begin{tabularx}{\textwidth}{l >{\raggedright\arraybackslash}p{4.2cm} c X}
    \toprule
    \textbf{Plugin CNI} & \textbf{Meccanismo di Trasporto} & \textbf{Incapsulamento} & \textbf{Caso d'Uso Elettivo} \\
    \midrule
    \textbf{Flannel} & Overlay L3 (VXLAN, host-gw) & Sì (VXLAN) & Cluster semplici, ambienti bare-metal senza requisiti complessi di policy. \\
    \addlinespace
    \textbf{Calico} & vRouter L3 con BGP nativo (o IP-in-IP) & Opzionale & Ambienti ad altissime prestazioni, data center multi-rack, Network Policies avanzate. \\
    \addlinespace
    \textbf{Weave Net} & Overlay mesh peer-to-peer & Sì & Installazioni rapide, ambienti distribuiti con cifratura automatica integrata. \\
    \addlinespace
    \textbf{AWS VPC CNI} & Diretto nativo su ENI VPC & No & Cluster eseguiti su AWS che richiedono integrazione con Security Groups e VPC routing. \\
    \bottomrule
\end{tabularx}
\end{center}
% AA Orchestrazione_1 [133-165] ------------------------------------------------------------

\section{Evoluzione Dei Runtime Di Container E Networking CNI: Docker vs. Containerd}

Nelle moderne architetture cloud-native orchestrate tramite Kubernetes, l'esecuzione dei carichi di lavoro poggia su un disaccoppiamento fondamentale tra il piano di orchestrazione e il runtime dei container. Storicamente, Docker rappresentava la soluzione monolitica predefinita; tuttavia, l'evoluzione dello standard \textbf{CRI} (\emph{Container Runtime Interface}) ha portato all'adozione prevalente di \textbf{containerd}, un runtime minimale e incentrato sulle specifiche aperte OCI (\emph{Open Container Initiative}). Tale transizione architetturale ha avuto un impatto diretto sul sottosistema di rete dei nodi di elaborazione.

\subsection{Richiamo Concettuale: Coppie Virtual Ethernet (veth) E Bridge Linux}
Prima di analizzare le differenze a livello di runtime, è essenziale rinfrescare il funzionamento del networking a livello di kernel Linux nei container:
\begin{itemize}[noitemsep]
    \item \textbf{Network Namespace (netns)}: Fornisce a un processo (o a un insieme di container facenti capo allo stesso Pod) uno stack di rete completamente isolato, provvisto delle proprie tabelle di instradamento, regole di filtro (iptables/nftables) e interfacce di rete.
    \item \textbf{Coppia veth (\emph{Virtual Ethernet Pair})}: Funge da ``cavo di rete virtuale'' bidirezionale punto-a-punto. Un estremo della coppia risiede all'interno del network namespace del Pod (convenzionalmente rinominato \texttt{eth0}), mentre l'altro estremo rimane nel namespace root dell'host (con nomi quali \texttt{veth...} o \texttt{cali...}). Qualsiasi frame inviato a un estremo emerge istantaneamente all'estremo opposto.
    \item \textbf{Linux Bridge}: Uno switch software di Livello 2 implementato nel kernel che connette diverse interfacce virtuali, inoltrando i frame Ethernet in base all'apprendimento degli indirizzi MAC (\emph{MAC learning}).
\end{itemize}

\subsection{Confronto Architetturale Tra Docker E Containerd}
La divergenza sostanziale tra Docker e containerd risiede nella rigidità della gestione delle interfacce e nell'aderenza allo standard \textbf{CNI} (\emph{Container Network Interface}).

% [IMMAGINE DEL FILE DI SLIDE CON NOME "Calico-Kubernetes-Services.pdf"] Tabella comparativa tra i runtime Docker e Containerd rispetto a Network Management, veth, Pause container, flessibilità e CNI — Slide 102 / Pagina 1 del PDF

\begin{table}[H]
\centering
\small
\begin{tabularx}{\textwidth}{l p{5.5cm} X}
\toprule
\textbf{Caratteristica} & \textbf{Docker} & \textbf{Containerd} \\
\midrule
\textbf{Network Management} & Utilizza per impostazione predefinita un bridge virtuale Linux (\texttt{docker0}). & Si affida interamente a plugin CNI configurabili (\emph{Flannel}, \emph{Calico}, \emph{Cilium}, ecc.). \\
\addlinespace
\textbf{Interfacce veth} & Associa le coppie \emph{veth} al bridge \texttt{docker0} per instradare il traffico. & Crea e collega le coppie \emph{veth}, ma la configurazione e il posizionamento dipendono dal plugin CNI. \\
\addlinespace
\textbf{Container Pause} & Mantiene un container di \emph{pause} dedicato per ciascun Pod per agganciare il netns. & Supporta il container di \emph{pause}, delegando l'infrastruttura di rete al plugin CNI. \\
\addlinespace
\textbf{Flessibilità} & Rigida: topologie predefinite basate principalmente su bridge nativo. & Elevata: ampiamente personalizzabile tramite topologie CNI (overlay, routed puro, eBPF). \\
\addlinespace
\textbf{Gestione Plugin CNI} & Possibile ma non nativa (richiedeva shim intermedi come \emph{dockershim}). & Nativa e integrata tramite specifiche CNI dirette. \\
\bottomrule
\end{tabularx}
\caption{Confronto sistematico tra i modelli di networking in Docker e Containerd.}
\label{tab:docker-vs-containerd}
\end{table}

Nel paradigma tradizionale di Docker, il demone crea all'avvio il bridge virtuale \texttt{docker0}, a cui assegna un indirizzo IP privato (tipicamente \texttt{172.17.0.1/16}). Quando un container viene instanziato, Docker vi connette un estremo della coppia \emph{veth} e colloca l'altro all'interno di \texttt{docker0}, mascherando il traffico uscente tramite regole di \textbf{Source NAT (SNAT/Masquerade)} gestite dal motore iptables dell'host. Sebbene questa configurazione consenta un'operatività \emph{out-of-the-box} per installazioni a singolo nodo, essa mal si adatta alle rigide specifiche di rete di Kubernetes, le quali prescrivono che:
\begin{enumerate}[noitemsep]
    \item Ogni Pod possegga un indirizzo IP unico e instradabile nell'intero cluster.
    \item Tutti i Pod possano comunicare direttamente tra loro senza ricorrere al NAT.
    \item Gli agenti di nodo (come il Kubelet) possano comunicare con tutti i Pod presenti sullo stesso host.
\end{enumerate}

Containerd elimina completamente l'imposizione del bridge statico \texttt{docker0}. Demandando la configurazione a una catena di plugin CNI eseguiti all'atto della creazione del namespace di rete del Pod (durante l'avvio del container di \emph{pause}), containerd permette all'amministratore di implementare topologie avanzate: bridge dedicati, interfacce routed dirette di Livello 3 o tunnel incapsulati overlay.

\newpage

% =============================================================================
% SEZIONE: CNI CALICO
% =============================================================================

\section{Architettura Del CNI Calico: Modelli Di Inoltro Intra-Nodo E Inter-Nodo}

\textbf{Project Calico} è uno dei plugin CNI più performanti e ampiamente adottati nell'ecosistema Kubernetes. A differenza delle soluzioni convenzionali basate su overlay incapsulati o bridge L2 complessi, Calico adotta primariamente un approccio di \textbf{routing puro di Livello 3 (L3)}, trasformando ciascun nodo del cluster in un vero e proprio router IP.

\subsection{Assegnazione Degli Indirizzi IP (IPAM)}
Calico assegna a ciascun nodo una porzione riservata dello spazio di indirizzamento complessivo dei Pod (\texttt{podCIDR}), tipicamente una maschera \texttt{/24} o \texttt{/26} ritagliata dal blocco aggregato del cluster (es. \texttt{10.244.0.0/16} o \texttt{10.233.64.0/18}). Ogni Pod creato su un determinato nodo riceve un indirizzo univoco appartenente al CIDR dedicato a quell'host.

\subsection{Comunicazione Intra-Nodo: Routing Puro Senza Bridge Virtuali}
Quando due Pod risiedono sul medesimo nodo, la comunicazione non richiede il passaggio attraverso la rete fisica esterna, ma non ricorre neppure a un tradizionale bridge Linux (\texttt{cbr0} o \texttt{docker0}) né a switch software avanzati come \textbf{Open vSwitch (OVS)}.

All'atto del bootstrap del Pod, Calico genera la coppia di interfacce \emph{veth}:
\begin{itemize}[noitemsep]
    \item Un estremo viene spostato nel namespace di rete del Pod e rinominato \texttt{eth0}.
    \item L'estremo opposto rimane nel namespace root dell'host e assume un identificatore univoco del tipo \texttt{caliXXXXXX} (ad esempio \texttt{cali0}, \texttt{cali1234} o hash casuali come \texttt{calie2861d758c8}).
\end{itemize}

% [IMMAGINE DEL FILE DI SLIDE CON NOME "Calico-Kubernetes-Services.pdf"] Schema architetturale del networking intra-nodo con Calico: interfaccia caliX direttamente attestata al layer IP dell'host senza bridge — Slide 104 / Pagina 3 del PDF

L'interfaccia \texttt{caliX} non viene collegata ad alcun bridge. Al contrario, il kernel dell'host opera come un router IP puro: per instradare i pacchetti tra i Pod locali, il kernel consulta direttamente la propria routing table di Livello 3, dove sono definite rotte puntuali (\texttt{/32} per IPv4) verso ciascun indirizzo di Pod, associate alla specifica interfaccia \texttt{caliX}.

I vantaggi architetturali di questa scelta rispetto a un bridge tradizionale o a OVS sono rilevanti:
\begin{itemize}[noitemsep]
    \item \textbf{Prestazioni elevate}: Il traffico attraversa le pipeline di \emph{IP forwarding} native del kernel Linux senza contestualizzazioni di commutazione L2 software.
    \item \textbf{Minimo overhead computazionale}: Vengono completamente eliminate le tabelle di inoltro MAC (\emph{Forwarding Information Base - MAC tables}) e la gestione delle regole complesse di \emph{OpenFlow}.
    \item \textbf{Semplificazione operativa}: L'intera visibilità del traffico e l'applicazione dei criteri di sicurezza (\emph{Network Policies}) sono gestite mediante le primitive convenzionali di routing e filtraggio del kernel Linux (iptables, ipset o eBPF).
\end{itemize}

\subsection{Comunicazione Inter-Nodo: Instradamento Nativo Con BGP}
Quando due Pod risiedono su nodi fisici distinti, Calico consente di trasportare il traffico senza alcuna incapsulazione, sfruttando il protocollo standard dell'infrastruttura Internet: il \textbf{Border Gateway Protocol (BGP)}.

\subsubsection{Componenti Del Control Plane BGP}
Su ogni nodo Kubernetes, Calico esegue un agente modulare incapsulato nel Pod DaemonSet \texttt{calico-node}:
\begin{itemize}[noitemsep]
    \item \textbf{Felix}: È il demone primario di Calico. Monitora il datastore (etcd o le API di Kubernetes) e programma la tabella di instradamento del kernel Linux (tramite chiamate Netlink), i filtri iptables/ipset e le interfacce virtuali sull'host.
    \item \textbf{Bird (o GoBGP)}: È un demone di routing BGP dinamico. Ciascuna istanza di Bird stabilisce sessioni di \emph{peering} BGP (in modalità \emph{full-mesh} tra i nodi o mediante \emph{BGP Route Reflector}) con gli altri nodi o con i router fisici del datacenter (\emph{Top-of-Rack switch}), annunciando le subnet dei Pod locali (\texttt{podCIDR}) assegnate al proprio nodo.
\end{itemize}

% [IMMAGINE DEL FILE DI SLIDE CON NOME "Calico-Kubernetes-Services.pdf"] Schema della propagazione delle rotte con BGP tra nodi Kubernetes tramite il demone Bird — Slide 105 / Pagina 4 del PDF

Il ciclo di vita dell'instradamento avviene attraverso quattro passaggi sequenziali:
\begin{enumerate}[noitemsep]
    \item \textbf{Assegnazione CIDR}: Il controller di rete assegna a ciascun nodo un blocco dedicato (es. \texttt{10.244.1.0/24} al Nodo 1 e \texttt{10.244.2.0/24} al Nodo 2).
    \item \textbf{Annuncio BGP}: Il demone Bird sul Nodo 1 diffonde un messaggio di annuncio BGP (\emph{Update}) ai suoi peer, notificando che la subnet \texttt{10.244.1.0/24} è raggiungibile attraverso il suo indirizzo IP fisico (es. \texttt{192.168.1.2}).
    \item \textbf{Aggiornamento della Routing Table Host}: Bird riceve gli annunci dai nodi remoti e istruisce il kernel Linux ad aggiungere una rotta statica nella tabella principale dell'host.
    \item \textbf{Inoltro Diretto}: Quando un Pod sul Nodo 1 genera un pacchetto per un Pod sul Nodo 2, il pacchetto esce dall'interfaccia fisica \texttt{eth0} dell'host mantenendo gli indirizzi IP sorgente e destinazione originari dei Pod, senza alcun header aggiuntivo.
\end{enumerate}

\subsubsection{Analisi Di Una Tabella Di Instradamento BGP}
Consideriamo l'output del comando \texttt{ip route} visualizzato su un nodo Kubernetes con indirizzo fisico \texttt{192.168.1.10} all'interno di una rete \texttt{192.168.0.0/16}:

\begin{lstlisting}[language=bash]
$ ip route
default via 192.168.1.1 dev eth0
10.244.1.0/24 via 192.168.1.2 dev eth0 # Route to Pods on another node
10.244.2.0/24 via 192.168.1.3 dev eth0 # Route to Pods on another node
192.168.0.0/16 dev eth0 proto kernel scope link src 192.168.1.10 # direct route
\end{lstlisting}

La decodifica ingegneristica delle voci di instradamento rivela quanto segue:
\begin{itemize}[noitemsep]
    \item \texttt{default via 192.168.1.1 dev eth0}: Rotta predefinita dell'host per raggiungere reti esterne al cluster mediante il gateway fisico \texttt{192.168.1.1}.
    \item \texttt{10.244.1.0/24 via 192.168.1.2 dev eth0}: Notifica che l'intero blocco di indirizzi \texttt{10.244.1.0/24}, associato ai Pod ospitati sul nodo remoto con IP \texttt{192.168.1.2}, deve essere inoltrato direttamente a quest'ultimo attraverso l'interfaccia fisica \texttt{eth0}.
    \item \texttt{10.244.2.0/24 via 192.168.1.3 dev eth0}: Specifica analogamente che i Pod della subnet \texttt{10.244.2.0/24} risiedono sul nodo \texttt{192.168.1.3}.
    \item \texttt{192.168.0.0/16 dev eth0 proto kernel scope link src 192.168.1.10}: Rotta di collegamento locale (\emph{scope link}) che definisce la raggiungibilità diretta di Livello 2 per tutti i nodi fisici appartenenti alla subnet del cluster, indicando \texttt{192.168.1.10} come indirizzo IP sorgente dell'host locale.
\end{itemize}

Questo modello garantisce le massime prestazioni di trasmissione (throughput pari alla velocità di linea del ferro e latenza minima), ma richiede che l'infrastruttura di rete sottostante (\emph{underlay network}) consenta il traffico di pacchetti con IP sorgente/destinazione privati dei Pod, oppure supporti sessioni di peering BGP con i router aziendali.

\subsection{Comunicazione Inter-Nodo Tramite Incapsulamento IP-in-IP}
Qualora i router fisici intermedi non supportino BGP, oppure l'infrastruttura di rete scarti i pacchetti con indirizzi IP privati dei Pod (comportamento frequente nei cloud provider pubblici a causa delle verifiche anti-spoofing), Calico offre la modalità di incapsulamento \textbf{IP-in-IP} (\emph{IPIP}).

\subsubsection{Principio Di Funzionamento Dell'Incapsulamento IPIP}
Nel modello IPIP, il pacchetto IP generato dal Pod sorgente viene interamente incapsulato all'interno di un secondo pacchetto IP esterno (protocollo IP numero 4):
\begin{align*}
    \text{Header IP Esterno: } & \left[ \text{Src: } \texttt{192.168.1.10} \ (\text{Nodo 1}), \ \text{Dst: } \texttt{192.168.1.20} \ (\text{Nodo 2}) \right] \\
    \text{Header IP Interno: } & \left[ \text{Src: } \texttt{10.244.1.10} \ (\text{Pod A}), \ \text{Dst: } \texttt{10.244.2.20} \ (\text{Pod B}) \right] \\
    \text{Payload: } & \left[ \text{Dati applicativi TCP/UDP} \right]
\end{align*}

Quando il pacchetto incapsulato raggiunge il nodo di destinazione (\texttt{192.168.1.20}), il modulo kernel dell'host rimuove l'header IP esterno (\emph{decapsulazione}) e reinstrada il pacchetto originario verso l'interfaccia virtuale \texttt{caliX} del Pod ricevente.

\subsubsection{Analisi Della Tabella Di Instradamento IP-in-IP}
La presenza del tunnel IPIP modifica radicalmente le voci del comando \texttt{ip route}:

\begin{lstlisting}[language=bash]
$ ip route
10.244.1.0/24 dev cali0 proto kernel scope link src 10.244.1.1 # pod network in the node
10.244.2.0/24 via 192.168.1.20 dev tunl0 onlink # route to pods in a different node
192.168.1.0/24 dev eth0 proto kernel scope link src 192.168.1.10 # direct route
\end{lstlisting}

Dall'analisi delle direttive emergono dettagli critici:
\begin{itemize}[noitemsep]
    \item \texttt{10.244.1.0/24 dev cali0 proto kernel scope link src 10.244.1.1}: La subnet locale dei Pod (\texttt{10.244.1.0/24}) è gestita localmente tramite l'interfaccia virtuale \texttt{cali0}, alla quale è assegnato l'indirizzo IP di gateway locale \texttt{10.244.1.1}.
    \item \texttt{10.244.2.0/24 via 192.168.1.20 dev tunl0 onlink}: Il traffico destinato alla subnet remota \texttt{10.244.2.0/24} viene instradato attraverso l'interfaccia di tunneling \textbf{\texttt{tunl0}}. L'opzione \textbf{\texttt{onlink}} indica al kernel di considerare il gateway di next-hop (\texttt{192.168.1.20}) come direttamente raggiungibile sul link sottostante, forzando l'invio del pacchetto attraverso il driver del tunnel senza richiedere una corrispondenza diretta nella tabella di routing convenzionale.
    \item \texttt{192.168.1.0/24 dev eth0 proto kernel scope link src 192.168.1.10}: Definisce l'adattatore fisico \texttt{eth0} per la comunicazione broadcast/unicast tra gli indirizzi IP dei nodi dell'infrastruttura.
\end{itemize}

\subsubsection{Distribuzione Delle Rotte In Assenza Di BGP}
In assenza del demone BGP (Bird), il piano di controllo di Calico adotta un meccanismo centralizzato basato sul datastore:
\begin{enumerate}[noitemsep]
    \item All'ingresso di un nuovo nodo, il Kubelet registra il nodo presso l'API server di Kubernetes, che gli alloca un \texttt{PodCIDR} (ad esempio \texttt{10.244.1.0/24}).
    \item Il componente \textbf{Calico Kube-Controllers} o il DaemonSet locale rileva la modifica e salva la mappatura \texttt{Nodo $\to$ podCIDR} all'interno del datastore di Calico (configurabile direttamente sulle risorse CRD dell'API server di Kubernetes oppure su un cluster \textbf{etcd} dedicato).
    \item L'agente locale \textbf{Felix} (eventualmente coadiuvato da \textbf{Typha}, un proxy che riduce il carico delle connessioni concorrenti verso il datastore nei grandi cluster) interroga continuamente tali mappature.
    \item Felix scrive direttamente nel kernel le rotte che associano ciascun \texttt{podCIDR} remoto all'interfaccia \texttt{tunl0}, specificando come next-hop l'indirizzo fisico del nodo corrispondente.
\end{enumerate}

\subsection{Comunicazione Inter-Nodo Tramite Overlay VXLAN}
Un'alternativa moderna a IP-in-IP è rappresentata da \textbf{VXLAN} (\emph{Virtual Extensible LAN}). Mentre IP-in-IP incapsula pacchetti L3 in L3, VXLAN crea una rete overlay di Livello 2 incapsulando frame Ethernet/IP all'interno di datagrammi \textbf{UDP} (generalmente indirizzati alla porta 4789).

\subsubsection{Architettura VXLAN In Calico}
Ciascun nodo crea un'interfaccia di tunneling dedicata denominata \textbf{\texttt{vxlan.calico}}. All'interno dell'header VXLAN risiede il campo \textbf{VNI} (\emph{Virtual Network Identifier}), uno spazio a 24 bit utilizzato tipicamente per isolare fino a 16 milioni di segmenti virtuali in contesti multi-tenant (sebbene in un cluster Kubernetes standard venga generalmente impiegato un singolo VNI condiviso per tutti i Pod).

% [IMMAGINE DEL FILE DI SLIDE CON NOME "Calico-Kubernetes-Services.pdf"] Schema del pacchetto incapsulato VXLAN tra Pod su nodi differenti con VNI — Slide 112 / Pagina 15 del PDF

Il flusso di comunicazione procede come segue:
\begin{enumerate}[noitemsep]
    \item Il Pod sorgente (\texttt{10.244.1.10}) genera un pacchetto per il Pod remoto (\texttt{10.244.2.20}).
    \item La routing table dell'host intercetta la destinazione e la invia all'interfaccia \texttt{vxlan.calico}.
    \item L'interfaccia \texttt{vxlan.calico} incapsula il pacchetto aggiungendo un header VXLAN (contenente il VNI), seguito da un header di trasporto UDP e un header IP esterno con:
    \begin{itemize}[noitemsep]
        \item \emph{Outer IP Source}: \texttt{192.168.1.10} (Nodo A).
        \item \emph{Outer IP Destination}: \texttt{192.168.1.20} (Nodo B).
    \end{itemize}
    \item Il pacchetto attraversa l'infrastruttura fisica mediante instradamento standard. Il binding tra l'IP del Pod di destinazione e l'IP fisico dell'host remoto viene risolto tramite ARP proxy, tabelle FDB (\emph{Forwarding Database}) gestite da Felix, oppure meccanismi multicast.
    \item Il nodo ricevente decapsula il datagramma UDP/VXLAN e recapita il pacchetto originale al Pod locale.
\end{enumerate}

\subsubsection{Analisi Della Tabella Di Instradamento VXLAN}
\begin{lstlisting}[language=bash]
$ ip route
10.244.1.0/24 dev cali0 proto kernel scope link src 10.244.1.1 # pod network in the node
10.244.2.0/24 dev vxlan.calico onlink # route to pods in a different node
192.168.1.0/24 dev eth0 proto kernel scope link src 192.168.1.10
\end{lstlisting}

In questo scenario, la voce relativa ai Pod remoti non specifica un singolo indirizzo di next-hop esplicito (\texttt{via IP}), bensì delega l'intero blocco \texttt{10.244.2.0/24} all'interfaccia \texttt{vxlan.calico onlink}. Sarà il sottosistema VXLAN del kernel, opportunamente popolato da Felix con le tabelle di adiacenza, a determinare l'IP di destinazione del tunnel UDP (\emph{VTEP - VXLAN Tunnel End Point}) corrispondente.

\subsection{Quadro Comparativo: BGP vs. IP-in-IP vs. VXLAN}
La scelta del modello di instradamento in Calico è guidata da precisi vincoli infrastrutturali:

\begin{table}[H]
\centering
\small
\begin{tabularx}{\textwidth}{l X X X}
\toprule
\textbf{Parametro} & \textbf{BGP Diretto} & \textbf{IP-in-IP (IPIP)} & \textbf{VXLAN} \\
\midrule
\textbf{Incapsulamento} & Nessuno (Routing nativo L3). & Header IP esterno (20 byte). & Header UDP + VXLAN (50 byte complessivi). \\
\addlinespace
\textbf{Overhead / MTU} & Zero overhead; MTU piena dell'interfaccia fisica. & Basso (richiede riduzione MTU di 20 byte). & Medio (richiede riduzione MTU di 50 byte). \\
\addlinespace
\textbf{Latenza} & Minima (prestazioni \emph{near-wire-speed}). & Leggera penalizzazione (CPU per incapsulamento). & Penalizzazione analoga a IPIP (overhead parsing UDP). \\
\addlinespace
\textbf{Supporto Rete} & Richiede abilitazione BGP o fabric L2 trasparente. & Funziona ovunque passi traffico IP standard. & Supportato nativamente su apparati moderni e cloud. \\
\addlinespace
\textbf{Casi D'Uso Tipici} & Installazioni \emph{bare-metal} e datacenter privati. & Reti multi-cloud prive di supporto BGP diretto. & Ambienti enterprise complessi o dove IPIP è bloccato. \\
\bottomrule
\end{tabularx}
\caption{Confronto sinottico tra le modalità di inoltro supportate dal CNI Calico.}
\label{tab:calico-modes}
\end{table}

\newpage

% =============================================================================
% SEZIONE: KUBERNETES SERVICE
% =============================================================================

\section{L'Astrazione Di Servizio In Kubernetes (Basic Objects: Service)}

\subsection{Il Problema Dell'Effimerità Dei Pod E Della Scalabilità Orizzontale}
Nel modello operativo di Kubernetes, i Pod sono considerati \textbf{risorse effimere}: nascono, possono essere distrutti o migrati su altri nodi per riallocazione delle risorse, fallimento hardware o eventi di ridimensionamento (\emph{horizontal autoscaling}).

Ciascun Pod riceve un indirizzo IP dinamico allocato dal CNI soltanto \emph{dopo} essere stato schedulato su un nodo e prima del suo avvio formale. Di conseguenza:
\begin{itemize}[noitemsep]
    \item L'indirizzo IP di un Pod non può essere previsto anticipatamente.
    \item Se un Pod termina e viene ricreato da un controller, riceverà un nuovo indirizzo IP.
    \item In presenza di più repliche dello stesso carico applicativo, ciascuna istanza possiede un proprio indirizzo distinto.
\end{itemize}

Nel modello convenzionale di amministrazione di sistema, un client verrebbe configurato cablando l'indirizzo IP o il nome host del server all'interno dei propri file di configurazione. In un'architettura a container distribuita, questo schema provocherebbe il collasso istantaneo della comunicazione al primo ciclo di riavvio o scalabilità.

\subsection{Il Concetto Di Reverse Proxy vs. Forward Proxy}
La soluzione architetturale adottata dai progettisti di Kubernetes consiste nell'introduzione dell'oggetto \textbf{Service}, che agisce concettualmente come un \textbf{reverse proxy} resiliente e distribuito.

% [IMMAGINE DEL FILE DI SLIDE CON NOME "Calico-Kubernetes-Services.pdf"] Schema concettuale che confronta il ruolo del Forward Proxy e del Reverse Proxy nel flusso di traffico — Slide 114 / Pagina 17 del PDF

Per comprendere a fondo la differenza architetturale:
\begin{itemize}[noitemsep]
    \item \textbf{Forward Proxy}: Si colloca tra un gruppo di client (es. i computer degli utenti in una rete aziendale) e Internet. Il client contatta il proxy, il quale inoltra la richiesta al web server esterno. Viene utilizzato primariamente per filtraggio dei contenuti, anonimizzazione o caching sul lato client.
    \item \textbf{Reverse Proxy}: Si colloca tra i client esterni e una \emph{farm} di server applicativi di backend. Il client si connette esclusivamente all'indirizzo del reverse proxy credendo di dialogare con il server reale. Il proxy intercetta la richiesta, seleziona una tra le istanze di backend disponibili secondo una politica di bilanciamento del carico, inoltra la chiamata e restituisce la risposta al client. Il reverse proxy agisce dunque \emph{per conto del server}.
\end{itemize}

L'oggetto \textbf{Service} di Kubernetes concretizza questo pattern: espone un endpoint stabile (indirizzo IP virtuale denominato \textbf{ClusterIP} e un record DNS canonico) e implementa il load balancing trasparente verso l'insieme dei Pod di backend sani. Il ciclo di vita dei Pod e quello del Service sono completamente indipendenti: i Pod possono morire e rigenerarsi con nuovi indirizzi IP senza che il client debba modificare la propria configurazione o perda la connettività.

\subsection{Implementazione Pratica: Deployment E Service Di Esempio}
Per illustrare il meccanismo, consideriamo l'attivazione di due Pod server HTTP tramite un controller di tipo \textbf{Deployment}, descritto nel manifest seguente:

\begin{lstlisting}[language=yaml, caption={Manifest di Deployment (\texttt{deployment-service-test.yml})}]
apiVersion: apps/v1
kind: Deployment
metadata:
  name: service-test
spec:
  replicas: 2
  selector:
    matchLabels:
      app: service_test_pod
  template:
    metadata:
      labels:
        app: service_test_pod
    spec:
      containers:
      - name: simple-http
        image: python:2.7
        command: ["/bin/bash"]
        args: ["-c", "echo \"<p>Hello from $(hostname)</p>\" > index.html; python -m SimpleHTTPServer 8080"]
        ports:
        - name: http
          containerPort: 8080
\end{lstlisting}

All'avvio, ogni container esegue uno script bash che genera un file \texttt{index.html} contenente il proprio hostname univoco (corrispondente al nome del Pod) e attiva un server HTTP minimale in ascolto sulla porta TCP \texttt{8080}. L'etichetta associata al template dei Pod è \texttt{app: service\_test\_pod}.

Per raggruppare stabilmente questi Pod dietro un unico punto di ingresso, si definisce il manifest dell'oggetto \textbf{Service}:

\begin{lstlisting}[language=yaml, caption={Manifest del Service ClusterIP associato (\texttt{deployment-service-test.yml})}]
apiVersion: v1
kind: Service
metadata:
  name: service-test
spec:
  selector:
    app: service_test_pod
  ports:
  - port: 80
    targetPort: http
\end{lstlisting}

\subsubsection{Distinzione Fondamentale Tra \texttt{port} E \texttt{targetPort}}
La configurazione dei campi di porta richiede estrema chiarezza formale:
\begin{itemize}[noitemsep]
    \item \textbf{\texttt{port}}: Rappresenta la porta astratta esposta dal \textbf{Service} all'interno del cluster. I client invieranno le proprie richieste TCP/UDP all'indirizzo ClusterIP su questa specifica porta (nel nostro esempio, la porta \texttt{80}).
    \item \textbf{\texttt{targetPort}}: Rappresenta la porta reale su cui il processo all'interno del container è in effettivo ascolto all'interno del proprio network namespace. Tale parametro può essere specificato come valore numerico (es. \texttt{8080}) oppure come stringa mnemonica (es. \texttt{http}), agganciandosi al campo \texttt{ports.name} dichiarato all'interno della specifica del container nel Deployment. Nel nostro caso, \texttt{targetPort: http} risolve puntualmente alla porta \texttt{8080} del container.
\end{itemize}

\begin{theorem}
L'oggetto \textbf{Endpoints} (o l'evoluzione \textbf{EndpointSlice}) è la risorsa Kubernetes che materializza l'associazione dinamica tra il Service e i Pod reali. Monitorando costantemente i Pod i cui label coincidono con il \texttt{selector} del Service, il control plane popola l'elenco degli indirizzi IP e porte effettivi dei container pronti a ricevere traffico:
\[
    \texttt{Endpoints: 10.233.75.36:8080, 10.233.75.48:8080}
\]
\end{theorem}

\subsubsection{Metodi Di Verifica Della Connettività}
L'accesso e la verifica del servizio possono avvenire tramite molteplici canali:
\begin{enumerate}[noitemsep]
    \item \textbf{Esecuzione diretta all'interno del Pod}:
    \begin{lstlisting}[language=bash]
$ kubectl exec -it service-test-12345-abcde -- curl http://localhost:8080
    \end{lstlisting}
    \item \textbf{Creazione di un tunnel tramite port-forwarding locale}:
    \begin{lstlisting}[language=bash]
$ kubectl port-forward service/service-test 8080:80
# In una shell separata:
$ curl http://localhost:8080
    \end{lstlisting}
    Il comando stabilisce un tunnel inoltrando le connessioni dalla porta \texttt{8080} della macchina di amministrazione alla porta \texttt{80} del Service nel cluster, che a sua volta instrada il traffico verso la \texttt{targetPort 8080} di uno dei backend.
    \item \textbf{Interrogazione diretta del ClusterIP o del nome di servizio}:
    \begin{lstlisting}[language=bash]
$ curl http://<clusterIP>:80
$ curl http://service-test:80
    \end{lstlisting}
\end{enumerate}

\newpage

% =============================================================================
% SEZIONE: SERVICE DISCOVERY E DNS
% =============================================================================

\section{Service Discovery E Risoluzione DNS Nel Cluster}

Affinché i componenti applicativi possano comunicare in modo trasparente, Kubernetes integra un sottosistema di \textbf{Service Discovery} nativo imperniato su un server DNS interno. Invece di dover conoscere l'indirizzo ClusterIP (anch'esso generato dinamicamente al momento dell'applicazione del manifest), i Pod client possono utilizzare nomi di dominio mnemonici e costanti.

\subsection{Architettura Del DNS Di Cluster: CoreDNS vs. Kube-DNS}
In passato, l'infrastruttura DNS di Kubernetes era affidata a \textbf{Kube-DNS}. A partire dalla versione 1.12, \textbf{CoreDNS} è diventato il server DNS predefinito e raccomandato, grazie alla sua architettura modulare scritta in Go, alla maggiore efficienza nell'uso della memoria e al supporto flessibile di plugin.

Entrambi gli add-on espongono il medesimo contratto architetturale:
\begin{itemize}[noitemsep]
    \item Vengono eseguiti come Pod gestiti all'interno del namespace \texttt{kube-system}.
    \item Vengono pubblicati mediante un Service ClusterIP con un indirizzo statico e immutabile (es. \texttt{10.96.0.10}).
    \item Per garantire la retrocompatibilità, il Service associato a CoreDNS mantiene formalmente il nome \texttt{kube-dns} nel campo \texttt{metadata.name}.
    \item All'atto del provisioning di ciascun container, il demone \textbf{Kubelet} inietta automaticamente la configurazione di risoluzione nel file \texttt{/etc/resolv.conf}, valorizzando i flag passati all'avvio del nodo: \texttt{--cluster-dns=<dns-service-ip>} e \texttt{--cluster-domain=<default-local-domain>}.
\end{itemize}

\subsection{Configurazione Di Risoluzione Nel Container: Il File \texttt{resolv.conf}}
Accedendo a una shell all'interno di un generico Pod di test (ad esempio basato su un'immagine minimale come \texttt{busybox}), è possibile ispezionare il resolver di sistema:

\begin{lstlisting}[language=bash]
$ kubectl run busybox --image=busybox -it --rm -- /bin/sh
/ # cat /etc/resolv.conf
search default.svc.cluster.local svc.cluster.local cluster.local
nameserver 10.96.0.10
options ndots:5
\end{lstlisting}

L'interpretazione dettagliata dei parametri svela il comportamento del client:
\begin{itemize}[noitemsep]
    \item \textbf{\texttt{nameserver 10.96.0.10}}: Rappresenta l'indirizzo ClusterIP del servizio DNS interno del cluster. Tutte le query generate dai processi del container vengono dirette a questo server sulla porta UDP/TCP 53.
    \item \textbf{\texttt{search default.svc.cluster.local svc.cluster.local cluster.local}}: Definisce la lista di ricerca per l'espansione dei suffissi di dominio qualificati (\emph{FQDN - Fully Qualified Domain Name}). Se un'applicazione tenta di risolvere il nome semplice \texttt{service-test}, il resolver accoderà progressivamente i suffissi specificati fino a trovare una corrispondenza positiva:
    \[
        \texttt{service-test.default.svc.cluster.local} \implies \text{Risoluzione con successo!}
    \]
    \item \textbf{\texttt{options ndots:5}}: Indica che se un nome interrogato contiene un numero di punti (\emph{dots}) inferiore a 5, il resolver considererà il nome come non-FQDN e tenterà preventivamente la risoluzione applicando i domini della direttiva \texttt{search}.
\end{itemize}

\subsection{Verifica Empirica Del Bilanciamento Del Carico}
Per testare la correttezza della risoluzione dei nomi e dimostrare l'azione di bilanciamento statistico del Service, consideriamo la creazione di un Pod client interattivo dedicato:

\begin{lstlisting}[language=yaml, caption={Manifest del Pod client di test (\texttt{pod-dns.yaml})}]
apiVersion: v1
kind: Pod
metadata:
  name: service-test-client
spec:
  restartPolicy: Never
  containers:
  - name: test-client2
    image: busybox
    command: ["/bin/sh"]
    args: ["-c", "echo -e 'GET / HTTP/1.1\r\n\r\n' | nc -v service-test 80"]
\end{lstlisting}

Il comando invia una stringa conforme al protocollo HTTP/1.1 (conclusa dalla sequenza di ritorno carrello e avanzamento riga \texttt{\textbackslash r\textbackslash n\textbackslash r\textbackslash n}, corrispondente allo standard ASCII \texttt{<CR><LF>}) tramite l'utility \textbf{Netcat} (\texttt{nc}), aprendo una connessione TCP verso la porta \texttt{80} del nome simbolico \texttt{service-test}.

Eseguendo ripetutamente il Pod o un ciclo di richieste similari, l'ispezione dei log:
\begin{lstlisting}[language=bash]
$ kubectl logs service-test-client
\end{lstlisting}
mostrerà in output le risposte alternate:
\begin{lstlisting}[language=html]
<p>Hello from service-test-586c84576b-2mxsf</p>
<!-- alla chiamata successiva: -->
<p>Hello from service-test-586c84576b-sjxrl</p>
\end{lstlisting}

Ciascun server backend riceve circa il $50\%$ delle richieste totali. Questo risultato evidenzia come il ClusterIP non instrada staticamente verso un unico container, ma attua un bilanciamento del carico tra tutte le repliche sane censite nel relativo oggetto Endpoints.

\newpage

% =============================================================================
% SEZIONE: CLUSTER IP MANAGEMENT E DATA PLANE
% =============================================================================

\section{Gestione Del ClusterIP E Meccanismi Di Inoltro Nel Data Plane: Kube-Proxy E Netfilter}

\subsection{La Natura Virtuale Della Rete Di Servizio (Service Network)}
Una delle peculiarità concettualmente più complesse del networking di Kubernetes riguarda la natura dell'indirizzo IP assegnato a un Service (\textbf{ClusterIP}).

Quando si esamina l'output di un'interrogazione di dettaglio del servizio:
\begin{lstlisting}[language=bash]
$ kubectl describe service/service-test
Name:              service-test
Namespace:         default
Selector:          app=service_test_pod
Type:              ClusterIP
IP:                10.233.51.173
Port:              <unset> 80/TCP
TargetPort:        http/TCP
Endpoints:         10.233.75.36:8080, 10.233.75.48:8080
\end{lstlisting}
si nota che al Service è associato l'IP \texttt{10.233.51.173}. Tuttavia, analizzando la tabella di routing del kernel Linux sul nodo host (\texttt{node1}):

\begin{lstlisting}[language=bash]
vncc@node1:~/esempi$ ip route
default via 192.168.43.1 dev enp0s3 proto static metric 100
10.233.75.0/26 via 10.233.75.0 dev vxlan.calico onlink
blackhole 10.233.102.128/26 proto 80
10.233.102.156 dev calie2861d758c8 scope link
10.233.102.157 dev cali079ca33d357 scope link
10.233.102.158 dev cali43e5dfd19a9 scope link
192.168.43.0/24 dev enp0s3 proto kernel scope link src 192.168.43.10 metric 100
\end{lstlisting}

Si riscontra un fatto sorprendente: \textbf{la rete \texttt{10.233.51.0} non compare da nessuna parte nella tabella di instradamento del kernel}.
\begin{itemize}[noitemsep]
    \item Non esiste alcuna interfaccia fisica o virtuale (come \texttt{eth0}, \texttt{cbr0} o \texttt{cali0}) che possieda l'indirizzo \texttt{10.233.51.173}.
    \item Non esistono rotte di rete per il subnet del ClusterIP.
    \item Se si esamina il gateway fisico di rete, esso non ha alcuna conoscenza di tale subnet.
\end{itemize}
Ciononostante, inviando un comando \texttt{curl http://10.233.51.173:80}, il sistema risponde prontamente restituendo la pagina HTML del Pod backend. Come è possibile che un pacchetto indirizzato a un'interfaccia inesistente raggiunga la sua destinazione?

\subsection{Evoluzione Di Kube-Proxy: Da Userspace A Kernel-Level}
La gestione di questa illusione architetturale è demandata a \textbf{kube-proxy}, un componente di Kubernetes eseguito come DaemonSet su ciascun nodo del cluster.

% [IMMAGINE DEL FILE DI SLIDE CON NOME "Calico-Kubernetes-Services.pdf"] Schema del pacchetto catturato in volo dal modulo netfilter e rediretto da kube-proxy — Slide 124 / Pagina 31 del PDF

Nel corso dell'evoluzione di Kubernetes, kube-proxy ha adottato diverse modalità operative:
\begin{enumerate}[noitemsep]
    \item \textbf{Userspace Proxy Mode (Storica)}: Kube-proxy apriva una porta reale in ascolto nello spazio utente sull'host per ciascun Service. Quando un client inviava traffico al ClusterIP, regole iptables rudimentali catturavano il pacchetto e lo ridirigevano a kube-proxy in user-space. Il processo leggeva i dati, selezionava un backend e ritrasmetteva la connessione al Pod reale. Questo approccio comportava un gravoso costo di \emph{context switch} continuo tra spazio kernel e spazio utente per ogni singolo pacchetto, limitando drasticamente la banda passante.
    \item \textbf{Iptables Mode (Predefinita e Attuale)}: Kube-proxy cessa di essere un proxy nel percorso dati. Esso si limita a operare come controllore nel piano di controllo locale: ascolta le modifiche ai Service ed Endpoints dall'API server e programma direttamente una serie di catene e tabelle all'interno del sottosistema **Netfilter** del kernel Linux. Il pacchetto non abbandona mai lo spazio kernel.
    \item \textbf{IPVS / eBPF Mode}: Nelle versioni avanzate per cluster su vastissima scala, la programmazione iptables (che opera con scansione sequenziale lineare $O(N)$) viene sostituita da hash table in IPVS o programmi eBPF caricati direttamente nei punti di aggancio del driver di rete, garantendo complessità di ricerca $O(1)$.
\end{enumerate}

\subsection{L'Infrastruttura Netfilter Ed Il Meccanismo Delle Catene Iptables}
\textbf{Netfilter} è un framework integrato nel kernel Linux che fornisce una serie di \emph{hook} (punti di aggancio) all'interno dello stack di rete per intercettare, manipolare, filtrare o alterare l'indirizzamento dei pacchetti IP prima e dopo le decisioni di routing. \textbf{Iptables} è l'interfaccia utente con cui vengono configurate tali regole.

Quando un processo client all'interno di un Pod genera una richiesta HTTP verso il ClusterIP (es. \texttt{10.3.241.152:80}):
\begin{enumerate}[noitemsep]
    \item Il Pod emette il pacchetto attraverso la sua interfaccia \texttt{eth0}.
    \item Il pacchetto attraversa la coppia \emph{veth} ed entra nello stack di rete dell'host.
    \item Poiché l'indirizzo IP \texttt{10.3.241.152} non appartiene a nessuna interfaccia locale, il kernel si preparerebbe a inoltrare il pacchetto al proprio default gateway (uscendo verso l'interfaccia fisica).
    \item Tuttavia, prima che ciò avvenga, il pacchetto attraversa gli hook di Netfilter (in particolare la tabella \textbf{NAT} nella catena \textbf{PREROUTING} o \textbf{OUTPUT}).
    \item Il pacchetto viene letteralmente ``catturato in volo'' (\emph{snagged in flight}) dalle regole pre-programmate da kube-proxy.
\end{enumerate}

% [IMMAGINE DEL FILE DI SLIDE CON NOME "Calico-Kubernetes-Services.pdf"] Diagramma di flusso dettagliato dell'attraversamento delle catene iptables KUBE-SERVICES, KUBE-SVC e KUBE-SEP con DNAT — Slide 125 / Pagina 33 del PDF

\subsection{Tracciamento Passo-Passo Del Flusso Iptables (DNAT Probabilistico)}
Kube-proxy struttura le regole della tabella NAT suddividendole in una gerarchia di catene dedicate:

\begin{figure}[H]
\centering
\small
\begin{tikzpicture}[node distance=1.5cm, auto,
    block/.style={rectangle, draw=theoryblue, fill=backcolour, text width=7.5cm, rounded corners, minimum height=1cm, align=left},
    arrow/.style={-Latex, thick, darkgreen}]

    \node [block] (input) {\textbf{Pacchetto in ingresso allo stack host:}\\Dst: \texttt{10.3.241.152:80} (ClusterIP)};
    \node [block, below of=input] (services) {\textbf{Catena \texttt{KUBE-SERVICES}}\\Intercetta le chiamate verso qualsiasi Service IP del cluster. Conduce a \texttt{KUBE-SVC-X}.};
    \node [block, below of=services, node distance=1.8cm] (svc) {\textbf{Catena \texttt{KUBE-SVC-X} (Bilanciamento Statistico)}\\Seleziona uno dei backend attivi mediante il modulo probabilistico del kernel.};
    \node [block, below of=svc, node distance=1.8cm] (sep) {\textbf{Catena \texttt{KUBE-SEP-Y} (Service Endpoint)}\\Applica la regola di \textbf{DNAT} verso il Pod reale: \texttt{to-destination 10.0.2.2:8080}.};
    \node [block, below of=sep, node distance=1.6cm] (output) {\textbf{Decisione di Routing del Kernel:}\\Dst alterata in: \texttt{10.0.2.2:8080}. Il pacchetto viene instradato verso il Pod corrispondente.};

    \path [arrow] (input) -- (services);
    \path [arrow] (services) -- node [right, font=\footnotesize] {Match ClusterIP} (svc);
    \path [arrow] (svc) -- node [right, font=\footnotesize] {Regola probabilistica} (sep);
    \path [arrow] (sep) -- (output);
\end{tikzpicture}
\caption{Flusso di risoluzione di un ClusterIP attraverso le catene iptables.}
\label{fig:iptables-flow}
\end{figure}

Nel dettaglio analitico del diagramma di slide 125:
\begin{enumerate}
    \item \textbf{Catena \texttt{KUBE-SERVICES}}:
    Il pacchetto transita nella catena primaria che valuta la corrispondenza dell'indirizzo IP di destinazione:
    \begin{lstlisting}[language=bash]
-A KUBE-SERVICES -d 10.3.241.152/32 -p tcp --dport 80 -j KUBE-SVC-X
    \end{lstlisting}
    Se l'indirizzo coincide con il ClusterIP e la porta coincide con la \texttt{port: 80}, il flusso salta alla catena specifica del servizio \texttt{KUBE-SVC-X}.

    \item \textbf{Catena \texttt{KUBE-SVC-X} e Modulo Statistico}:
    All'interno di questa catena risiede l'algoritmo di bilanciamento del carico tra le $N$ repliche disponibili (nel nostro caso $N=2$). Il kernel valuta la scelta tramite il modulo statistico (\texttt{-m statistic --mode random}):
    \begin{lstlisting}[language=bash]
-A KUBE-SVC-X -m statistic --mode random --probability 0.5 -j KUBE-SEP-1
-A KUBE-SVC-X -j KUBE-SEP-2
    \end{lstlisting}
    La prima regola viene valutata con probabilità $p = \frac{1}{N} = 0.5$. Se la condizione casuale ha successo, il salto porta alla catena del primo endpoint (\texttt{KUBE-SEP-1}); altrimenti, il pacchetto prosegue linearmente e incappa nella seconda regola, che cattura deterministicamente il restante $50\%$ del traffico inviandolo a \texttt{KUBE-SEP-2}.

    \item \textbf{Catena \texttt{KUBE-SEP-Y} e Traduzione DNAT}:
    Ciascuna catena di endpoint applica l'operazione di \textbf{Destination NAT (DNAT)}:
    \begin{lstlisting}[language=bash]
-A KUBE-SEP-1 -p tcp -j DNAT --to-destination 10.0.2.2:8080
    \end{lstlisting}
    L'header IP del pacchetto viene riscritto in tempo reale: l'indirizzo virtuale di destinazione \texttt{10.3.241.152:80} viene sostituito con l'indirizzo reale del Pod backend \texttt{10.0.2.2:8080}.

    \item \textbf{Instradamento Verso il Backend}:
    A questo punto, la pipeline di routing del kernel Linux si attiva \emph{dopo} il DNAT. Esaminando la nuova destinazione (\texttt{10.0.2.2}), il kernel riscontra una rotta valida: se il Pod risiede sullo stesso nodo, lo instraderà direttamente all'interfaccia \texttt{caliX}; se si trova su un nodo remoto, applicherà le regole BGP o l'incapsulamento IPIP/VXLAN analizzate nelle sezioni precedenti.

    \item \textbf{Gestione Del Flusso Di Ritorno (\emph{Conntrack})}:
    Al momento della ricezione della risposta da parte del server backend, il modulo \textbf{conntrack} (\emph{Connection Tracking}) del kernel Linux intercetta il pacchetto di ritorno e compie l'operazione inversa in modo del tutto simmetrico e trasparente: riscrive l'IP sorgente da \texttt{10.0.2.2:8080} a \texttt{10.3.241.152:80}. Il client originario riceve così una risposta formalmente coerente con l'indirizzo verso cui aveva avviato la sessione TCP.
\end{enumerate}

Grazie a questo raffinato connubio tra controllo logico (Kubelet e Kube-proxy) e accelerazione a livello di kernel (Netfilter e CNI Calico), l'infrastruttura di Kubernetes garantisce un'astrazione di rete scalabile, dinamica e resiliente, svincolando completamente gli sviluppatori e i carichi applicativi dalla complessità dell'infrastruttura sottostante.


% AA Orchestrazione_1 [166-198] ------------------------------------------------------------

\section{Meccanismi Avanzati Di Data Plane Per ClusterIP: Kube-Proxy, Netfilter E Conntrack}

Nel modello di rete di Kubernetes, l'astrazione di un \textbf{Service} di tipo \textbf{ClusterIP} rappresenta un punto di accesso logico e immutabile verso un insieme dinamico di Pod. Sebbene a livello logico il ClusterIP si comporti come un indirizzo IP convenzionale, a livello di sistema operativo esso non corrisponde ad alcuna interfaccia di rete fisica o virtuale, né è associato ad alcuna scheda di rete dell'host. L'intero meccanismo di instradamento è realizzato mediante una manipolazione dinamica delle tabelle del kernel Linux gestita da \textbf{kube-proxy}.

\subsection{Architettura Di Inoltro Nel Kernel: Il Ruolo Di Kube-Proxy}
Nelle prime versioni di Kubernetes, kube-proxy operava come un proxy a livello applicativo nello spazio utente (\emph{userspace proxy}). In tale modalità, il demone apriva una porta TCP/UDP locale sull'host per ciascun Service registrato e utilizzava regole di filtro elementari per convogliare le connessioni verso di sé. Kube-proxy leggeva i pacchetti in arrivo, selezionava un backend disponibile ed effettuava una nuova connessione verso il Pod di destinazione. Tale modello soffriva di un costo computazionale proibitivo: ogni singolo pacchetto richiedeva molteplici cambi di contesto (\emph{context switch}) tra spazio kernel e spazio utente, introducendo latenze elevate e un forte collo di bottiglia sul throughput.

% [IMMAGINE DEL FILE DI SLIDE CON NOME "Calico-Kubernetes-Services.pdf"] Schema architetturale del flusso ClusterIP con kube-proxy fuori dal data path e manipolazione kernel via iptables — Slide 126 / Pagina 1 del PDF

Nelle implementazioni moderne basate su \textbf{iptables} o \textbf{IPVS}, kube-proxy è stato completamente rimosso dal percorso dati (\emph{data path}):
\begin{itemize}[noitemsep]
    \item Kube-proxy agisce unicamente nel \textbf{piano di controllo locale} (\emph{control plane}), eseguito come \textbf{DaemonSet} su ciascun nodo del cluster.
    \item Interroga continuamente l'API server di Kubernetes per rilevare la creazione, l'aggiornamento o la distruzione di oggetti \texttt{Service} ed \texttt{Endpoints} (o \texttt{EndpointSlices}).
    \item All'atto di ogni variazione topologica, programma e aggiorna istantaneamente le catene e le tabelle del framework **Netfilter** all'interno del kernel Linux.
    \item L'intero ciclo di elaborazione dei pacchetti (traslazione degli indirizzi, selezione del backend e instradamento) viene eseguito nativamente dal sottosistema di rete del kernel a velocità di linea (\emph{wire-speed}).
\end{itemize}

\subsection{Ispezione Delle Catene Iptables: Il Flusso Da \texttt{KUBE-SERVICES} A \texttt{KUBE-SEP}}
Per comprendere la reale sequenza operativa con cui il kernel processa una richiesta destinata a un ClusterIP, è possibile esaminare direttamente le tabelle Netfilter generate sul nodo mediante il comando \texttt{iptables}.

% [IMMAGINE DEL FILE DI SLIDE CON NOME "Calico-Kubernetes-Services.pdf"] Diagramma a blocchi delle catene KUBE-SERVICES, KUBE-SVC e KUBE-SEP con la logica di smistamento e load balancing — Slide 128 / Pagina 4 del PDF

Consideriamo un servizio di test denominato \texttt{service-test} associato al ClusterIP \texttt{10.96.199.67} sulla porta TCP \texttt{80}, supportato da due repliche di Pod con indirizzi IP \texttt{10.244.0.8:8080} e \texttt{10.244.0.9:8080}. L'ispezione della tabella NAT rivela la seguente gerarchia:

\begin{lstlisting}[language=bash, caption={Ispezione delle regole iptables tramite kube-proxy}]
$ kubectl exec -n kube-system kube-proxy-t2nnw -- iptables -t nat -L KUBE-SERVICES -n
KUBE-SVC-GBTMB5ZTD56ECL5F  6  --  0.0.0.0/0  10.96.199.67  /* default/service-test cluster IP */ tcp dpt:80

$ kubectl exec -n kube-system kube-proxy-t2nnw -- iptables -t nat -L KUBE-SVC-GBTMB5ZTD56ECL5F -n
Chain KUBE-SVC-GBTMB5ZTD56ECL5F (1 references)
target                     prot opt source          destination
KUBE-MARK-MASQ             6    -- !10.244.0.0/16   10.96.199.67  /* default/service-test cluster IP */ tcp dpt:80
KUBE-SEP-E2TCQP2POLEEKAHJ  0    --  0.0.0.0/0       0.0.0.0/0     /* default/service-test -> 10.244.0.8:8080 */ statistic mode random probability 0.50000000000
KUBE-SEP-4RNZELUSBQ4FXW7B  0    --  0.0.0.0/0       0.0.0.0/0     /* default/service-test -> 10.244.0.9:8080 */
\end{lstlisting}

L'analisi strutturata dei parametri di queste direttive illustra la logica formale:
\begin{enumerate}
    \item \textbf{Ingresso nella catena \texttt{KUBE-SERVICES}}:
    Tutto il traffico in transito o generato localmente attraversa preliminarmente questa catena. La regola controlla il protocollo IP di livello 4 (il parametro numerico \texttt{6} identifica formalmente il protocollo \textbf{TCP}), l'indirizzo di destinazione (\texttt{10.96.199.67}) e la porta di destinazione (\texttt{dpt:80}). In caso di corrispondenza, il flusso salta immediatamente alla catena specifica del servizio, identificata dall'hash univoco \texttt{KUBE-SVC-GBTMB5ZTD56ECL5F}.

    \item \textbf{Valutazione del traffico sorgente e marcatura (\texttt{KUBE-MARK-MASQ})}:
    La prima riga all'interno della catena di servizio analizza la provenienza del pacchetto: se la sorgente \emph{non} appartiene alla subnet complessiva dei Pod (\texttt{!10.244.0.0/16}), il pacchetto viene indirizzato alla catena \texttt{KUBE-MARK-MASQ}, la quale applica un flag di marcatura Netfilter (convenzionalmente il bitmask \texttt{0x4000}). Questa marcatura predispone il pacchetto alla successiva operazione di mascheramento (\textbf{Source NAT / MASQUERADE}) nella fase di post-routing.

    \item \textbf{Bilanciamento del carico probabilistico nei \texttt{KUBE-SEP-*}}:
    Le regole successive implementano il bilanciamento del carico a livello di kernel utilizzando l'estensione \texttt{statistic} di iptables in modalità casuale (\texttt{mode random}):
    \begin{itemize}[noitemsep]
        \item Per $N$ endpoint disponibili, la $k$-esima regola viene programmata con probabilità condizionata:
        \[
            P(\text{Match al passo } k) = \frac{1}{N - k + 1}, \quad \text{per } k = 1, \dots, N-1
        \]
        \item Nel caso di due repliche ($N=2$), la prima regola adotta una probabilità $p = 0.5$. Se il generatore pseudo-casuale del kernel valida la condizione, il pacchetto salta alla catena \texttt{KUBE-SEP-E2TC...}.
        \item Se la prima regola fallisce, il pacchetto scende linearmente alla regola successiva. Essendo l'ultima della catena, essa non presenta vincoli probabilistici e cattura deterministicamente il restante $50\%$ delle connessioni verso \texttt{KUBE-SEP-4RNZ...}.
    \end{itemize}

    \item \textbf{Applicazione del Destination NAT (\texttt{DNAT})}:
    Ciascuna catena \texttt{KUBE-SEP-*} (\emph{Service EndPoint}) contiene una singola direttiva che riscrive l'header IP del pacchetto, alterando la tupla di destinazione:
    \[
        \underbrace{\texttt{10.96.199.67:80}}_{\text{ClusterIP virtuale}} \xrightarrow{\text{DNAT}} \underbrace{\texttt{10.244.0.8:8080} \quad \text{oppure} \quad \texttt{10.244.0.9:8080}}_{\text{IP reale del Pod e porta del Container}}
    \]
\end{enumerate}

\subsection{La Necessità Del Source NAT (SNAT): Traffico Interno vs. Esterno}
L'applicazione del Source NAT (SNAT) risponde a un preciso vincolo di simmetria del protocollo TCP. Si distinguono due scenari operativi:

\subsubsection{Caso 1: Comunicazione Pod-to-Service Interna (Nessun SNAT)}
Quando un Pod sorgente (\texttt{10.244.0.5}) invia una richiesta al ClusterIP \texttt{10.96.199.67:80}:
\begin{itemize}[noitemsep]
    \item Il kernel del nodo applica unicamente il DNAT, riscrivendo la destinazione verso il Pod di backend (\texttt{10.244.0.8:8080}).
    \item L'indirizzo IP sorgente rimane inalterato (\texttt{10.244.0.5}).
    \item Poiché sia il client sia il server appartengono alla medesima rete overlay/routed dei Pod (\texttt{10.244.0.0/16}), il Pod di destinazione riceve il pacchetto e risponde direttamente all'IP sorgente \texttt{10.244.0.5}.
    \item Il modulo di tracciamento delle connessioni del nodo di partenza re-intercetta il traffico di risposta e ripristina come mittente il ClusterIP, chiudendo il cerchio TCP in modo coerente. In questo flusso \textbf{non è richiesto alcun SNAT}.
\end{itemize}

\subsubsection{Caso 2: Comunicazione Esterna o Da Host (Necessità Di SNAT/Masquerade)}
Consideriamo invece una richiesta originata da un host fisico o da una macchina esterna alla rete dei Pod (\texttt{10.100.0.2}) diretta al ClusterIP:
\begin{itemize}[noitemsep]
    \item Se il nodo applicasse esclusivamente il DNAT verso il Pod \texttt{10.244.0.8}, il pacchetto consegnato al Pod avrebbe:
    \[
        \text{Sorgente: } \texttt{10.100.0.2}, \quad \text{Destinazione: } \texttt{10.244.0.8}
    \]
    \item Al momento di emettere il segmento di risposta \texttt{TCP SYN-ACK}, il Pod consulterebbe la propria tabella di routing e instraderebbe il pacchetto direttamente verso l'indirizzo del client \texttt{10.100.0.2}.
    \item Il client riceverebbe un pacchetto con IP sorgente \texttt{10.244.0.8}. Tuttavia, il socket sul client è in attesa di risposta dall'indirizzo a cui ha inviato la richiesta iniziale (\texttt{10.96.199.67}). Riscontrando un disallineamento della tupla di connessione (\emph{address mismatch}), il client invierebbe immediatamente un pacchetto \texttt{TCP RST}, abbattendo la connessione.
\end{itemize}

Per prevenire questo fenomeno di instradamento asimmetrico, la regola \texttt{KUBE-MARK-MASQ} intercetta il traffico non appartenente a \texttt{10.244.0.0/16} e attiva la riscrittura SNAT: l'indirizzo sorgente viene sostituito con l'indirizzo IP del nodo Kubernetes da cui transita il traffico. In tal modo, il Pod invierà obbligatoriamente la risposta al nodo, il quale potrà invertire sia lo SNAT sia il DNAT prima di recapitare i dati al client esterno.

\subsection{Il Modulo Di Connessione Di Stato: Conntrack}
L'intera architettura NAT di Kubernetes collasserebbe se il kernel non disponesse di una memoria di stato in grado di gestire i flussi bidirezionali. Tale funzione è garantita dal modulo **conntrack** (\emph{Connection Tracking}) del kernel Linux.

Per ciascuna sessione attiva che attraversa il sottosistema di rete, conntrack alloca e mantiene una voce all'interno di una tabella hash in memoria, registrando:
\begin{itemize}[noitemsep]
    \item \textbf{Tupla identificativa del flusso originario}: Protocollo di trasporto (TCP, UDP, ICMP), indirizzo IP sorgente, indirizzo IP destinazione, porta sorgente e porta destinazione.
    \item \textbf{Stato della connessione}: Ad esempio \texttt{NEW} (richiesta iniziale), \texttt{ESTABLISHED} (flusso bidirezionale confermato) o \texttt{RELATED} (flussi ausiliari).
    \item \textbf{Trasformazioni NAT applicate}: Registra la mappatura esatta delle riscritture di indirizzo e porta effettuate al passaggio del primo pacchetto.
\end{itemize}

Quando un pacchetto di risposta attraversa a ritroso il percorso di rete, il modulo conntrack riconosce la sessione associata, bypassa la rivalutazione sequenziale delle catene iptables e applica automaticamente la traslazione inversa (\emph{Reverse NAT}), garantendo trasparenza assoluta ed elevata efficienza prestazionale.

\newpage

% =============================================================================
% SEZIONE: ESPOSIZIONE DI LIVELLO 4 (NODEPORT E LOADBALANCER)
% =============================================================================

\section{Esposizione Del Traffico Esterno Di Livello 4: NodePort E LoadBalancer}

\subsection{Limiti Architetturali Dell'Esposizione Diretta Di ClusterIP}
Poiché un servizio ClusterIP non è raggiungibile dall'esterno del cluster, sorge l'esigenza di consentire ai client esterni di consumare i carichi di lavoro ospitati sui Pod.

Una soluzione apparentemente intuitiva potrebbe consistere nell'assegnare al ClusterIP una rotta statica sui router fisici per instradare il traffico verso uno dei nodi del cluster. Tuttavia, questa configurazione presenta criticità insormontabili:
\begin{itemize}[noitemsep]
    \item \textbf{Effimerità dei nodi}: Sebbene meno volatili dei Pod, i nodi fisici o le macchine virtuali del cluster possono andare incontro a riavvii, migrazioni o guasti hardware. I router fisici non possiedono la visibilità dello stato applicativo interno e continuerebbero a instradare traffico verso un nodo disconnesso, interrompendo il servizio.
    \item \textbf{Collo di bottiglia e asimmetria}: Concentrare l'intero traffico di ingresso su un singolo nodo saturerebbe la relativa interfaccia di rete e degraderebbe le prestazioni globali.
    \item \textbf{Isolamento di sicurezza}: I nodi di controllo (\emph{master/control-plane}) sono deputati alla gestione del cluster e non devono in alcun caso processare traffico applicativo proveniente dall'esterno.
\end{itemize}

% [IMMAGINE DEL FILE DI SLIDE CON NOME "Calico-Kubernetes-Services.pdf"] Schema del routing con rotta statica manuale verso un singolo nodo e relative criticità architetturali — Slide 130 / Pagina 7 del PDF

Per risolvere tale problematica in maniera controllata e nativa, Kubernetes implementa le estensioni di servizio di Livello 4: **NodePort** e **LoadBalancer**.

\subsection{L'Oggetto Di Servizio NodePort}
Il servizio di tipo \textbf{NodePort} rappresenta l'estensione primaria di un ClusterIP: quando viene creato un servizio NodePort, Kubernetes alloca automaticamente anche un sottostante ClusterIP (e il relativo set di Endpoints).

% [IMMAGINE DEL FILE DI SLIDE CON NOME "Calico-Kubernetes-Services.pdf"] Architettura di rete del servizio NodePort con Load Balancer esterno che distribuisce il traffico sui nodi fisici — Slide 131 / Pagina 9 del PDF

\subsubsection{Meccanismo Di Allocazione E Funzionamento}
All'atto della definizione del servizio, il control plane assegna una porta dedicata, denominata \textbf{NodePort}, ritagliata all'interno di un intervallo riservato a livello di cluster, configurato mediante il flag \texttt{--service-node-port-range} (il cui valore predefinito è compreso tra \textbf{30000} e \textbf{32767}).

Le proprietà cardine di questo modello sono:
\begin{itemize}[noitemsep]
    \item La porta assegnata viene aperta ed esposta sull'indirizzo IP fisico di \textbf{tutti i nodi} del cluster (siano essi nodi di elaborazione che ospitano i Pod target o nodi privi di istanze locali).
    \item Qualsiasi pacchetto inviato a \texttt{<Indirizzo\_IP\_Nodo>:<NodePort>} viene intercettato dalle regole Netfilter di quel nodo, incapsulato nella logica del corrispondente ClusterIP, tradotto tramite DNAT verso uno dei backend sani e, qualora il Pod selezionato risieda su un host differente, inoltrato attraverso la rete overlay (es. VXLAN o Calico routed).
\end{itemize}

\subsubsection{Manifest Di Un Servizio NodePort}
\begin{lstlisting}[language=yaml, caption={Definizione di un servizio NodePort (\texttt{nodeport1.yml})}]
apiVersion: v1
kind: Service
metadata:
  name: service-test
spec:
  type: NodePort
  selector:
    app: service_test_pod
  ports:
  - port: 80
    targetPort: http
    nodePort: 32213  # Parametro opzionale; se omesso, viene allocato dinamicamente
\end{lstlisting}

Interrogando il cluster dopo l'applicazione:
\begin{lstlisting}[language=bash]
$ kubectl get svc service-test
NAME           TYPE       CLUSTER-IP      EXTERNAL-IP   PORT(S)        AGE
service-test   NodePort   10.233.51.173   <none>        80:32213/TCP   30h
\end{lstlisting}
Un client esterno può raggiungere il servizio contattando indifferentemente l'IP del Nodo 1 o del Nodo 2:
\[
    \texttt{http://10.100.0.2:32213} \quad \text{oppure} \quad \texttt{http://10.100.0.3:32213}
\]

\subsection{Tassonomia Delle Porte In Kubernetes}
La coesistenza di molteplici livelli di astrazione richiede una distinzione analitica tra le cinque tipologie di porte coinvolte nella comunicazione client-server:

% [IMMAGINE DEL FILE DI SLIDE CON NOME "Calico-Kubernetes-Services.pdf"] Schema grafico della correlazione tra Application Listening Port, ContainerPort, TargetPort, Service Port e NodePort — Slide 134 / Pagina 13 del PDF

\begin{enumerate}
    \item \textbf{Application Listening Port}: È la porta effettiva su cui il demone applicativo all'interno del container si mette in ascolto (\emph{bind} del socket di rete, es. porta \texttt{5000} o \texttt{8080}).
    \item \textbf{Container Port (\texttt{containerPort})}: Dichiarata all'interno della specifica del container nel Deployment (\texttt{spec.containers[].ports[].containerPort}). Funge da metadato documentativo e descrittivo all'interno del cluster per indicare la porta esposta dal container.
    \item \textbf{Target Port (\texttt{targetPort})}: Specificata nel manifest del Service (\texttt{spec.ports[].targetPort}). Indica la porta del Pod verso cui il Service deve redirigere il traffico. Se omessa, il suo valore assume di default lo stesso numero specificato nel campo \texttt{port}. Può essere dichiarata come numero intero o come stringa (riferimento a un nome di porta).
    \item \textbf{Service Port (\texttt{port})}: È la porta logica interna del Service (\texttt{spec.ports[].port}) associata all'indirizzo ClusterIP. Tutti gli altri carichi interni al cluster inviano richieste a questa porta (tipicamente la porta standard \texttt{80} per HTTP o \texttt{443} per HTTPS).
    \item \textbf{NodePort (\texttt{nodePort})}: È la porta esposta su tutte le interfacce fisiche esterne dei nodi (\texttt{spec.ports[].nodePort}), allocata nell'intervallo \texttt{30000-32767}.
\end{enumerate}

\begin{lstlisting}[language=yaml, caption={Esempio integrato di correlazione delle porte}]
# Estratto dal Deployment applicativo:
spec:
  containers:
  - name: my-app-container
    image: my-app-image
    ports:
    - containerPort: 5000  # Coincide con la porta reale del processo

---
# Estratto dal Service di tipo NodePort:
spec:
  type: NodePort
  ports:
  - port: 80             # Porta ClusterIP per i client interni
    targetPort: 5000     # Reindirizza alla containerPort dei Pod
    nodePort: 30080      # Porta aperta sui nodi per l'esterno
    protocol: TCP
\end{lstlisting}

\subsection{L'Oggetto Di Servizio LoadBalancer}
Sebbene il modello NodePort consenta l'accesso dall'esterno, trasferisce sul client l'onere di conoscere gli indirizzi IP dei singoli nodi e di gestire gli eventuali fallimenti hardware. Per ovviare a ciò, negli ambienti cloud pubblici o in infrastrutture con controller integrati viene impiegato il servizio di tipo \textbf{LoadBalancer}.

\subsubsection{Integrazione Cloud-Native}
Il servizio LoadBalancer è una sovrastruttura che include tutte le proprietà del NodePort e del ClusterIP:
\begin{enumerate}[noitemsep]
    \item Kubernetes alloca internamente un ClusterIP.
    \item Kubernetes alloca una porta NodePort su tutti i nodi fisici.
    \item Il componente **Cloud Controller Manager (CCM)** intercetta la risorsa e richiama, mediante chiamate API asincrone, l'infrastruttura del cloud provider (\textbf{AWS ELB/NLB}, \textbf{GCP Cloud Load Balancing}, \textbf{Azure Load Balancer}).
    \item Il cloud provider alloca un bilanciatore di carico hardware/software esterno, gli assegna un indirizzo IP pubblico instradabile su Internet e configura i propri target pool per inoltrare il traffico verso le porte \texttt{NodePort} dei nodi worker.
\end{enumerate}

\begin{lstlisting}[language=yaml, caption={Definizione di un servizio LoadBalancer (\texttt{loadbalancer1.yml})}]
apiVersion: v1
kind: Service
metadata:
  name: service-test
spec:
  type: LoadBalancer
  selector:
    app: service_test_pod
  ports:
  - port: 80
    targetPort: http
\end{lstlisting}

\subsubsection{Ispezione E Compromessi Architetturali}
Una volta completato il provisioning asincrono, il comando \texttt{kubectl get svc} mostrerà l'IP esterno:
\begin{lstlisting}[language=bash]
$ kubectl get svc service-test
NAME           TYPE           CLUSTER-IP     EXTERNAL-IP     PORT(S)        AGE
service-test   LoadBalancer   10.233.47.52   192.168.1.240   80:30557/TCP   12s
\end{lstlisting}

Nonostante l'automazione trasparente, l'adozione incontrollata di servizi di tipo LoadBalancer solleva problematiche economiche e gestionali:
\begin{itemize}[noitemsep]
    \item \textbf{Costo}: Ciascun LoadBalancer allocato presso i cloud provider prevede un canone orario/mensile fisso oltre al costo per gigabyte scambiato. Un'architettura a microservizi con centinaia di servizi comporterebbe spese ingenti.
    \item \textbf{Spreco di indirizzi IPv4}: Ogni servizio consuma un IP pubblico dedicato.
    \item \textbf{Mancanza di intelligenza di Livello 7}: I LoadBalancer infrastrutturali operano prettamente a Livello 4 (TCP/UDP), risultando incapaci di instradare le richieste in base al path URL (\emph{path-based routing}) o all'hostname (\emph{host-based routing}).
\end{itemize}

\newpage

% =============================================================================
% SEZIONE: SERVIZI SENZA SELETTORE ED EXTERNALNAME
% =============================================================================

\section{Servizi Senza Selettore Ed ExternalName}

Sebbene l'impiego canonico dei Service preveda l'aggancio dinamico ai Pod mediante il costrutto \texttt{selector}, esistono scenari architetturali in cui il backend non è costituito da Pod interni al cluster.

\subsection{Servizi Senza Label Selector}
Se si definisce un Service omettendo deliberatamente il blocco \texttt{spec.selector}, il control plane di Kubernetes non genererà automaticamente l'oggetto \textbf{Endpoints} (o \textbf{EndpointSlice}). Spetta all'amministratore definire manualmente tale risorsa, associandovi indirizzi IP e porte arbitrari.

% [IMMAGINE DEL FILE DI SLIDE CON NOME "Calico-Kubernetes-Services.pdf"] Schema architetturale di un Service senza selettore che punta a risorse esterne al cluster — Slide 137 / Pagina 16 del PDF

I contesti applicativi tipici includono:
\begin{itemize}[noitemsep]
    \item Connessione a cluster di database esterni (es. Oracle RAC, PostgreSQL bare-metal) eseguiti al di fuori dell'infrastruttura container per motivi prestazionali o di licenza.
    \item Indirizzamento verso servizi residenti in un differente Namespace o in un altro cluster Kubernetes.
    \item Fasi transitorie di migrazione dei carichi da architetture monolitiche legacy verso il cloud.
\end{itemize}

\subsection{Il Tipo Di Servizio ExternalName}
Quando l'integrazione con l'esterno deve avvenire a livello di nome DNS piuttosto che tramite indirizzi IP statici, Kubernetes fornisce il tipo \textbf{ExternalName}.

A differenza di tutte le altre tipologie di servizio:
\begin{itemize}[noitemsep]
    \item Non viene allocato alcun indirizzo ClusterIP virtuale.
    \item Non vengono programmate regole Netfilter o iptables sul kernel dei nodi.
    \item \textbf{Kube-proxy viene interamente scavalcato} (\emph{bypass}).
    \item La risoluzione opera esclusivamente a livello di server DNS interno (**CoreDNS**).
\end{itemize}

\begin{lstlisting}[language=yaml, caption={Definizione di un servizio ExternalName (\texttt{external1.yml})}]
apiVersion: v1
kind: Service
metadata:
  name: my-database-svc
  namespace: prod
spec:
  type: ExternalName
  externalName: my.database.example.com
  ports:
  - port: 3306
    targetPort: 3306  # Puramente informativo, non influenza kube-proxy
\end{lstlisting}

\subsubsection{Meccanismo Di Risoluzione DNS CNAME}
Quando un'applicazione interna al namespace \texttt{prod} tenta di connettersi al database utilizzando il nome canonico locale:
\[
    \texttt{my-database-svc.prod.svc.cluster.local}
\]
CoreDNS intercetta la richiesta e non restituisce un record di tipo A (indirizzo IP), bensì un record di tipo \textbf{CNAME} (\emph{Canonical Name}) che mappa direttamente il nome sull'indirizzo esterno configurato:
\[
    \texttt{my-database-svc.prod.svc.cluster.local.} \quad \text{CNAME} \quad \texttt{my.database.example.com.}
\]
Il resolver del Pod client riceve il record CNAME ed esegue una risoluzione standard verso l'IP pubblico del database esterno, instradando il traffico TCP \textbf{direttamente} dal Pod all'endpoint remoto sulla porta \texttt{3306}, senza alcuna mediazione o overhead da parte dei nodi del cluster.

\newpage

% =============================================================================
% SEZIONE: INGRESS CONTROLLER E ROUTING L7
% =============================================================================

\section{Routing Applicativo Di Livello 7 Tramite La Risorsa Ingress}

I limiti strutturali ed economici associati all'esposizione di molteplici servizi tramite LoadBalancer di Livello 4 hanno guidato la progettazione della risorsa **Ingress**. Ingress agisce come un punto di ingresso unico centralizzato al cluster, operando al **Livello 7 (Applicazione - HTTP/HTTPS)** dello stack OSI.

% [IMMAGINE DEL FILE DI SLIDE CON NOME "Calico-Kubernetes-Services.pdf"] Schema del routing applicativo Layer 7 tramite Ingress verso molteplici servizi interni — Slide 139 / Pagina 18 del PDF

\subsection{Funzionalità Cardine Dell'Ingress}
La specifica Ingress fornisce nativamente:
\begin{itemize}[noitemsep]
    \item \textbf{Routing basato su Host (\emph{Name-based Virtual Hosting})}: Consente di instradare il traffico proveniente da molteplici nomi di dominio (es. \texttt{api.example.com} e \texttt{web.example.com}) condividendo un unico indirizzo IP pubblico. Il componente esamina l'header HTTP \texttt{Host} inviato dal client nel layer applicativo.
    \item \textbf{Routing basato su Path (\emph{Path-based routing})}: Instrada le richieste a servizi differenti in base al prefisso dell'URI invocato (es. \texttt{/api} inoltrato al servizio di backend, \texttt{/static} inoltrato al frontend web).
    \item \textbf{Terminazione crittografica SSL/TLS (\emph{TLS Termination})}: Il traffico HTTPS cifrato viene terminato e decifrato a monte dall'Ingress, inoltrando richieste HTTP in chiaro all'interno della rete del cluster, sgravando i Pod applicativi dal pesante overhead computazionale degli handshake crittografici.
\end{itemize}

\subsection{Disaccoppiamento Tra Ingress Resource Ed Ingress Controller}
Un principio fondante dell'architettura di Kubernetes è la netta separazione tra la definizione dichiarativa delle regole e la loro effettiva esecuzione:
\begin{itemize}[noitemsep]
    \item \textbf{Ingress Resource}: È il documento manifest YAML (\texttt{kind: Ingress}) in cui lo sviluppatore descrive l'insieme astratto delle regole di routing, i domini e i percorsi associati. La semplice creazione di un oggetto Ingress nell'API server \textbf{non produce alcun effetto operativo sulla rete}.
    \item \textbf{Ingress Controller}: È l'applicazione attiva (generalmente un reverse proxy avanzato come \textbf{NGINX}, \textbf{Traefik}, \textbf{HAProxy} o controller integrati cloud come \textbf{GCE Ingress Controller}) che deve essere preventivamente installata nel cluster. L'Ingress Controller esegue un loop di riconciliazione: monitora le Ingress Resource create, ne interpreta le specifiche e riconfigura dinamicamente il proprio file di configurazione interno (es. \texttt{nginx.conf}), ricaricando i processi a caldo.
\end{itemize}

\subsection{Pattern Di Esposizione Dell'Ingress Controller}
Affinché l'Ingress Controller possa ricevere il traffico proveniente da Internet, il Pod che esegue il proxy deve essere a sua volta esposto. Esistono tre pattern architetturali per raggiungere tale obiettivo:

% [IMMAGINE DEL FILE DI SLIDE CON NOME "Calico-Kubernetes-Services.pdf"] Confronto tra i metodi di esposizione dell'Ingress Controller: LoadBalancer vs NodePort vs hostNetwork — Slide 141-144 / Pagine 21-24 del PDF

\subsubsection{Pattern 1: Esposizione Tramite LoadBalancer Service (Ambienti Cloud)}
In infrastrutture cloud gestite (AWS, GCP, Azure), l'Ingress Controller viene pubblicato tramite un singolo servizio di tipo \texttt{LoadBalancer}:

\begin{lstlisting}[language=yaml, caption={Ingress Controller esposto via LoadBalancer}]
apiVersion: v1
kind: Service
metadata:
  name: ingress-nginx-controller
  namespace: ingress-nginx
spec:
  type: LoadBalancer
  ports:
  - name: http
    port: 80
    targetPort: 80
  - name: https
    port: 443
    targetPort: 443
  selector:
    app.kubernetes.io/name: ingress-nginx
\end{lstlisting}

\begin{itemize}[noitemsep]
    \item \textbf{Vantaggi}: Automazione completa, rilascio di un IP pubblico stabile, scalabilità elastica e monitoraggio di integrità (\emph{health check}) gestito direttamente dall'infrastruttura cloud.
    \item \textbf{Svantaggi}: Canone economico mensile per il bilanciatore cloud, rischio di \emph{vendor lock-in} e tempi di provisioning dell'ordine di alcuni minuti.
\end{itemize}

\subsubsection{Pattern 2: Esposizione Tramite NodePort Service (Ambienti On-Premise)}
In cluster fisici o datacenter privati privi di API di integrazione cloud, il controller viene associato a un servizio NodePort:

\begin{lstlisting}[language=yaml, caption={Ingress Controller esposto via NodePort}]
apiVersion: v1
kind: Service
metadata:
  name: ingress-nginx-controller
  namespace: ingress-nginx
spec:
  type: NodePort
  ports:
  - name: http
    port: 80
    targetPort: 80
    nodePort: 30080
  - name: https
    port: 443
    targetPort: 443
    nodePort: 30443
  selector:
    app.kubernetes.io/name: ingress-nginx
\end{lstlisting}

\begin{itemize}[noitemsep]
    \item \textbf{Vantaggi}: Nessun costo aggiuntivo di licenza o cloud, massima flessibilità e compatibilità con qualsiasi ambiente Linux.
    \item \textbf{Svantaggi}: Richiede l'apertura di porte non standard (\texttt{30080}, \texttt{30443}). In scenari di produzione, impone la configurazione e manutenzione manuale di un proxy esterno addizionale (es. F5 BIG-IP, HAProxy fisico) per inoltrare il traffico standard 80/443 verso i nodi.
\end{itemize}

\subsubsection{Pattern 3: Esposizione Diretta Tramite \texttt{hostNetwork: true}}
Per eliminare completamente l'overhead di rete, è possibile eseguire il Pod dell'Ingress Controller legandolo direttamente al namespace di rete dell'host mediante la direttiva \texttt{hostNetwork: true}:

\begin{lstlisting}[language=yaml, caption={Deployment Ingress Controller in modalità \texttt{hostNetwork}}]
apiVersion: apps/v1
kind: Deployment
metadata:
  name: ingress-nginx-controller
  namespace: ingress-nginx
spec:
  template:
    spec:
      hostNetwork: true
      dnsPolicy: ClusterFirstWithHostNet
      containers:
      - name: controller
        image: registry.k8s.io/ingress-nginx/controller:v1.8.0
        ports:
        - name: http
          containerPort: 80
          hostPort: 80
        - name: https
          containerPort: 443
          hostPort: 443
\end{lstlisting}

\begin{itemize}[noitemsep]
    \item \textbf{Vantaggi}: Massime prestazioni possibili; eliminazione totale dei layer intermedi di kube-proxy, coppie \emph{veth} e traduzioni NAT; impegno nativo delle porte privilegiate \texttt{80} e \texttt{443} direttamente sulle schede di rete fisiche dei nodi selezionati.
    \item \textbf{Svantaggi}: \textbf{Conflitto di porte}: su ciascun nodo può essere eseguito un solo Pod di Ingress Controller, poiché la porta fisica 80 non può essere condivisa. Comporta inoltre seri rischi di sicurezza, in quanto il Pod ottiene piena visibilità sul traffico dell'host, e richiede l'impiego rigido di \texttt{nodeSelector} o \texttt{tolerations} per confinare l'esecuzione su nodi dedicati (\emph{edge nodes}).
\end{itemize}

\subsection{Sintassi E Matching Delle Regole Ingress}
La definizione di una risorsa Ingress complessa (conosciuta come \emph{Multi-Service Ingress}) struttura le regole analizzando nome host, path e tipologia di matching:

\begin{lstlisting}[language=yaml, caption={Manifest di Ingress con routing multi-servizio e annotazioni}]
apiVersion: networking.k8s.io/v1
kind: Ingress
metadata:
  name: simple-fanout-example
  annotations:
    nginx.ingress.kubernetes.io/rewrite-target: /
spec:
  rules:
  - host: foo.bar.com
    http:
      paths:
      - path: /foo
        pathType: Prefix
        backend:
          service:
            name: service1
            port:
              number: 4200
      - path: /bar
        pathType: Exact
        backend:
          service:
            name: service2
            port:
              number: 8080
\end{lstlisting}

\subsubsection{Tipologie Di Matching Del Path (\texttt{pathType})}
Lo standard Kubernetes prescrive tre modalità rigorose per l'interpretazione dei percorsi:
\begin{enumerate}[noitemsep]
    \item \textbf{\texttt{Exact}}: Il percorso della richiesta HTTP deve coincidere esattamente con il valore specificato, rispettando la distinzione tra maiuscole e minuscole (\emph{case-sensitive}). Ad esempio, \texttt{/bar} accetta solo \texttt{/bar} e rifiuta \texttt{/bar/}.
    \item \textbf{\texttt{Prefix}}: Il matching avviene elemento per elemento, segmentando l'URI tramite il carattere delimitatore \texttt{/}. Ad esempio, la regola \texttt{/foo} corrisponde con successo a \texttt{/foo}, \texttt{/foo/} e \texttt{/foo/schede}, ma rifiuta \texttt{/foobar}.
    \item \textbf{\texttt{ImplementationSpecific}}: L'interpretazione dei caratteri jolly o delle espressioni regolari (regex) è delegata alla logica interna dello specifico Ingress Controller in uso (definito dalla classe \texttt{ingressClassName}).
\end{enumerate}

\subsubsection{Vincoli Sui Nomi Di Dominio (RFC 1123)}
I campi \texttt{host} devono rispettare i criteri formali dei sottodomini DNS stabiliti dallo standard \textbf{RFC 1123}:
\begin{itemize}[noitemsep]
    \item Lunghezza massima non superiore a 253 caratteri.
    \item Composizione limitata esclusivamente a caratteri alfanumerici minuscoli, trattini (\texttt{-}) e punti (\texttt{.}).
    \item Inizio e fine della stringa vincolati a caratteri alfanumerici (non possono iniziare o finire con un trattino o un punto).
    \item Supporto per corrispondenze esatte (es. \texttt{foo.bar.com}) o wildcard con prefisso (es. \texttt{*.bar.com}).
\end{itemize}

\subsection{Sicurezza E Terminazione TLS Mediante Oggetti Secret}
Per implementare la cifratura HTTPS, l'Ingress si integra con la risorsa **Secret** di tipo \texttt{kubernetes.io/tls}.

% [IMMAGINE DEL FILE DI SLIDE CON NOME "Calico-Kubernetes-Services.pdf"] Diagramma dell'integrazione tra Ingress e risorsa Secret per la terminazione TLS — Slide 147 / Pagina 29 del PDF

Il Secret custodisce in formato codificato \textbf{Base64} il certificato digitale X.509 (\texttt{tls.crt}) e la corrispondente chiave privata RSA/ECC (\texttt{tls.key}):

\begin{lstlisting}[language=yaml, caption={Manifest della risorsa Secret di tipo TLS}]
apiVersion: v1
kind: Secret
metadata:
  name: testsecret-tls
  namespace: default
type: kubernetes.io/tls
data:
  tls.crt: LS0tLS1CRUdJTiBDRVJUSUZJQ0FURS0tLS0tC... (base64)
  tls.key: LS0tLS1CRUdJTiBSU0EgUFJJVkFURSBLRVktLS... (base64)
\end{lstlisting}

L'Ingress richiama la risorsa mediante il blocco \texttt{spec.tls}:

\begin{lstlisting}[language=yaml, caption={Configurazione TLS nell'Ingress}]
spec:
  tls:
  - hosts:
    - sslexample.foo.com
    secretName: testsecret-tls
  rules:
  - host: sslexample.foo.com
    http:
      paths:
      - path: /
        pathType: Prefix
        backend:
          service:
            name: service1
            port:
              number: 80
\end{lstlisting}

\subsubsection{Diagnostica Dei Problemi Comuni Di Configurazione TLS}
Durante la messa in esercizio di endpoint HTTPS protetti da Ingress, gli errori più frequenti includono:
\begin{itemize}[noitemsep]
    \item \textbf{Disallineamento dei Namespace}: L'oggetto \texttt{Secret} contenente i certificati risiede in un namespace differente rispetto a quello in cui è allocata la risorsa \texttt{Ingress}. Per vincolo architetturale, l'Ingress può referenziare esclusivamente Secret collocati nel suo medesimo namespace.
    \item \textbf{Mancata corrispondenza del CN/SAN}: Il Common Name o il \emph{Subject Alternative Name} registrato all'interno del certificato X.509 non coincide con l'hostname specificato nel campo \texttt{spec.rules[].host}, scatenando errori di certificato non valido (\texttt{SSL\_ERROR\_BAD\_CERT\_DOMAIN}) nel browser.
    \item \textbf{Errori di codifica Base64}: Caratteri spuri, inclusioni errate di newline o troncamenti durante la conversione da file PEM a Base64 rendono la chiave privata illeggibile per il parser OpenSSL dell'Ingress Controller.
\end{itemize}

\newpage

% =============================================================================
% SEZIONE: GATEWAY API
% =============================================================================

\section{La Nuova Generazione Del Networking: Le Gateway API}

Sebbene ampiamente adottata, la specifica Ingress manifesta limiti strutturali rilevanti nei cluster enterprise:
\begin{itemize}[noitemsep]
    \item \textbf{Mancanza di standardizzazione}: L'assenza di campi avanzati ha indotto i produttori a proliferare annotazioni proprietarie disomogenee (\texttt{annotations}) per supportare funzionalità quali il riscrittura dei path (\emph{URL rewriting}), il rate limiting o il bilanciamento pesato.
    \item \textbf{Monolitismo dei ruoli}: Un unico file manifest unisce la configurazione infrastrutturale (indirizzi IP, porte, certificati TLS) e le regole di routing applicativo dello sviluppatore, violando la separazione delle responsabilità (\emph{Role-Based Access Control - RBAC}).
\end{itemize}

Per superare tali limitazioni, la comunità di Kubernetes ha introdotto la specifica **Gateway API** (evoluzione formale dello standard Ingress), imperniata su un'architettura modulare a tre livelli gerarchici.

% [IMMAGINE DEL FILE DI SLIDE CON NOME "Calico-Kubernetes-Services.pdf"] Schema architetturale della Gateway API con separazione dei ruoli tra Gateway, HTTPRoute e Service — Slide 148 / Pagina 32 del PDF

\subsection{Architettura E Separazione Dei Ruoli Nelle Gateway API}
Le Gateway API disaccoppiano le responsabilità su tre risorse distinte:
\begin{enumerate}
    \item \textbf{\texttt{GatewayClass}} (Gestita dal \emph{Cluster Operator}): Definisce l'implementazione del controller di rete sottostante (es. NGINX, Traefik, Istio, Envoy).
    \item \textbf{\texttt{Gateway}} (Gestita dal \emph{Platform / Infrastructure Engineer}): Dichiara il punto di ingresso dell'infrastruttura di rete, definendo gli indirizzi IP fisici, le porte esposte e i listener di protocollo:
\begin{lstlisting}[language=yaml, caption={Definizione del punto di ingresso Gateway}]
apiVersion: gateway.networking.k8s.io/v1
kind: Gateway
metadata:
  name: my-gateway
spec:
  gatewayClassName: nginx
  listeners:
  - name: http
    protocol: HTTP
    port: 80
\end{lstlisting}
    \item \textbf{\texttt{HTTPRoute}} (Gestita dall'\emph{Application Developer}): Specifica in maniera atomica e standardizzata come il traffico debba essere instradato verso i backend applicativi, referenziando il genitore infrastrutturale mediante il blocco \texttt{parentRefs}:
\begin{lstlisting}[language=yaml, caption={Definizione delle regole di routing con HTTPRoute}]
apiVersion: gateway.networking.k8s.io/v1
kind: HTTPRoute
metadata:
  name: web-route
spec:
  parentRefs:
  - name: my-gateway
  rules:
  - matches:
    - path:
        type: PathPrefix
        value: /
    backendRefs:
    - name: web
      port: 80
\end{lstlisting}
\end{enumerate}

Grazie a questo disaccoppiamento, gli ingegneri infrastrutturali possono vincolare e proteggere le porte fisiche e i certificati crittografici, mentre i team applicativi mantengono piena autonomia nel distribuire, modificare e aggiornare le proprie regole di routing senza rischiare interferenze cross-tenant.

\newpage

% =============================================================================
% SEZIONE: HORIZONTAL POD AUTOSCALER (HPA)
% =============================================================================

\section{Scalabilità Dinamica: Horizontal Pod Autoscaler (HPA)}

Nelle architetture cloud ad alto carico dinamico, l'allocazione statica delle repliche applicative è insufficiente: sottodimensionare il carico conduce al degrado del servizio (\emph{starvation} delle risorse e incremento delle code di attesa), mentre sovradimensionarlo comporta un ingiustificato dispendio economico. Per adeguare elasticamente le risorse al fabbisogno computazionale, Kubernetes integra il controller **Horizontal Pod Autoscaler (HPA)**.

\subsection{Scaling Orizzontale vs. Scaling Verticale}
È cruciale chiarire la divergenza teorica tra i due paradigmi di scalabilità:
\begin{itemize}[noitemsep]
    \item \textbf{Vertical Scaling (\emph{Scaling Up/Down})}: Incrementa o decrementa la quota di risorse computazionali (core CPU o memoria RAM) allocate ai singoli container preesistenti. In Kubernetes, questa procedura (gestita dal \emph{Vertical Pod Autoscaler - VPA}) è complessa e storicamente imponeva il riavvio del container per aggiornare i limiti imposti nei cgroups del kernel Linux.
    \item \textbf{Horizontal Scaling (\emph{Scaling Out/In})}: Varia dinamicamente il \textbf{numero totale di repliche (Pod)} attive per uno specifico carico di lavoro, mantenendo invariata l'allocazione di risorse dei singoli individui. È il modello ideale per carichi applicativi \emph{stateless}, distribuendo il traffico in arrivo su un fronte più ampio.
\end{itemize}

\begin{theorem}
L'autoscaling orizzontale può essere applicato esclusivamente a risorse replicabili scalabili dichiarativamente tramite un campo \texttt{spec.replicas} (quali \textbf{Deployment}, \textbf{ReplicaSet} o \textbf{StatefulSet}). Esso \textbf{non può} essere applicato ad oggetti architetturali vincolati alla topologia hardware del cluster, come i \textbf{DaemonSet} (la cui esecuzione prevede tassativamente un solo Pod per nodo).
\end{theorem}

\subsection{Il Ciclo Di Controllo Chiuso Dell'HPA}
L'HPA è implementato come un ciclo di controllo continuo a retroazione (\emph{feedback loop}) eseguito all'interno del **Kubernetes Controller Manager**.

% [IMMAGINE DEL FILE DI SLIDE CON NOME "Calico-Kubernetes-Services.pdf"] Schema del ciclo di retroazione di HPA: Metrics Server, Controller, Replica Calculation e Scaling — Slide 149 / Pagina 33 del PDF

L'algoritmo si articola in cinque fasi sequenziali iterate periodicamente (con frequenza predefinita tipica di 15 secondi):
\begin{enumerate}
    \item \textbf{Campionamento delle metriche}: Il controller HPA interroga l'infrastruttura di aggregazione delle metriche, primariamente il **Metrics Server**, che raccoglie l'utilizzo istantaneo di CPU e memoria aggregando i dati emessi dal sottosistema \textbf{cAdvisor} integrato nei Kubelet di ciascun nodo.
    \item \textbf{Elaborazione dello scostamento}: Il controllore calcola la media aritmetica dell'utilizzo della metrica tra tutti i Pod che concorrono al carico di lavoro e la confronta con il valore di soglia desiderato (\emph{target utilization}).
    \item \textbf{Computazione del numero di repliche desiderate}: Applicando una formula proporzionale a ciclo chiuso, l'HPA stima il numero ottimale di repliche necessarie a ricondurre l'utilizzo della metrica al valore obiettivo.
    \item \textbf{Attuazione dello scaling}: Qualora il valore computato differisca dal numero attuale di repliche (oltre una soglia di tolleranza predefinita, tipicamente il $10\%$, introdotta per prevenire oscillazioni repentine note come \emph{flapping} o \emph{thrashing}), il controller aggiorna il campo \texttt{spec.replicas} del Deployment.
    \item \textbf{Riconciliazione del carico}: Il controller del Deployment rileva la modifica e istruisce il ReplicaSet a creare o distruggere le istanze di Pod necessarie. L'HPA attende un tempo di stabilizzazione (\emph{cool-down period}) e ripete il ciclo.
\end{enumerate}

\subsection{Formulazione Matematica Dell'Algoritmo Di Autoscaling}
Il dimensionamento attuato dal controller risponde a una legge analitica proporzionale rigorosa.

Sia:
\begin{itemize}[noitemsep]
    \item $R_{\text{corrente}}$ il numero di repliche attualmente in esecuzione;
    \item $M_{\text{corrente}}$ il valore medio misurato della metrica di utilizzo (es. percentuale di utilizzo CPU rispetto al \texttt{request} dichiarato);
    \item $M_{\text{target}}$ il valore obiettivo desiderato configurato nella specifica dell'HPA.
\end{itemize}

Il numero di repliche obiettivo $R_{\text{desiderato}}$ è computato secondo l'operatore di approssimazione per eccesso (\emph{ceiling}):
\[
    \boxed{R_{\text{desiderato}} = \left\lceil R_{\text{corrente}} \times \left( \frac{M_{\text{corrente}}}{M_{\text{target}}} \right) \right\rceil}
\]

\subsubsection{Dimostrazione Ingegneristica Di Un Ciclo Di Calcolo}
Ipotizziamo uno scenario operativo in cui un Deployment di un microservizio web sia attualmente composto da:
\[
    R_{\text{corrente}} = 3 \text{ repliche}
\]
Ciascun container ha una richiesta di CPU dichiarata (\texttt{cpu request}) pari a $200\,\text{mCPU}$. L'amministratore impone una soglia target di utilizzo:
\[
    M_{\text{target}} = 50\% \quad (\text{ovvero } 100\,\text{mCPU per replica})
\]
A seguito di un picco imprevisto di traffico HTTP in ingresso, il Metrics Server campiona un consumo istantaneo complessivo pari a $600\,\text{mCPU}$ distribuito sulle tre istanze, rilevando un valore medio misurato pari a:
\[
    M_{\text{corrente}} = \frac{600\,\text{mCPU}}{3 \times 200\,\text{mCPU}} = 100\%
\]

Applicando l'equazione di bilancio dell'HPA:
\begin{align*}
    R_{\text{desiderato}} &= \left\lceil 3 \times \left( \frac{100\%}{50\%} \right) \right\rceil \\
    &= \left\lceil 3 \times 2.0 \right\rceil \\
    &= \mathbf{6 \text{ repliche}}
\end{align*}

Il controller HPA aggiorna il Deployment imponendo l'attivazione immediata di 3 nuovi Pod, distribuendo il carico complessivo su 6 repliche e riportando il consumo medio pro-capite all'esatto valore di soglia stazionario del $50\%$. Qualora il traffico dovesse successivamente estinguersi con una caduta del consumo medio al $25\%$, l'applicazione della formula determinerebbe:
\[
    R_{\text{desiderato}} = \left\lceil 6 \times \left( \frac{25\%}{50\%} \right) \right\rceil = \lceil 3.0 \rceil = \mathbf{3 \text{ repliche}}
\]
ripristinando la conformazione minima iniziale nel pieno rispetto dei vincoli di efficienza e ottimizzazione delle risorse computazionali del datacenter.


% AA Orchestrazione_1 [199-230] ------------------------------------------------------------

\section{Scalabilità Dinamica Avanzata: Horizontal Pod Autoscaler (HPA)}

Nelle architetture cloud-native distribuite, la capacità di reagire elasticamente alle fluttuazioni del carico di lavoro rappresenta un requisito essenziale per soddisfare gli accordi sul livello di servizio (\emph{Service Level Agreements}, SLA) minimizzando contestualmente i costi operativi infrastrutturali. In Kubernetes, tale funzionalità è demandata all'**Horizontal Pod Autoscaler (HPA)**, un controller a retroazione che adegua dinamicamente il numero di repliche dei carichi di lavoro applicativi.

\subsection{Richiamo Concettuale: Cicli Di Controllo Nei Sistemi Distribuiti}
Un sistema di controllo a ciclo chiuso (\emph{closed-loop control system}) misura con regolarità lo stato osservabile di un processo, lo confronta con un valore di riferimento desiderato (\emph{setpoint}) e calcola un segnale di correzione atto ad annullare l'errore.
In Kubernetes, questo principio è implementato dal **control loop** dell'HPA:
\begin{itemize}[noitemsep]
    \item \textbf{Stato osservato}: L'utilizzo medio delle risorse computazionali (ad esempio, milli-core di CPU o byte di memoria RAM) oppure metriche applicative custom (come il numero di richieste HTTP al secondo o la profondità di una coda di messaggi).
    \item \textbf{Valore di riferimento}: La soglia target stabilita dichiarativamente nel manifest dell'HPA (es. il $60\%$ della richiesta di CPU dichiarata o un valore assoluto prefissato).
    \item \textbf{Azione di controllo}: La modifica del campo \texttt{spec.replicas} all'interno della risorsa target (\textbf{Deployment}, \textbf{ReplicaSet} o \textbf{StatefulSet}).
\end{itemize}

A differenza dello scaling verticale (\emph{Vertical Pod Autoscaler}, VPA), che altera i vincoli fisici di risorse (cgroups) assegnati ai singoli container imponendone storicamente il riavvio, lo scaling orizzontale preserva l'immutabilità dei Pod esistenti, istanziandone o terminandone di nuovi. L'HPA è pertanto ideale per microservizi \emph{stateless}. Esso \textbf{non è applicabile} a risorse vincolate alla topologia dei nodi fisici, come i \textbf{DaemonSet}.

\subsection{Architettura Della Resource Metrics Pipeline}
L'HPA non monitora direttamente i processi nei container: per assumere decisioni di dimensionamento, esso dipende da una catena architetturale gerarchica denominata **Resource Metrics Pipeline**, la quale campiona, normalizza ed espone i dati di telemetria.

% [IMMAGINE DEL FILE DI SLIDE CON NOME "Kubernetes-HPA-RBAC-Deployment.pdf"] Diagramma architetturale della Resource Metrics Pipeline da destra verso sinistra: cAdvisor, Kubelet Summary API, Metrics Server e API Server — Slide 151 / Pagina 3 del PDF

L'architettura del flusso metrico si sviluppa attraverso i seguenti componenti:
\begin{enumerate}
    \item \textbf{cAdvisor (\emph{Container Advisor})}: È un demone di telemetria a basso livello integrato nativamente all'interno del binario del **Kubelet** su ciascun nodo worker. cAdvisor interroga direttamente le interfacce del kernel Linux (specificamente i sottosistemi di contabilità dei \textbf{cgroups}, quali \texttt{cpuacct} e \texttt{memory}) per raccogliere metriche storiche e istantanee sull'utilizzo di CPU, memoria, filesystem e interfacce di rete per ogni container attivo.
    \item \textbf{Kubelet (Summary API)}: L'agente di nodo raccoglie i campioni elaborati da cAdvisor e li rende disponibili all'interno del cluster attraverso endpoint HTTP interni protetti:
    \begin{itemize}[noitemsep]
        \item \texttt{/metrics/resource}: Espone le metriche ``core'' (CPU e memoria) formattate per il consumo standardizzato.
        \item \texttt{/stats/summary}: Espone la **Summary API**, una vista aggregata e dettagliata dello stato del nodo e dei relativi Pod in formato JSON.
    \end{itemize}
    \item \textbf{Metrics Server}: È un componente di cluster (eseguito tipicamente come Deployment singolo nel namespace \texttt{kube-system}) che funge da aggregatore in-memory ad alte prestazioni. A intervalli periodici, il Metrics Server interroga l'endpoint \texttt{/metrics/resource} di ciascun Kubelet presente nel cluster, normalizza le misurazioni e le mantiene volatili in memoria (non implementa un database storico permanente).
    \item \textbf{API Server (Metrics API)}: Il Metrics Server registra se stesso come un **API Service** aggregato (\texttt{metrics.k8s.io}). L'API Server centrale di Kubernetes funge da intermediario unico: riceve le chiamate esterne o interne, valida l'autenticazione/autorizzazione tramite RBAC e reindirizza le richieste alla Metrics API verso il Metrics Server.
\end{enumerate}

\subsubsection{Estensione Alle Metriche Avanzate E Custom Metrics}
Mentre il Metrics Server gestisce unicamente le metriche ``core'' essenziali per il funzionamento di base dell'HPA e del comando \texttt{kubectl top}, la telemetria complessa (metriche di storage I/O, latenza di rete, \emph{Pressure Stall Information} - PSI o metriche applicative L7) richiede strumenti dedicati:
\begin{itemize}[noitemsep]
    \item \textbf{Prometheus}: Tramite appositi moduli di esportazione (come il \emph{Node Exporter} o sonde dirette sui Pod), estrae i dati dalla Summary API (accessibile localmente aprendo un canale sicuro con \texttt{kubectl proxy} all'indirizzo \texttt{http://localhost:8080/api/v1/nodes/\{node-name\}/proxy/stats/summary}).
    \item \textbf{Custom Metrics API / External Metrics API}: Consentono all'HPA di ridimensionare i carichi di lavoro basandosi su parametri esterni al runtime (es. metriche esposte da ingress controller, code RabbitMQ o bilanciatori cloud).
\end{itemize}

\subsection{Installazione E Configurazione Operativa Del Metrics Server}
Negli ambienti di laboratorio o su cluster bare-metal inizializzati tramite strumenti a basso livello, il Metrics Server non è installato di default. La sua distribuzione avviene applicando il manifest ufficiale del repository di progetto:

\begin{lstlisting}[language=bash]
$ kubectl apply -f https://github.com/kubernetes-sigs/metrics-server/releases/latest/download/components.yaml
\end{lstlisting}

\subsubsection{Risoluzione Delle Problematiche TLS Nei Certificati Di Nodo}
Nei cluster di sviluppo, i Kubelet generano frequentemente certificati di servizio auto-firmati privi di una catena di certificazione verificabile da una CA pubblica riconosciuta. In tale scenario, il Metrics Server fallisce la connessione HTTPS verso i Kubelet sulla porta sicura \texttt{10250}. Per consentire l'acquisizione delle metriche, è necessario intervenire sul manifest del container aggiungendo direttive specifiche nel blocco \texttt{args}:

\begin{lstlisting}[language=yaml, caption={Estratto di configurazione per Metrics Server (\texttt{components.yaml})}]
template:
  metadata:
    labels:
      k8s-app: metrics-server
  spec:
    containers:
    - name: metrics-server
      image: k8s.gcr.io/metrics-server/metrics-server:v0.4.2
      args:
      - --cert-dir=/tmp
      - --secure-port=10250
      - --kubelet-insecure-tls  # Disabilita la verifica della CA per certificati Kubelet
      - --kubelet-preferred-address-types=InternalIP,ExternalIP,Hostname
      - --kubelet-use-node-status-port
\end{lstlisting}

L'argomento \texttt{--kubelet-preferred-address-types} impone l'ordine di priorità con cui il Metrics Server risolve gli indirizzi dei nodi, privilegiando l'interfaccia interna (\texttt{InternalIP}) per evitare problemi di risoluzione DNS sui nomi host.

\subsection{Algoritmo E Formulazione Matematica Dell'HPA}
L'HPA opera in maniera intermittente. L'intervallo temporale tra due valutazioni successive è governato dal parametro \texttt{--horizontal-pod-autoscaler-sync-period} configurato sul componente \textbf{kube-controller-manager} (il cui valore predefinito è pari a \textbf{15 secondi}).

A ogni ciclo di sincronizzazione, il controller:
\begin{enumerate}[noitemsep]
    \item Identifica la risorsa target dichiarata nel campo \texttt{scaleTargetRef}.
    \item Seleziona tutti i Pod appartenenti al carico di lavoro tramite le etichette (\texttt{.spec.selector}).
    \item Interroga la Metrics API per ottenere l'utilizzo istantaneo di ciascun Pod target.
    \item Calcola la media aritmetica dell'utilizzo tra tutti i Pod considerati sani e pronti (\emph{Ready}).
    \item Calcola il nuovo numero di repliche necessarie mediante l'algoritmo di proporzionalità diretta.
\end{enumerate}

\subsubsection{Formulazione Analitica}
Siano:
\begin{itemize}[noitemsep]
    \item $R_{\text{curr}}$ il numero di repliche attualmente attive;
    \item $M_{\text{curr}}$ il valore medio misurato della metrica (espresso in valore assoluto o percentuale);
    \item $M_{\text{target}}$ il valore target desiderato configurato nell'HPA;
    \item $\lceil \cdot \rceil$ l'operatore matematico di approssimazione all'intero superiore (\emph{ceiling}).
\end{itemize}

La legge di scala fondamentale è espressa da:
\begin{equation}
    \boxed{R_{\text{desired}} = \left\lceil R_{\text{curr}} \times \left( \frac{M_{\text{curr}}}{M_{\text{target}}} \right) \right\rceil}
    \label{eq:hpa_core}
\end{equation}

\subsubsection{Banda Di Tolleranza E Smorzamento Delle Oscillazioni}
Per evitare reazioni eccessive dovute a transitori microscopici o rumore nelle misurazioni, il control plane implementa una **banda di tolleranza globale** (configurabile, pari di default a $\varepsilon = 0.1$, ovvero il $10\%$):
\begin{equation}
    \text{Se } \left| \frac{M_{\text{curr}}}{M_{\text{target}}} - 1.0 \right| \le \varepsilon \implies R_{\text{desired}} = R_{\text{curr}}
\end{equation}
Se il rapporto tra valore corrente e valore desiderato ricade nell'intervallo compreso tra $0.9$ e $1.1$, l'HPA considera il sistema in equilibrio stazionario e scarta qualsiasi operazione di ridimensionamento.

\subsubsection{Esempi Di Calcolo Dettagliati}
\begin{itemize}
    \item \textbf{Esempio di Scale-Out (Picco di carico)}:
    Consideriamo un Deployment inizializzato con $R_{\text{curr}} = 2$ Pod. La politica impone che ciascun Pod consumi idealmente una quota target pari a $M_{\text{target}} = 200\,\text{m}$ (milli-core di CPU). A seguito di un afflusso massivo di richieste, il Metrics Server rileva un consumo medio pro-capite pari a $M_{\text{curr}} = 1000\,\text{m}$.
    Il rapporto delle metriche è:
    \[
        \frac{M_{\text{curr}}}{M_{\text{target}}} = \frac{1000\,\text{m}}{200\,\text{m}} = 5.0
    \]
    Poiché $|5.0 - 1.0| = 4.0 > 0.1$, il controller applica l'equazione~\eqref{eq:hpa_core}:
    \[
        R_{\text{desired}} = \left\lceil 2 \times 5.0 \right\rceil = \mathbf{10 \text{ repliche}}
    \]
    Il controller HPA aggiorna il Deployment, il quale invoca il ReplicaSet per schedulare istantaneamente $10 - 2 = 8$ nuovi Pod.

    \item \textbf{Esempio di Scale-In (Rilascio del carico)}:
    Ipotizziamo ora una flotta con $R_{\text{curr}} = 10$ repliche, in cui la metrica target desiderata sia $M_{\text{target}} = 100\,\text{m}$. Al calare della domanda, il consumo medio precipita a $M_{\text{curr}} = 50\,\text{m}$.
    Il rapporto risulta:
    \[
        \frac{M_{\text{curr}}}{M_{\text{target}}} = \frac{50\,\text{m}}{100\,\text{m}} = 0.5
    \]
    Essendo $|0.5 - 1.0| = 0.5 > 0.1$, il calcolo determina:
    \[
        R_{\text{desired}} = \left\lceil 10 \times 0.5 \right\rceil = \mathbf{5 \text{ repliche}}
    \]
    La flotta viene dimezzata terminando ordinatamente 5 Pod in eccesso.
\end{itemize}

\subsection{Specifica Dichiarativa HPA E Controllo Avanzato Del Comportamento}
Nel gruppo API \texttt{autoscaling/v2}, l'HPA consente di specificare metriche multiple basate su utilizzo percentuale o valori grezzi assoluti:

\begin{lstlisting}[language=yaml, caption={Configurazione metrica di risorsa nell'HPA (\texttt{autoscaling/v2})}]
type: Resource
resource:
  name: cpu
  target:
    type: Utilization
    averageUtilization: 60  # Mantiene l'utilizzo medio al 60% del request dichiarato
\end{lstlisting}

\subsubsection{Prevenzione Del Flapping Tramite Il Blocco \texttt{behavior}}
Nei sistemi distribuiti, una rapida oscillazione del carico può innescare il fenomeno del **flapping** (o \emph{thrashing}): la creazione immediata di Pod seguita, pochi secondi dopo, dalla loro distruzione, con conseguente spreco di cicli CPU per il bootstrap dei runtime e instabilità nel routing.
Per mitigare tale problematica, l'API \texttt{autoscaling/v2} espone il blocco \texttt{behavior}, che permette di definire politiche asimmetriche per le fasi di scale-up e scale-down.

\begin{lstlisting}[language=yaml, caption={Configurazione predefinita del blocco \texttt{behavior} nell'HPA}]
behavior:
  scaleDown:
    stabilizationWindowSeconds: 300  # Finestra di stabilizzazione (5 minuti)
    policies:
    - type: Percent
      value: 100                    # Consente di eliminare fino al 100% delle repliche correnti
      periodSeconds: 15
  scaleUp:
    stabilizationWindowSeconds: 0    # Reazione immediata ai picchi
    policies:
    - type: Percent
      value: 100                    # Consente di raddoppiare i Pod correnti
      periodSeconds: 15
    - type: Pods
      value: 4                      # Oppure consente di aggiungere un massimo di 4 Pod assoluti
      periodSeconds: 15
    selectPolicy: Max                # Sceglie la politica che produce la variazione maggiore
\end{lstlisting}

L'analisi dei parametri evidenzia la logica progettuale:
\begin{itemize}[noitemsep]
    \item \textbf{\texttt{scaleDown.stabilizationWindowSeconds: 300}}: Quando il carico decresce, il controller non termina immediatamente i Pod, ma analizza retroattivamente i dati degli ultimi 300 secondi (5 minuti). Il dimensionamento viene effettuato sul valore massimo raccomandato registrato all'interno di tale finestra, prevenendo riduzioni premature se il carico è soggetto a pause momentanee.
    \item \textbf{\texttt{scaleUp.stabilizationWindowSeconds: 0}}: In presenza di picchi di carico, l'infrastruttura non può attendere; pertanto, non vi è alcuna finestra di ritardo e l'espansione avviene istantaneamente.
    \item \textbf{\texttt{selectPolicy: Max}}: Posto a confronto l'incremento percentuale ($100\%$ dei Pod correnti) e l'incremento assoluto ($4$ Pod ogni 15 secondi), il sistema seleziona il vincolo che consente la crescita più rapida verso lo stato stazionario desiderato.
\end{itemize}

\subsection{Laboratorio Applicativo: Verifica Sperimentale Dell'HPA}
Per verificare sul campo il comportamento del controller, si predispone un ambiente di test minimale.

\subsubsection{Passo 1: Deployment Applicativo Con Risorse Vincolate}
Affinché l'HPA possa calcolare l'utilizzo percentuale, la specifica del Pod deve obbligatoriamente definire le risorse richieste (\texttt{resources.requests}):

\begin{lstlisting}[language=yaml, caption={Manifest di Deployment applicativo (\texttt{nginx-deploy.yaml})}]
apiVersion: apps/v1
kind: Deployment
metadata:
  name: nginx-deploy
  labels:
    type: dev
spec:
  replicas: 1
  selector:
    matchLabels:
      type: dev
  template:
    metadata:
      labels:
        type: dev
    spec:
      containers:
      - name: nginx
        image: nginx
        ports:
        - containerPort: 80
        resources:
          limits:
            cpu: 500m
          requests:
            cpu: 200m
\end{lstlisting}

Il carico viene distribuito tramite il comando:
\begin{lstlisting}[language=bash]
$ kubectl apply -f nginx-deploy.yaml
\end{lstlisting}

\subsubsection{Passo 2: Definizione Dell'Oggetto HPA}
È possibile definire la risorsa HPA tramite approccio dichiarativo o imperativo:
\begin{itemize}
    \item \textbf{Approccio Imperativo}:
\begin{lstlisting}[language=bash]
$ kubectl autoscale deployment nginx-deploy --cpu-percent=10 --min=1 --max=5
\end{lstlisting}
    \item \textbf{Approccio Dichiarativo (\texttt{hpa-test.yaml})}:
\begin{lstlisting}[language=yaml, caption={Manifest della risorsa HPA (\texttt{hpa-test.yaml})}]
apiVersion: autoscaling/v1
kind: HorizontalPodAutoscaler
metadata:
  name: nginx-deploy
spec:
  minReplicas: 1
  maxReplicas: 5
  scaleTargetRef:
    apiVersion: apps/v1
    kind: Deployment
    name: nginx-deploy
  targetCPUUtilizationPercentage: 10
\end{lstlisting}
\end{itemize}

Applicando il manifest e ispezionando lo stato iniziale:
\begin{lstlisting}[language=bash]
$ kubectl apply -f hpa-test.yaml
$ kubectl get hpa
NAME           REFERENCE                 TARGETS         MINPODS   MAXPODS   REPLICAS   AGE
nginx-deploy   Deployment/nginx-deploy   cpu: <unknown>/10%   1         5         1          27s
\end{lstlisting}
La dicitura \texttt{<unknown>} permane per alcuni secondi durante la fase di bootstrap iniziale, finché il Metrics Server non completa il primo ciclo di polling dai Kubelet.

\subsubsection{Passo 3: Esposizione E Generazione Del Carico Di Stress}
Per consentire l'invio di traffico, si espone il Deployment tramite un ClusterIP, quindi si instanzia un Pod client in loop continuo:
\begin{lstlisting}[language=bash]
# Esposizione del servizio
$ kubectl expose deployment nginx-deploy --type=ClusterIP --port=80 --target-port=80

# Avvio del generatore di carico sintetico (loop busybox)
$ kubectl run -i --tty load-generator --rm --image=busybox:1.28 --restart=Never -- /bin/sh -c "while sleep 0.01; do wget -q -O- http://nginx-deploy; done"
\end{lstlisting}

In una seconda finestra di terminale, il monitoraggio costante dei Pod evidenzierà la scalabilità automatica:
\begin{lstlisting}[language=bash]
$ kubectl get pods | grep nginx
nginx-deploy-5c4b8c6d7-abc12   1/1     Running   0          5m
nginx-deploy-5c4b8c6d7-def34   1/1     Running   0          45s
nginx-deploy-5c4b8c6d7-ghi56   1/1     Running   0          30s
nginx-deploy-5c4b8c6d7-jkl78   1/1     Running   0          15s
nginx-deploy-5c4b8c6d7-mno90   1/1     Running   0          15s
\end{lstlisting}
L'HPA incrementa la flotta fino al limite massimo configurato (\texttt{maxReplicas: 5}). All'interruzione del comando \texttt{wget}, decorso il periodo di stabilizzazione di 300 secondi, il controller provvederà al de-provisioning graduale dei Pod fino a raggiungere la cardinalità minima (\texttt{minReplicas: 1}).

\newpage

% =============================================================================
% SEZIONE: SICUREZZA E CONTROLLO ACCESSI (RBAC)
% =============================================================================

\section{Sicurezza E Controllo Degli Accessi: Il Modello RBAC Di Kubernetes}

In contesti aziendali multi-tenant, in cui team eterogenei (sviluppatori, operatori di sistema, auditor di sicurezza) condividono la medesima infrastruttura di calcolo, l'accesso indiscriminato alle primitive di cluster rappresenta una grave vulnerabilità.
Per regolare in modo deterministico le operazioni consentite, Kubernetes adotta un framework integrato denominato **Role-Based Access Control (RBAC)**.

\subsection{Principi Guida Del Controllo Degli Accessi}
Il framework RBAC risponde ai seguenti requisiti di governance:
\begin{itemize}[noitemsep]
    \item \textbf{Principio del Minimo Privilegio (\emph{Principle of Least Privilege})}: Ciascuna entità deve possedere esclusivamente i permessi strettamente indispensabili allo svolgimento delle proprie mansioni.
    \item \textbf{Isolamento Multi-Tenant}: Confinamento dei carichi di lavoro e delle credenziali all'interno di partizioni logiche distinte (\textbf{Namespace}).
    \item \textbf{Tracciabilità e Conformità (\emph{Audit \& Compliance})}: Registrazione immutabile delle azioni eseguite a fini di sicurezza forense.
\end{itemize}

\begin{theorem}
È fondamentale ribadire la netta separazione tra due concetti ortogonali:
\begin{enumerate}
    \item \textbf{Autenticazione (\emph{Authentication - Who you are})}: Identifica l'identità del chiamante che invia una richiesta all'API server.
    \item \textbf{Autorizzazione (\emph{Authorization - What you can do})}: Stabilisce se l'identità autenticata possiede i diritti necessari per invocare uno specifico verbo su una determinata risorsa. Il sottosistema RBAC governa \textbf{esclusivamente la fase di autorizzazione}.
\end{enumerate}
\end{theorem}

\subsection{Tipologie Di Utenti In Kubernetes: Normal Users vs. ServiceAccounts}
Nel database distribuito di Kubernetes (\textbf{etcd}) \textbf{non esiste un oggetto o una risorsa denominata ``User''}. Il sistema distingue formalmente due categorie di soggetti:

% [IMMAGINE DEL FILE DI SLIDE CON NOME "Kubernetes-HPA-RBAC-Deployment.pdf"] Diagramma architetturale che confronta Human Users (esterni) e ServiceAccounts (interni al cluster) verso l'API Server — Slide 163 / Pagina 17 del PDF

\begin{enumerate}
    \item \textbf{Utenti Normali (\emph{Human Users / External Accounts})}: Sono entità umane o software esterni (amministratori, tool di CI/CD) residenti al di fuori del cluster. La loro identità è validata mediante meccanismi crittografici o federati esterni. Non sono elencabili tramite comandi come \texttt{kubectl get users}.
    \item \textbf{Account di Servizio (\emph{ServiceAccounts})}: Sono entità di prima classe gestite nativamente dall'API di Kubernetes, concepite per assegnare un'identità formale ai processi e ai Pod che devono interagire con l'API server dall'interno del cluster. I ServiceAccount appartengono tassativamente a un determinato Namespace.
\end{enumerate}

\subsubsection{Meccanismi Di Autenticazione Degli Utenti Umani}
La validazione delle identità esterne può avvenire mediante quattro modalità:
\begin{itemize}
    \item \textbf{Certificati Client X.509}: È il metodo più comune. L'utente genera una coppia di chiavi asimmetriche e invia una richiesta di firma (\textbf{CSR}, \emph{Certificate Signing Request}) all'autorità di certificazione del cluster (\textbf{Cluster CA}, memorizzata in \texttt{/etc/kubernetes/pki/ca.crt} e \texttt{ca.key} su installazioni \emph{kubeadm}). L'API server interpreta:
    \begin{itemize}[noitemsep]
        \item Il campo **Common Name (CN)** come nome utente univoco (es. \texttt{CN=mario});
        \item Il campo **Organization (O)** come gruppo di appartenenza (es. \texttt{O=sviluppatori}).
    \end{itemize}
    \emph{Limitazione architetturale}: Kubernetes non supporta nativamente liste di revoca dei certificati (CRL) né il protocollo OCSP. Se una chiave privata viene compromessa, il certificato rimane valido fino alla sua naturale scadenza temporale.
    \item \textbf{OpenID Connect (OIDC)}: Lo standard de-facto per ambienti enterprise. Consente il Single Sign-On (SSO) federando l'autenticazione su \emph{Identity Provider} (IdP) quali Google Workspace, Microsoft Entra ID (Azure AD), Okta o Keycloak. L'utente esegue il login nel browser e riceve un \textbf{JSON Web Token (JWT)} firmato crittograficamente dall'IdP. Quando \texttt{kubectl} trasmette il token all'API server, quest'ultimo si limita a validare la firma digitale dell'IdP mediante la chiave pubblica esposta da quest'ultimo, ricavando i claim di identità ed i gruppi (\texttt{groups\_claim}) senza dover censire l'utente nel cluster.
    \item \textbf{Webhook Di Autenticazione}: L'API server delega la validazione a un servizio HTTP esterno inviando un payload di tipo \texttt{TokenReview}, a cui il server risponde confermando o rifiutando l'identità.
    \item \textbf{File Di Token Statici (\emph{Legacy - Non sicuro})}: Un file CSV locale sul master contenente la tupla \texttt{token,user,uid,"group1,group2"}. Richiede il riavvio manuale dell'API server a ogni variazione ed è severamente sconsigliato in produzione.
\end{itemize}

\subsubsection{Configurazione Del Contesto Locale In \texttt{kubectl}}
L'accesso tramite certificati X.509 si formalizza nel file di configurazione locale (\texttt{\textasciitilde/.kube/config}) definendo un **Context**, ovvero l'unione di tre elementi:
\begin{lstlisting}[language=bash]
# 1. Definizione delle credenziali utente (chiave e certificato firmato dalla CA)
$ kubectl config set-credentials mario --client-certificate=mario.crt --client-key=mario.key

# 2. Definizione del contesto operativo (Cluster + Utente + Namespace predefinito)
$ kubectl config set-context mario-context --cluster=mio-cluster-prod --user=mario --namespace=sviluppo

# 3. Attivazione del contesto
$ kubectl config use-context mario-context
\end{lstlisting}

\subsection{Gestione Avanzata Dell'Identità Dei ServiceAccount: Pod-Bound Tokens}
Nelle versioni storiche di Kubernetes (precedenti alla v1.21/v1.24), la creazione di un ServiceAccount generava automaticamente un token JWT statico a tempo indeterminato, memorizzato in una risorsa di tipo \texttt{Secret}. Se un attaccante esfiltrava tale Secret, otteneva un accesso permanente all'API server.

Nelle versioni moderne, il modello è stato integralmente riscritto adottando i **Short-Lived Pod-Bound Tokens**:
\begin{itemize}[noitemsep]
    \item I token non risiedono più come Secret persistenti nel database etcd.
    \item Il Kubelet richiede un token dinamico per ciascun Pod in fase di avvio mediante la **TokenRequest API**.
    \item La validità temporale del token è limitata (valore predefinito: \textbf{1 ora}).
    \item Il Kubelet monitora costantemente il ciclo di vita del token: al raggiungimento dell'$80\%$ della sua vita utile, provvede a rinnovarlo automaticamente senza imporre il riavvio del container.
\end{itemize}

\subsubsection{Meccanismo Dei Projected Volumes}
Il token generato non viene più inserito tramite un comune secret-volume, bensì iniettato nel filesystem del container tramite un **Projected Volume**:
\begin{lstlisting}[language=bash]
/var/run/secrets/kubernetes.io/serviceaccount/
|-- token        # Token JWT effimero e rinnovato automaticamente
|-- ca.crt       # Certificato della CA per validare l'identita dell'API Server
\-- namespace    # Stringa che definisce il namespace di appartenenza
\end{lstlisting}

% [IMMAGINE DEL FILE DI SLIDE CON NOME "Kubernetes-HPA-RBAC-Deployment.pdf"] Flusso completo di generazione, rilascio e iniezione del token Pod-bound tra Kubelet, TokenRequest API e API Server — Slide 168 / Pagina 24 del PDF

\subsubsection{Anatomia Del Payload Di Un Token Pod-Bound}
Se si estrae e si decodifica il token JWT iniettato nel container, è possibile riscontrare i vincoli di confinamento (\emph{cryptographic binding}) posti a livello di payload:

\begin{lstlisting}[caption={Payload decodificato del token JWT di un ServiceAccount}]
{
  "kubernetes.io": {
    "namespace": "my-namespace",
    "pod": {
      "name": "my-pod-5c4b8c6d7-abc12",
      "uid": "5e0bd49b-f040-43b0-99b7-22765a53f7f3"
    },
    "serviceaccount": {
      "name": "my-serviceaccount",
      "uid": "14ee3fa4-a7e2-420f-9f9a-dbc4507c3798"
    }
  },
  "exp": 1729605240,
  "sub": "system:serviceaccount:my-namespace:my-serviceaccount"
}
\end{lstlisting}

Questo vincolo implica che il token è valido **esclusivamente** se presentato dal Pod con quello specifico UID univoco: qualora il Pod venga distrutto, il token cessa istantaneamente di essere riconosciuto dall'API server, neutralizzando l'impatto di un furto di credenziali.

\subsection{I Quattro Pilastri Del Modello RBAC}
L'infrastruttura di autorizzazione poggia su quattro primitive complementari, differenziate in base al loro raggio d'azione (\emph{scope}):

\begin{table}[H]
\centering
\small
\begin{tabularx}{\textwidth}{l l X}
\toprule
\textbf{Risorsa RBAC} & \textbf{Ambito (\emph{Scope})} & \textbf{Finalità Operativa} \\
\midrule
\textbf{Role} & Singolo Namespace & Dichiara una collezione di regole di accesso (permessi) confinate a un namespace. \\
\addlinespace
\textbf{ClusterRole} & Intero Cluster & Dichiara regole di accesso non confinate, applicabili a risorse globali o cross-namespace. \\
\addlinespace
\textbf{RoleBinding} & Singolo Namespace & Assegna i permessi definiti in un \texttt{Role} a soggetti (utenti, gruppi, ServiceAccount). \\
\addlinespace
\textbf{ClusterRoleBinding} & Intero Cluster & Assegna i permessi di un \texttt{ClusterRole} a soggetti sull'intero perimetro del cluster. \\
\bottomrule
\end{tabularx}
\caption{I quattro pilastri architetturali del modello RBAC in Kubernetes.}
\label{tab:rbac_pillars}
\end{table}

% [IMMAGINE DEL FILE DI SLIDE CON NOME "Kubernetes-HPA-RBAC-Deployment.pdf"] Rappresentazione grafica dello scope di un Role confinato al singolo Namespace — Slide 169 / Pagina 25 del PDF

% [IMMAGINE DEL FILE DI SLIDE CON NOME "Kubernetes-HPA-RBAC-Deployment.pdf"] Rappresentazione grafica di un ClusterRole con visibilità globale cross-namespace sull'intero cluster — Slide 170 / Pagina 26 del PDF

\subsubsection{Composizione Delle Regole: API Groups, Resources E Verbs}
Una regola di accesso definita all'interno di un \texttt{Role} o \texttt{ClusterRole} è formalizzata da tre attributi fondamentali:
\begin{enumerate}
    \item \textbf{\texttt{apiGroups}}: Definisce la famiglia logica a cui appartiene la risorsa per evitare collisioni di denominazione:
    \begin{itemize}[noitemsep]
        \item \texttt{""} (stringa vuota): Identifica il **Core API Group**, che racchiude le primitive storiche di base: \texttt{pods}, \texttt{services}, \texttt{namespaces}, \texttt{nodes}, \texttt{configmaps}, \texttt{secrets}, \texttt{persistentvolumes}.
        \item \texttt{"apps"}: Racchiude i controller di orchestrazione: \texttt{deployments}, \texttt{statefulsets}, \texttt{daemonsets}, \texttt{replicasets}.
        \item \texttt{"batch"}: Racchiude i carichi a termine: \texttt{jobs}, \texttt{cronjobs}.
        \item \texttt{"networking.k8s.io"}: Racchiude primitive di rete avanzate quali \texttt{ingresses} e \texttt{networkpolicies}.
    \end{itemize}
    \item \textbf{\texttt{resources}}: L'elenco al plurale degli oggetti a cui si applica la regola (es. \texttt{["pods", "services"]}). Può indirizzare anche sotto-risorse specifiche (es. \texttt{pods/log}, \texttt{pods/exec}).
    \item \textbf{\texttt{verbs}}: L'insieme delle azioni consentite sul piano HTTP REST: \texttt{get}, \texttt{list}, \texttt{watch} (operazioni in sola lettura); \texttt{create}, \texttt{update}, \texttt{patch}, \texttt{delete} (operazioni di mutazione); oppure il carattere jolly \texttt{"*"} per concedere qualsiasi verbo.
\end{enumerate}

\subsubsection{ClusterRole Come Template Di Namespace}
Oltre a governare risorse prive di namespace (come \texttt{nodes}, \texttt{storageclasses} o \texttt{persistentvolumes}), un \texttt{ClusterRole} può essere referenziato all'interno di un comune **RoleBinding** di namespace. In tal caso, il ClusterRole agisce come un modello riutilizzabile: i permessi in esso dichiarati vengono concessi al soggetto \textbf{esclusivamente all'interno del namespace specifico} in cui risiede il RoleBinding. Questo pattern evita la duplicazione di decine di Role identici tra namespace distinti.

% [IMMAGINE DEL FILE DI SLIDE CON NOME "Kubernetes-HPA-RBAC-Deployment.pdf"] Diagramma del legame tra User/ServiceAccount, RoleBinding e Role all'interno di un Namespace — Slide 171 / Pagina 27 del PDF

% [IMMAGINE DEL FILE DI SLIDE CON NOME "Kubernetes-HPA-RBAC-Deployment.pdf"] Diagramma di associazione globale tramite ClusterRoleBinding tra gruppi e ClusterRole — Slide 172 / Pagina 28 del PDF

\subsection{Verifica Empirica E Strumenti Di Diagnostica RBAC}
Kubernetes offre primitive integrate per convalidare e simulare i privilegi associati a un'identità senza dover materialmente eseguire l'azione potenzialmente distruttiva.

\subsubsection{Il Comando \texttt{kubectl auth can-i}}
Il comando consente di interrogare il motore di autorizzazione dell'API server mediante il protocollo \texttt{SelfSubjectAccessReview}:
\begin{lstlisting}[language=bash]
# Verifica se l'utente corrente puo creare Pod nel namespace default
$ kubectl auth can-i create pods --namespace=default

# Verifica dell'identita attiva
$ kubectl auth whoami

# Simulazione di accesso mediante impersonificazione (richiede privilegi di impersonate)
$ kubectl auth can-i list nodes --as=gianluca@example.com

# Simulazione dei permessi attribuiti a uno specifico ServiceAccount
$ kubectl auth can-i get pods --as=system:serviceaccount:default:deploy-automation
\end{lstlisting}

\subsubsection{Tracciamento Degli Accessi Tramite Audit Logging}
Poiché gli utenti umani non risiedono in etcd, l'unico meccanismo forense per ricostruire la cronologia delle chiamate autenticate consiste nell'attivazione degli **Audit Logs** sul componente \textbf{kube-apiserver}.
Configurando una policy di audit, il server registra ciascuna transazione archiviando:
\begin{itemize}[noitemsep]
    \item \texttt{user}: Il Common Name o l'identificativo che ha inviato la richiesta;
    \item \texttt{groups}: L'insieme dei gruppi aziendali ereditati dal token OIDC o dal certificato;
    \item \texttt{sourceIPs}: L'indirizzo IP di provenienza della connessione TCP;
    \item \texttt{stage}: Lo stadio di elaborazione della richiesta (\emph{RequestReceived}, \emph{ResponseComplete}).
\end{itemize}

\subsection{Laboratorio Operativo: Isolamento Di Un Team Di Sviluppo}
Consideriamo il caso d'uso in cui si desideri autorizzare una sviluppatrice (Maria, identificata da \texttt{maria@example.com}) a gestire liberamente Pod, Service e Deployment all'interno del namespace \texttt{app-dev}, precludendole categoricamente l'accesso al resto del cluster.

\subsubsection{Definizione Del Ruolo Locale (\texttt{Role})}
\begin{lstlisting}[language=yaml, caption={Definizione del ruolo con privilegi completi sulle risorse applicative}]
apiVersion: rbac.authorization.k8s.io/v1
kind: Role
metadata:
  namespace: app-dev
  name: dev-full-access
rules:
- apiGroups: ["apps", ""]  # Copre sia i Deployment ("apps") sia Pod e Service ("")
  resources: ["deployments", "pods", "services"]
  verbs: ["*"]             # Concede get, list, watch, create, update, delete
\end{lstlisting}

\subsubsection{Associazione Del Ruolo All'Utente (\texttt{RoleBinding})}
\begin{lstlisting}[language=yaml, caption={Associazione del ruolo dev-full-access all'utente Maria}]
apiVersion: rbac.authorization.k8s.io/v1
kind: RoleBinding
metadata:
  namespace: app-dev
  name: maria-dev-binding
subjects:
- kind: User
  name: maria@example.com
  apiGroup: rbac.authorization.k8s.io
roleRef:
  kind: Role
  name: dev-full-access
  apiGroup: rbac.authorization.k8s.io
\end{lstlisting}

Se Maria tenta di eseguire operazioni al di fuori di \texttt{app-dev}:
\begin{lstlisting}[language=bash]
$ kubectl get pods --namespace=production --as=maria@example.com
Error from server (Forbidden): pods is forbidden: User "maria@example.com" cannot list resource "pods" in API group "" in the namespace "production"
\end{lstlisting}
L'operazione viene respinta immediatamente allo stadio di autorizzazione nel kernel dell'API server, garantendo l'isolamento multi-tenant.

\newpage

\section{Strumenti Di Bootstrap E Automazione Del Cluster: Kubeadm vs. Kubespray}

L'installazione e la gestione del ciclo di vita di un cluster Kubernetes rappresentano un compito ingegneristico complesso, implicante la generazione di catene di certificati crittografici, la configurazione sicura dei componenti di controllo (\emph{etcd}, \emph{api-server}, \emph{controller-manager}, \emph{scheduler}) e l'orchestrazione degli agenti di nodo. La comunità open-source ha formalizzato due strumenti d'elezione per rispondere a differenti requisiti operativi.

% [IMMAGINE DEL FILE DI SLIDE CON NOME "Kubernetes-HPA-RBAC-Deployment.pdf"] Panoramica degli strumenti di deployment del cluster con il logo ufficiale di Kubeadm — Slide 175 / Pagina 31 del PDF

\subsection{Kubeadm: Il Tool Ufficiale Di Bootstrap Minimale}
**Kubeadm** è lo strumento ufficiale manutenuto dal gruppo \emph{SIG Cluster Lifecycle} di Kubernetes, concepito per semplificare le operazioni di inizializzazione (\texttt{kubeadm init}) e annessione dei nodi worker (\texttt{kubeadm join}).

\subsubsection{Perimetro Operativo E Filosofia Di Progettazione}
La filosofia ingegneristica di kubeadm è rigorosamente orientata alla realizzazione di un cluster essenziale e funzionale (\emph{Minimum Viable Cluster}):
\begin{itemize}[noitemsep]
    \item \textbf{Cosa fa}: Genera in modo conforme le chiavi e i certificati TLS in \texttt{/etc/kubernetes/pki}; crea i manifest statici per avviare il control plane come container gestiti direttamente dal Kubelet locale; inizializza il database distribuito etcd; rilascia il file di configurazione amministrativo protetto \texttt{admin.conf}; genera il token crittografico con cui i nodi worker effettuano il bootstrap sicuro.
    \item \textbf{Cosa NON fa (Out of Scope)}: Kubeadm non si occupa del provisioning dell'infrastruttura sottostante (non alloca macchine virtuali su VMware o istanze EC2); non installa componenti opzionali o interfacce grafiche (come la Kubernetes Dashboard); non installa stack di monitoraggio o logging; soprattutto, \textbf{non include alcun plugin di rete CNI predefinito}. Al termine di \texttt{kubeadm init}, i nodi permangono nello stato \texttt{NotReady} fino a quando l'amministratore non provvede alla distribuzione manuale di una soluzione CNI (es. Calico, Flannel o Cilium).
\end{itemize}

\subsection{Kubespray: Automazione Multi-Nodo Enterprise Basata Su Ansible}
Per superare i limiti di kubeadm nell'automazione di infrastrutture complesse di livello enterprise (datacenter privati, scenari multi-cloud o ambienti bare-metal eterogenei), la comunità adotta **Kubespray**.

Kubespray è un framework di orchestrazione completo basato su playbook **Ansible**, che combina flessibilità dichiarativa e controllo granulare dell'infrastruttura:
\begin{itemize}[noitemsep]
    \item \textbf{Automazione End-to-End}: Si occupa della preparazione dell'intero sistema operativo host (configurazione dei parametri del kernel come \texttt{sysctl net.bridge.bridge-nf-call-iptables=1}, disattivazione permanente dello spazio di swap, installazione e configurazione del runtime dei container come \texttt{containerd}).
    \item \textbf{Integrazione CNI Avanzata}: Permette di scegliere e configurare dichiarativamente, tramite variabili YAML d'inventario, il plugin di rete desiderato (Calico con instradamento BGP o incapsulamento IP-in-IP/VXLAN, Cilium con dataplane eBPF, Flannel), gestendone i parametri topologici fin dal primo bootstrap.
    \item \textbf{Alta Affidabilità (\emph{High Availability - HA})}: Configura automaticamente architetture di control plane altamente disponibili multi-master, integrando bilanciatori di carico locali (come NGINX o HAProxy combinati con Keepalived) per distribuire le chiamate verso le istanze multiple di \texttt{kube-apiserver}.
    \item \textbf{Gestione Del Ciclo Di Vita}: Fornisce playbook dedicati non solo per l'installazione iniziale, ma anche per l'aggiornamento a caldo senza interruzione di servizio (\emph{zero-downtime rolling upgrades}), l'aggiunta o la rimozione controllata di nodi dal cluster, e l'integrazione di componenti accessori pre-validati (Metrics Server, Ingress NGINX, Cert-Manager).
\end{itemize}

In sintesi, mentre kubeadm costituisce il mattoncino elementare di bootstrap logico adottato anche all'interno di piattaforme gestite, Kubespray rappresenta la soluzione integrata a livello di produzione per automatizzare l'intero ciclo di vita del datacenter moderno.


\end{document}
